\documentclass[11pt]{amsart}
\usepackage{amsmath,amssymb,amsthm,mathtools,mathrsfs}
\usepackage[margin=1.08in]{geometry}
\usepackage{booktabs}
\usepackage{array}
\usepackage{longtable}
\usepackage{tabularx}
\usepackage[hypertexnames=false]{hyperref}
\hypersetup{hidelinks}
\usepackage[final,expansion=false]{microtype}
\usepackage{url}
\setlength{\parskip}{4pt plus 1pt}
\emergencystretch=2.5em

\newtheorem{theorem}{Theorem}[section]
\newtheorem{proposition}[theorem]{Proposition}
\newtheorem{lemma}[theorem]{Lemma}
\newtheorem{corollary}[theorem]{Corollary}
\theoremstyle{definition}
\newtheorem{definition}[theorem]{Definition}
\newtheorem{problem}[theorem]{Problem}
\theoremstyle{remark}
\newtheorem{remark}[theorem]{Remark}

\newcommand{\C}{\mathbb C}
\newcommand{\R}{\mathbb R}
\newcommand{\N}{\mathbb N}
\newcommand{\Z}{\mathbb Z}
\newcommand{\Q}{\mathbb Q}
\newcommand{\Rel}{\operatorname{Re}}
\newcommand{\Iml}{\operatorname{Im}}
\newcommand{\supp}{\operatorname{supp}}
\newcommand{\TV}{\operatorname{TV}}
\newcommand{\Kp}{\mathcal K^{\mathrm{prim}}_{2T}}
\newcommand{\Kw}{\mathcal K_{2T}}
\newcommand{\PT}{\mathcal P_T}
\newcommand{\IT}{I_T}
\newcommand{\lamP}{\lambda^{\mathrm{prim}}_T}
\newcommand{\lamW}{\lambda^{\mathrm W}_T}
\newcommand{\lamA}{\lambda_{\mathrm A}(T)}
\newcommand{\sT}{\sigma_T}
\newcommand{\gh}[1]{\widehat g_{#1}}
\newcommand{\nuh}{\widehat\nu}
\newcommand{\wh}{\widehat w}
\newcommand{\Fh}{\widehat F}
\newcommand{\cF}{c_F}
\newcommand{\Lev}{\mathfrak E}
\newcommand{\Wcl}{\mathcal W_T}
\newcommand{\NT}{\mathcal N_T}
\newcommand{\Aconst}{\mathsf A_T}
\newcommand{\Ainf}{\mathcal A_\infty}
\newcommand{\PiT}{\Pi_T}
\newcommand{\mT}{\mathfrak m_T}
\newcommand{\laminf}{\lambda^{\infty}_{\min}}
\newcommand{\laminfP}{\lambda^{\infty,\mathrm P}_{\min}}
\newcolumntype{Y}{>{\raggedright\arraybackslash}X}

\begin{document}

\title[Weil positivity on a finite window]
{Weil Positivity on a Finite Window:\\
Duality, Certificates and a Finite Reduction}

\author{David Alejandro Trejo Pizzo}
\thanks{Independent Researcher. \texttt{dtrejopizzo@gmail.com}. The certificate data and verification
scripts of the computer-assisted results accompany the paper.}

\date{\today}

\begin{abstract}
For $T>0$ let $\mathfrak a_T$ be the truncated Weil quadratic form on $L^2(-T,T)$, $\PT$ the
subspace on which the two Tate moments vanish, $\lamA$ the bottom of $\mathfrak a_T$ on $\PT$, and
$\lamW$ the bottom of the complete form $W_T=\mathfrak a_T+2\Rel(A\bar B)$ on $L^2(-T,T)$. We study
the window positivity problem $\lamA\ge0$ for every $T>0$. A certified computation in ball
arithmetic gives $\lamA\ge2.5\cdot10^{-12}$ and $\lamW\ge10^{-18}$ for all $T\le\frac12\log5$. For
general $T$ we develop the dual side of the problem. The archimedean multiplier is a L\'evy exponent
with an explicit completely monotone density; the bottom is bounded below by the value of a
semi-infinite linear program whose dual variables are measures supported outside $[-2T,2T]$; strong
duality holds with finitely atomic dual sequences, the dual value is never attained, and every
certificate of the complete form is, up to an arbitrarily small uniform correction, the explicit formula
evaluated on modulated flat-top test functions. An explicit certificate with $437$ atoms gives
$\lambda_{\mathrm A}(\frac12\log3)\ge8.8577\cdot10^{-5}$. The function $T\mapsto\lamA$ is continuous and
nonincreasing; at one window the complete and the primitive bottoms differ by a diagonal $2\times2$
polar capacity; for every $T$ the sign of $\lamA$ is the sign of an explicit symmetric matrix of size
$2k-2$ indexed by piecewise Tate exponentials; and a first failure of positivity, if any, occurs at a
window $T_*>\frac12\log5$ carrying a primitive null vector of definite parity. We also prove limits of
the method: ceilings and class impossibilities for lattice certificates, the exact decay
$\lamA\le4\cdot10^9Te^{-15.449T}$ of the margin, a doubly exponential envelope threshold, and the
failure at explicit windows of every argument that uses only properties of the primes stable under
finite modification of the weights. The remaining problem is stated precisely.
\end{abstract}

\maketitle
\tableofcontents

\section{Introduction}\label{sec:intro}

\subsection{The problem}
Throughout, $T>0$, $\IT=(-T,T)$, and for $F\in L^2(\IT)$ (extended by zero) we write
$\Fh(\tau)=\int F(t)e^{i\tau t}\,dt$ and $\cF(u)=\int F(t)\overline{F(t-u)}\,dt$ for the
autocorrelation, so that $\cF$ is continuous, supported in $[-2T,2T]$, and
$\widehat{\cF}=|\Fh|^2$ on $\R$. Put
\begin{equation}\label{eq:symbol}
m(\tau)=\Rel\psi\!\left(\tfrac14+\tfrac{i\tau}{2}\right)-\log\pi,\qquad
\sT(\tau)=m(\tau)-\sum_{n\in\NT}v_n\cos(\tau\log n),\qquad v_n=\frac{2\Lambda(n)}{\sqrt n},
\end{equation}
where $\psi$ is the digamma function, $\NT=\{n\ge2:\ \log n<2T\}$ is the set of \emph{contacts} and
$\Aconst:=\sum_{n\in\NT}v_n$ is the total mass of the prime comb. The truncated Weil form is
\begin{equation}\label{eq:form}
\mathfrak a_T(F)=\frac1{2\pi}\int_\R\sT(\tau)\,|\Fh(\tau)|^2\,d\tau
\end{equation}
on the form domain $\mathcal D_T=\{F\in L^2(\IT):\int\log(2+|\tau|)|\Fh|^2<\infty\}$. The two \emph{Tate
moments} are $A(F)=\int F(t)e^{t/2}dt$ and $B(F)=\int F(t)e^{-t/2}dt$, the \emph{primitive subspace} is
$\PT=\{F\in L^2(\IT):A(F)=B(F)=0\}$, and the \emph{complete Weil form} on the window is
\begin{equation}\label{eq:W}
W_T(F)=\mathfrak a_T(F)+2\,\Rel\big(A(F)\overline{B(F)}\big).
\end{equation}
By the explicit formula (\cite{Weil52}, \cite[Thm.~8.4]{Paper45}; Lemma~\ref{lem:roughEF} below for
tests that are not smooth), $W_T(F)=\sum_\rho\Phi_F(\rho)\Phi_F(1-\rho)$ with
$\Phi_F(s)=\int F(t)e^{(s-\frac12)t}dt$, the sum over the nontrivial zeros of $\zeta$. We write
\[
\lamA=\inf\{\mathfrak a_T(F):F\in\PT\cap\mathcal D_T,\ \|F\|=1\},\qquad
\lamW=\inf\{W_T(F):F\in\mathcal D_T,\ \|F\|=1\};
\]
$\lamW\le\lamA$ because the two forms agree on $\PT$. The object of this paper is the following
statement, which we call \emph{window positivity}:
\begin{equation}\label{eq:target}
\lamA\ \ge\ 0\qquad\text{for every }T>0 ,
\end{equation}
together with its complete version $\lamW\ge0$ for every $T>0$. Since $T\mapsto\lamA$ is nonincreasing
(Lemma~\ref{lem:monoT}), \eqref{eq:target} is equivalent to $\lim_{T\to\infty}\lamA\ge0$; it is a
statement about every window at once, and every window carries a finite amount of arithmetic data, the
vector $(\Lambda(n))_{n<e^{2T}}$ (Proposition~\ref{prop:finite}).

Positivity of the complete form is classical for $T\le\frac12\log2$, where the comb is empty
\cite{Yoshida1992,Bombieri2000}, and was certified by Zhu for $T\le0.8$ \cite[Thm.~1.2]{Zhu2026}. A
third bottom, $\lamP$, the value of a linear program over autocorrelations (Definition~\ref{def:PD}),
satisfies $\lamW\le\lamP\le\lamA$ (Theorem~\ref{thm:A}, Remark~\ref{rem:Kprim}); it is the quantity
that dual certificates bound from below.

\subsection{Results}
\emph{Positivity on an initial range.} A certified computation in ball arithmetic, built on a
frequency-truncated lower form, a Legendre section of $140$ modes per parity and a Schur and Haynsworth
inertia count, gives
\[
\lamA\ \ge\ 2.5\cdot10^{-12},\qquad \lamW\ \ge\ 10^{-18}\qquad\text{for every }0<T\le\tfrac12\log5
\]
(Theorems~\ref{thm:log5A} and \ref{thm:log5W}). The window $\frac12\log5$ is the last one with contact
set $\{2,3,4\}$. For comparison, converged Galerkin computations (numerical, not certified; Galerkin
values are upper bounds) give $\lambda_{\mathrm A}(\frac12\log3)\approx9.448\cdot10^{-4}$,
$\lambda_{\mathrm A}(\log2)\approx9.83\cdot10^{-8}$ and $\lambda_{\mathrm A}(\frac12\log5)\approx3.10\cdot10^{-12}$
(Appendix~\ref{app:num}). Independently, at $T=\frac12\log3$ an explicit dual feasible point with $437$
atoms proves $\lambda_{\mathrm A}(\frac12\log3)\ge\lambda^{\mathrm{prim}}_{\frac12\log3}\ge8.8577\cdot10^{-5}$
(Theorem~\ref{thm:C}).

\emph{Duality.} The archimedean multiplier is a L\'evy exponent:
$m(\tau)=C_0+\int_0^\infty(1-\cos\tau u)K(u)\,du$ with $C_0=\psi(\frac14)-\log\pi$ and the completely
monotone density $K(u)=\frac12/\sinh(u/2)+\frac12/\cosh(u/2)$; equivalently $m-C_0=\varphi(\tau^2)$ with
$\varphi$ a complete Bernstein function (Theorem~\ref{thm:S}). The bottom $\lamP$ is the value of a
semi-infinite linear program whose dual is: maximize $\inf_\tau[\sT+\nuh+\beta_0\gh0+\beta_1\gh1]$ over
finite even measures $\nu$ supported in $\{|u|\ge2T\}$ and $\beta\in\R^2$. Strong duality holds with
finitely atomic $\nu$ (Theorem~\ref{thm:D1}), so positivity with a margin is always provable by a finite
certificate; the dual is never attained (Theorem~\ref{thm:O}); exponential columns exhaust it
(Theorem~\ref{thm:E}); and $\beta=(2,0)$ is the unique multiplier for which a certificate bounds the
complete form (Proposition~\ref{prop:P}).

\emph{The dual variable is the explicit formula.} For an even weight $w$ with $w\equiv1$ on $[-2T,2T]$
and sufficient decay, an explicit measure $\nu_w$ outside the window is dual feasible and
$\sT+\nuh_w+2\gh0=f_w$ with $f_w(\tau)=\sum_\rho\wh(\tau-\gamma_\rho)$ (Theorem~\ref{thm:Z}); every
certificate of the complete form is of this kind up to an arbitrarily small uniform correction, and
$\sup_w\inf f_w=\lamW$ (Theorem~\ref{thm:Zexh}). A flat-top weight never has a nonnegative Fourier
transform (Lemma~\ref{lem:Rrig}), and no finite set of zeros carries a flat-top certificate
(Theorem~\ref{thm:finZ}). At each window the certificates are the absolutely continuous solutions of a
Kre\u{\i}n extension problem (Theorems~\ref{thm:KEP1}--\ref{thm:KEPdich}).

\emph{Structure at one window.} The function $T\mapsto\lamA$ is nonincreasing and continuous
(Lemmas~\ref{lem:monoT} and \ref{lem:contA}). The complete and the primitive form at one window are related
by a diagonal $2\times2$ matrix $\mathfrak G_T$, the polar capacity: $\lamW>0$ if and only if $\lamA>0$
and $\mathfrak G_T\succ0$ (Theorem~\ref{thm:parityschur}, Proposition~\ref{prop:polarcap}), and the
condition at one window controls all smaller ones (Proposition~\ref{prop:onewindow}). The archimedean form
alone is positive on $\PT$ up to a threshold $T^*_{\infty,\mathrm P}\in[0.58,0.60]$ and its negative index
is computed exactly for $T\le0.90$ outside two short intervals (Theorem~\ref{thm:index},
Corollary~\ref{cor:Tarch}); beyond the threshold the primes supply positivity.

\emph{A finite reduction.} Split $\IT$ into $k$ equal blocks and let $\mathcal V_k\subset\PT$ be the
functions satisfying both Tate conditions on every block. The archimedean form is coercive on
$\mathcal V_k$ with an explicit constant, $\mathfrak m_T\ge(\log\frac kT-C_B(T/k))\|\cdot\|^2$
(Theorem~\ref{thm:coer-B}), so for $k\ge k_0(T)$ the form $\mathfrak a_T$ is coercive there, and
\[
\lamA>0\iff S_{T,k}\succ0,\qquad \lamA\ge0\iff S_{T,k}\succeq0,
\]
where $S_{T,k}$ is the Schur complement of $\mathfrak a_T$ on the $(2k-2)$-dimensional space of piecewise
Tate exponentials (Theorem~\ref{thm:coer-schur}). Thus \eqref{eq:target} at each window is the positive
semidefiniteness of an explicit finite matrix; $k=25$ suffices at $T=\frac12\log5$.

\emph{First contact.} If \eqref{eq:target} fails, then by monotonicity, continuity and the certified range
there is a first window $T_*>\frac12\log5$ with $\lambda_{\mathrm A}(T_*)=0$, and the kernel of the
primitive operator there contains a nonzero vector of definite parity (Lemma~\ref{lem:firstcontact}). For
such a vector we prove a finite Poisson--Chebyshev operator identity combining Poisson summation with
$\Lambda*1=\log$ (Theorem~\ref{thm:PCfinite}), a faithful observation by the odd poles of the Gamma factor,
an exact balance identity, and a strict contraction of the Euler channel (Proposition~\ref{prop:GEgap}).

\subsection{Barriers}
Several results limit what any proof of \eqref{eq:target} can look like. The exchange between escaped
primal mass and value is exactly logarithmic (Theorem~\ref{thm:Vp}), and no pointwise envelope can lower
the verification threshold $2\pi e^{\Aconst}$, doubly exponential in $T$ (Proposition~\ref{prop:env}).
Lattice certificates have an explicit ceiling, one small Toeplitz eigenvalue (Proposition~\ref{prop:TK});
the ceilings converge to $\lamP$ (Theorem~\ref{thm:R}), but no lattice of order $M\le8$ certifies at
$\frac12\log3$ and none of order $M\le24$ at $T=0.68$ or $T=\log2$ (Theorems~\ref{thm:N}, \ref{thm:Np}),
and sampling a flat-top certificate on a lattice costs exactly $\Theta(1/M)$
(Theorem~\ref{thm:latedge}). No partial prime comb of mass below $-C_0$ is positive on a large window
(Theorem~\ref{thm:SC}). The margin decays, $\lamA\le4\cdot10^9Te^{-15.449T}$ for $T\ge1$
(Corollary~\ref{cor:margin1}). Modifications of $\Lambda$ that preserve its asymptotics, the zeros and the
pole of the associated Dirichlet series violate positivity at explicit windows (Theorem~\ref{thm:finmod}), so
no argument that uses only such properties of the primes decides any window $T>0.75$
(Corollary~\ref{cor:void}); the same modifications break the per-window implication from primitive to
complete positivity (Corollary~\ref{cor:convfail}). Every bounded remainder after subtracting a positive
archimedean part has infinite negative index once a contact is present (Theorem~\ref{thm:inertia-nogo}).

\subsection{What remains open}
Apart from the certified range, none of the results gives the sign of $\mathfrak a_T$ at a window
$T>\frac12\log5$. In the terms of this paper, the open problem is: \emph{for every $T>\frac12\log5$, and
some (equivalently every) $k\ge k_0(T)$, the matrix $S_{T,k}$ of Theorem~\ref{thm:coer-schur} is positive
semidefinite}; equivalently, no window carries a first-contact vector. The precise sub-problems
(Gamma-moment annihilation at first contact, class-dependent value bounds for lattice certificates, the
polar capacity beyond $\frac12\log5$) are collected in \S\ref{sec:open}.

\subsection{Related work}\label{subsec:related}
Positivity of the Weil functional on a fixed window goes back to Yoshida~\cite{Yoshida1992}, who proved
strict positivity of the Hermitian form at $a=\tfrac12\log 2$ by an approximate finite-matrix computation.
Bombieri~\cite{Bombieri2000} gave the variational theory: the infimum on the unit sphere of $L^2(E)$ is
attained \cite[Thm.~3]{Bombieri2000}, the window minima are continuous and decreasing
\cite[Thm.~5]{Bombieri2000}, and the number of negative eigenvalues equals the number of off-line conjugate
pairs \cite[Thm.~8]{Bombieri2000}. Suzuki's work on screw functions
\cite{Suzuki2023jlms,Suzuki2023screwline,Suzuki2025hilbert} recasts the form through Kre\u{\i}n's class
$\mathcal G_a$; Theorem~\ref{thm:finZ} is a finite-rank counterpart of that picture. Groskin~\cite{Groskin2026}
proves an exact finite Guinand--Weil dictionary, and Zhu~\cite{Zhu2026} reduces the window problem to one
finite matrix by a pointwise envelope, certifies $Q(f)\ge 8.9\cdot 10^{-18}\|f\|_2^2$ for
$\supp f\subseteq[-0.8,0.8]$, and proves a doubly exponential barrier for pointwise-envelope certificates
\cite[Thm.~1.4]{Zhu2026}. The archimedean form alone was studied by Connes and Consani
\cite{ConnesConsani2021}. The Delsarte and Cohn--Elkies linear programming method \cite{CohnElkies2003} is
the structural analogue of Theorem~\ref{thm:D1}. What is new here is the dual formulation and its
identification with the explicit formula, the structural results at one window, the finite reduction, the
first-contact analysis, and the extension of the certified range to $\frac12\log5$.

\subsection{Conventions}\label{subsec:conv}
$\widehat h(\tau)=\int_\R h(u)e^{i\tau u}du$; for even $h$ this is $\int h\cos(\tau u)$. For a finite even
measure $\nu$, $\nuh(\tau)=\int\cos(\tau u)\,d\nu(u)$, and for a probability measure $\mu$ on $\R$,
$\check\mu(u)=\int\cos(\tau u)\,d\mu(\tau)$. The letter $\tau$ always denotes a frequency, and $u$ a
variable dual to it. The right shift is $(V_aF)(t)=F(t-a)$, the reflection is $(JF)(t)=F(-t)$, and $Q_T$
denotes restriction to $\IT$ followed by extension by zero. We write $\mT(F)=\frac1{2\pi}\int m|\Fh|^2$
for the archimedean form, $\PiT(F)=-\sum_{n\in\NT}v_n\Rel\cF(\log n)$ for the prime form, so that
$\mathfrak a_T=\mT+\PiT$, and $\Ainf=\mT+2\Rel(A\bar B)$ for the complete archimedean form. The functions
\begin{equation}\label{eq:g}
g_0(u)=\cosh\tfrac u2\,\mathbf 1_{|u|<2T},\qquad g_1(u)=u\sinh\tfrac u2\,\mathbf 1_{|u|<2T},
\end{equation}
have transforms $\gh0,\gh1$ (Lemma~\ref{lem:gbounds}). For a parity label
$\mathrm{par}\in\{\mathrm{ev},\mathrm{odd}\}$, $\chi_{\mathrm{ev}}=\cosh\frac t2\,\mathbf 1_{\IT}$ and
$\chi_{\mathrm{odd}}=\sinh\frac t2\,\mathbf 1_{\IT}$. Letters not listed here are local to the proof or
the subsection in which they occur.

Statements whose proof includes a computer verification in interval or ball arithmetic say so; the
methods are in Appendix~\ref{app:encl}. Galerkin eigenvalues are upper bounds for bottoms and are quoted
only as numerical values, never as inputs to a proof (Appendix~\ref{app:num}). The explicit formula is used
in the form \cite[Thm.~8.4]{Paper45}, i.e.\ Weil's formula \cite{Weil52}, \cite[Thm.~5.12]{IK}, in the
normalization above.


\section{Preliminaries: the archimedean multiplier}\label{sec:levy}

\begin{lemma}[the L\'evy density]\label{lem:K}
For $u>0$ put $K(u)=\dfrac{2e^{3u/2}}{e^{2u}-1}=\dfrac{1}{2\sinh(u/2)}+\dfrac{1}{2\cosh(u/2)}
=2\sum_{k\ge0}e^{-(2k+\frac12)u}$. Then $K$ is positive, completely monotone,
$K(u)=u^{-1}+O(1)$ as $u\downarrow0$, $K(u)=2e^{-u/2}(1+O(e^{-u}))$ as $u\to\infty$, and
\begin{equation}\label{eq:m}
m(\tau)=\Rel\psi\!\left(\tfrac14+\tfrac{i\tau}2\right)-\log\pi
= C_0+\int_0^\infty(1-\cos\tau u)\,K(u)\,du,\quad C_0=\psi(\tfrac14)-\log\pi=-5.3721834\ldots
\end{equation}
\end{lemma}
\begin{proof}
The three expressions for $K$ agree: $\frac{1}{2\sinh(u/2)}+\frac1{2\cosh(u/2)}
=\frac{\cosh(u/2)+\sinh(u/2)}{\sinh u}=\frac{2e^{u/2}}{e^u-e^{-u}}=\frac{2e^{3u/2}}{e^{2u}-1}$, and
$\frac{2e^{3u/2}}{e^{2u}-1}=2e^{-u/2}\sum_{k\ge0}e^{-2ku}$. Complete monotonicity is that of a
convergent series of decaying exponentials with positive coefficients. For \eqref{eq:m} use
$\psi(z)=-\gamma+\int_0^\infty\frac{e^{-t}-e^{-zt}}{1-e^{-t}}dt$ ($\Rel z>0$): with $z=\frac14+\frac{i\tau}2$,
\[
\Rel\psi(z)-\psi(\tfrac14)=\int_0^\infty\frac{e^{-t/4}\big(1-\cos(\tau t/2)\big)}{1-e^{-t}}\,dt
=\int_0^\infty\frac{2e^{-u/2}(1-\cos\tau u)}{1-e^{-2u}}\,du\qquad(t=2u),
\]
and $\frac{2e^{-u/2}}{1-e^{-2u}}=\frac{2e^{3u/2}}{e^{2u}-1}=K(u)$. The integral converges since
$1-\cos\tau u=O(u^2)$ at $0$ and $K=O(e^{-u/2})$ at $\infty$.
\end{proof}

\begin{theorem}[complete Bernstein structure]\label{thm:S}
Put $a_k=k+\frac14$ and $\varphi(x)=\sum_{k\ge0}\dfrac{x}{a_k(x+4a_k^2)}$ for $x\ge0$. Then
\emph{(i)} $\Rel\psi(\frac14+\frac{i\tau}2)-\psi(\frac14)=\varphi(\tau^2)$ and $\varphi$ is a
complete Bernstein function; \emph{(ii)} the L\'evy measure of the negative-definite function
$\tau\mapsto\varphi(\tau^2)$ on $\R$ is $\frac12K(|u|)du$, which is completely monotone on each
half-line, so the associated L\'evy process is a pure-jump subordinate Brownian motion
$B_{S_t}$ with a special subordinator $S$ of Laplace exponent $\varphi$; in particular it has no
Gaussian component; \emph{(iii)} $\int_{2T}^\infty K(u)\,du=\log\coth\frac T2+\pi-2\arctan e^{T}$.
\end{theorem}
\begin{proof}
(i) From $\psi(z)=-\gamma+\sum_{k\ge0}\big[\frac1{k+1}-\frac1{k+z}\big]$,
$\Rel\big[\frac1{a_k}-\frac1{a_k+i\tau/2}\big]=\frac1{a_k}-\frac{a_k}{a_k^2+\tau^2/4}
=\frac{\tau^2}{a_k(4a_k^2+\tau^2)}$, which is the stated series at $x=\tau^2$. Each summand
$x\mapsto a_k^{-1}\,x/(x+4a_k^2)$ is a complete Bernstein function (Stieltjes representation with
a point mass at $4a_k^2$), the summands are $O(k^{-3})$ locally uniformly, and the complete
Bernstein functions form a convex cone closed under pointwise limits
\cite[Ch.~6--7]{SSV2012}. (ii) By Lemma~\ref{lem:K}, $\varphi(\tau^2)=\int_\R(1-\cos\tau u)\,
\frac12K(|u|)du$, so the L\'evy measure is $\frac12K(|u|)du$; it is completely monotone on
$(0,\infty)$ and $\varphi(x)/x\to0$, so there is no Gaussian part; the representation as a
subordinate Brownian motion with a special subordinator is \cite[Thm.~13.10, Cor.~13.12]{SSV2012}.
(iii) $\int\frac{du}{2\sinh(u/2)}=\log\tanh\frac u4$ and $\int\frac{du}{2\cosh(u/2)}=2\arctan e^{u/2}$.
\end{proof}

\begin{theorem}[the form as a L\'evy energy]\label{thm:A}
For every $F\in L^2(\IT)$, with $\Lev(\Phi):=\int_0^\infty(\Phi(0)-\Rel\Phi(u))K(u)\,du\in[0,\infty]$
(for complex $F$ the autocorrelation is Hermitian and only its even part $\Rel\cF$ enters; the
prime term below is likewise $\Rel\cF(\log n)$, and from \S\ref{sec:duality} on all measures are even),
\begin{equation}\label{eq:Aform}
\mathfrak a_T(F)=C_0\|F\|^2+\Lev(\cF)-\sum_{n\in\NT}v_n\cF(\log n),
\end{equation}
both sides being finite exactly when $\int\log(2+|\tau|)|\Fh|^2<\infty$. Consequently, with
$\Kw=\{\cF:F\in L^2(\IT),\|F\|=1\}$ and the primitive class
\begin{equation}\label{eq:Kprim}
\Kp=\Big\{\Phi\in\Kw:\ \int\Phi(u)\cosh\tfrac u2\,du=0,\ \int\Phi(u)\,u\sinh\tfrac u2\,du=0\Big\},
\end{equation}
one has
\begin{align*}
\lamP&:=\inf_{\Phi\in\Kp}\Big[C_0+\Lev(\Phi)-\sum_{n\in\NT}v_n\Phi(\log n)\Big]=C_0-\Lambda_T,\\
\Lambda_T&:=\sup_{\Phi\in\Kp}\Big[\sum_{n\in\NT}v_n\Phi(\log n)-\Lev(\Phi)\Big],
\end{align*}
and $\lamP\le\lamA$.
\end{theorem}
\begin{proof}
By \eqref{eq:m} and Tonelli (both $1-\cos\tau u$ and $K$ are nonnegative),
$\frac1{2\pi}\int m|\Fh|^2=C_0\|F\|^2+\int_0^\infty K(u)\Big[\frac1{2\pi}\int(1-\cos\tau u)|\Fh|^2d\tau\Big]du$,
and $\frac1{2\pi}\int\cos(\tau u)|\Fh(\tau)|^2d\tau=\cF(u)$ by Plancherel; likewise for the
prime terms. Finiteness: $\int_0^\infty(1-\cos\tau u)K(u)du=\log|\tau|+O(1)$ as $|\tau|\to\infty$
by Lemma~\ref{lem:K}, since $\int_0^1(1-\cos\tau u)\frac{du}u=\log|\tau|+O(1)$. For the last
statements: the two integrals in \eqref{eq:Kprim} are $\Rel(A\bar B)$ and $\Rel(A'\bar B-A\bar B')$
with $A'=\int tF(t)e^{t/2}dt$, $B'=\int tF(t)e^{-t/2}dt$ (expand $\cosh\frac{t-s}2$ and
$(t-s)\sinh\frac{t-s}2$ in exponentials), so $\{\cF:F\in\PT,\|F\|=1\}\subseteq\Kp$, whence
$\lamP\le\lamA$. In terms of $\widehat\Phi$, an entire function of exponential type $2T$ with
$\widehat\Phi(i/2)=\overline{\widehat\Phi(-i/2)}$, the two conditions say that
$\Rel\widehat\Phi(i/2)=0$ and $\Rel\big[i(\widehat\Phi'(i/2)-\widehat\Phi'(-i/2))\big]=0$; when
$F\in\PT$ one has the stronger $\widehat\Phi(\pm i/2)=\widehat\Phi'(\pm i/2)=0$, i.e.\ $\widehat\Phi$
vanishes to order $\ge2$ at $\pm i/2$ (Remark~\ref{rem:Kprim}).
\end{proof}

\begin{remark}\label{rem:Kprim}
$\Kp$ is slightly larger than the set of primitive autocorrelations: for real $F$ its two
conditions read $AB=0$ and $A'B-AB'=0$, i.e.\ $F\in\PT$, or $A=A'=0$, or $B=B'=0$. Every
certificate below bounds $\lamP$, hence $\lamA$; the equality $\lamP=\lamA$ is not needed
anywhere and is not asserted. Since $W_T(F)=\mathfrak a_T(F)$ whenever $A\bar B$ is purely
imaginary, also $\lamW\le\lamP$.
\end{remark}

\begin{corollary}[monotonicity; upper bound]\label{cor:mono}
With the contact set held fixed, $T\mapsto\Lambda_T$ is nondecreasing and $T\mapsto\lamP$
nonincreasing. Moreover
$\lamA\le C_0+\Aconst+\varphi\big((3\pi/2T)^2\big)$, which equals $1.291$ at $T=\frac12\log3$ and
$2.327$ at $T=\log 2$.
\end{corollary}
\begin{proof}
Monotonicity: $\Kp$ grows with $T$ and $\Lev$ does not depend on $T$. Upper bound: the form
$\Lev(\cF)$ is the Dirichlet form of the L\'evy process of Theorem~\ref{thm:S} killed on
leaving $\IT$; its eigenvalues are bounded by those of the subordinate killed Brownian motion,
$\lambda_n\le\varphi((n\pi/2T)^2)$ \cite[Prop.~2.3]{SongVondracek2003}, the Brownian motion being
normalized with characteristic exponent $\xi^2$; $\PT$ has codimension
$2$, so the bottom of $\Lev$ on $\PT$ is at most $\lambda_3$; and $-\sum v_n\cF(\log n)\le\Aconst\|F\|^2$.
\end{proof}

The following elementary bounds on $m$ and on the prime comb are used throughout.

\begin{lemma}[elementary bounds]\label{lem:mbounds}
\emph{(i)} $m$ is increasing on $[0,\infty)$. \emph{(ii)} For every $\tau>0$,
\[
\log\frac{\tau}{2\pi}-\frac1\tau-\frac1{2\tau^2}\ \le\ m(\tau)\ \le\ \log\frac{\tau}{2\pi}+\frac1\tau+\frac12\log\Big(1+\frac1{4\tau^2}\Big),
\]
and the lower bound, which we denote by $m_{\mathrm{lo}}(\tau)$, is increasing in $\tau$. \emph{(iii)} $-\sum_{n\in\NT}v_n\cos(\tau\log n)\ge-\Aconst$, with
$\sup_\tau\sum_{n\in\NT}v_n\cos(\tau\log n)=\Aconst$.
\end{lemma}
\begin{proof}
(i) In $\Rel\psi(\frac14+iy)=-\gamma+\sum_{n\ge0}\big[\frac1{n+1}-\frac{n+\frac14}{(n+\frac14)^2+y^2}\big]$
every summand increases with $|y|$. (ii) Use $\psi(z)=\log z-\frac1{2z}-\int_0^\infty g(t)e^{-zt}dt$
for $\Rel z>0$, $g(t)=\frac1{e^t-1}-\frac1t+\frac12$ \cite[5.9.13]{DLMF}, at $z=\frac14+iy$, $y=\tau/2$.
Here $g(0^+)=0$, $g'(t)=t^{-2}-\frac14\sinh^{-2}(t/2)\in(0,\infty)$ and $g(\infty)=\frac12$, so one
integration by parts gives $|\int_0^\infty g e^{-zt}dt|=|z|^{-1}|\int_0^\infty g'e^{-zt}dt|\le\frac1{2|z|}\le\frac1{2y}$.
Also $\Rel\log z=\frac12\log(\frac1{16}+y^2)\in[\log y,\log y+\frac12\log(1+\frac1{16y^2})]$ and
$\Rel\frac1{2z}=\frac{1/8}{1/16+y^2}\in[0,\frac1{8y^2}]$. Substituting $y=\tau/2$ and $-\log\pi$
gives (ii); the derivative of the lower bound is $\tau^{-1}+\tau^{-2}+\tau^{-3}>0$. (iii) The
first is trivial. For the second, the numbers $\log p$, $p$ prime, are linearly independent over
$\Q$, so by Kronecker's theorem the vector $(\tau\log p)_{p<e^{2T}}$ is dense modulo $2\pi$; at
$\tau$ where all $\cos(\tau\log p)>1-\varepsilon$ also $\cos(\tau k\log p)>1-k^2\varepsilon$, so the
sum over all $n\in\NT$ exceeds $\Aconst-O(\varepsilon)$.
\end{proof}

The hypothesis that the contact set is fixed can be removed from Corollary~\ref{cor:mono}, for
all three bottoms at once.

\begin{lemma}[monotonicity across contact thresholds]\label{lem:monoT}
The functions $T\mapsto\lamA$, $T\mapsto\lamP$ and $T\mapsto\lamW$ are nonincreasing on $(0,\infty)$,
including at the points $T=\frac12\log n$ where a contact enters.
\end{lemma}
\begin{proof}
Let $T<T'$ and extend $F\in L^2(\IT)$ by zero to $I_{T'}$. The norm, the Tate moments $A(F),B(F)$ and
the autocorrelation $\cF$ are unchanged, and $\cF$ vanishes on $|u|\ge2T$. A contact
$n\in\mathcal N_{T'}\setminus\NT$ has $\log n\ge2T$, so $\cF(\log n)=0$, and by \eqref{eq:Aform}
$\mathfrak a_{T'}(F)=\mathfrak a_T(F)$. Hence $\PT\subset\mathcal P_{T'}$ isometrically with the same
form values, which gives the statement for $\lamA$; the same computation with $W_T=\mathfrak a_T+2\Rel(A\bar B)$
gives it for $\lamW$, and for $\lamP$ the class $\Kp$ grows with $T$, $\Lev$ is $T$-free and the new
contacts contribute $\Phi(\log n)=0$.
\end{proof}

For the complete form, continuity is \cite[Thm.~5]{Bombieri2000}:
Bombieri proves that the even and odd minima $\mu_\pm$ of the complete functional
on $L^2([M^{-1},M])$ are continuous and decreasing. Here $W_T$ splits into its values on the even and
odd parts of $F$: the symbol is even, and $\Rel(A\bar B)=|C|^2-|S|^2$ with $C=\int F\cosh\frac t2$, which
sees only the even part, and $S=\int F\sinh\frac t2$, which sees only the odd part. So, after the change of
variables $x=e^t$, $\lamW=\min(\mu_+,\mu_-)$ at $M=e^T$. A direct proof in the present normalization follows.

\begin{lemma}[continuity of the complete bottom]\label{lem:contT}
The function $T\mapsto\lamW$ is continuous on $(0,\infty)$.
\end{lemma}
\begin{proof}
Fix $T$ and let $r\in[\frac12,2]$. The dilation $U_rF(t)=r^{-1/2}F(t/r)$ is unitary from $L^2(\IT)$ onto
$L^2(I_{rT})$, with $\widehat{U_rF}(\tau)=r^{1/2}\Fh(r\tau)$ and $c_{U_rF}(u)=\cF(u/r)$. It maps the form
domain onto the form domain, since $\log(2+|s|/r)\le\log(2+|s|)+\log2$. Hence
$\lambda^{\mathrm W}_{rT}=\inf\{W_{rT}(U_rF):F\text{ in the form domain},\ \|F\|=1\}$. By \eqref{eq:form}
and \eqref{eq:W},
\[
W_{rT}(U_rF)=\frac1{2\pi}\int_\R m(s/r)|\Fh(s)|^2ds-\sum_{2\le n<e^{4T}}v_n\Rel\cF(\tfrac{\log n}r)
+2\Rel\big(A(U_rF)\overline{B(U_rF)}\big),
\]
where the comb may be summed over $n<e^{4T}$ for both windows, because $\cF$ vanishes on $|u|\ge2T$. This
also takes care of contacts that enter or leave. We compare the three terms with those of $W_T(F)$.
\emph{(i) Archimedean.} $m$ is even and $m'(x)=-\frac12\Iml\psi'(\frac14+\frac{ix}2)$ with
$\psi'(z)=O(1/|z|)$ on $\Rel z=\frac14$. So $M_1:=\sup_x|xm'(x)|<\infty$ and
$|m(s/r)-m(s)|\le M_1|\log r|$, and the difference of the archimedean terms is at most
$M_1|\log r|\,\|F\|^2$.
\emph{(ii) Comb.} $|\Rel\cF(u)-\Rel\cF(u')|\le\frac1{2\pi}\int\min(2,|\tau||u-u'|)|\Fh|^2d\tau$. Put
$\delta=4T|1/r-1|$ and $\|F\|_Q^2:=\frac1{2\pi}\int\log(2+|\tau|)|\Fh|^2d\tau$. Splitting at
$|\tau|=\delta^{-1/2}$, the difference of the comb terms is at most
$S_T\big(\delta^{1/2}\|F\|^2+2\|F\|_Q^2/\log(2+\delta^{-1/2})\big)$, where $S_T=\sum_{n<e^{4T}}v_n$.
\emph{(iii) Polar.} $A(U_rF)=r^{1/2}\int F(s)e^{rs/2}ds$, so
$|A(U_rF)-A(F)|\le\|F\|\,\|r^{1/2}e^{rs/2}-e^{s/2}\|_{L^2(\IT)}$, and likewise for $B$. Both moments are
bounded by $C_T\|F\|$ uniformly in $r\in[\frac12,2]$. So the difference of the polar terms is at most
$\delta_P(r)\|F\|^2$ with $\delta_P(r)\to0$ as $r\to1$.
By Lemma~\ref{lem:mbounds}(ii) there is $C_1$ with $\log(2+|\tau|)\le m(\tau)+C_1$. By
Lemma~\ref{lem:mbounds}(iii) and $|2A\bar B|\le4\sinh T\,\|F\|^2$ (Cauchy--Schwarz), this gives
$\|F\|_Q^2\le W_T(F)+C'\|F\|^2$ with $C'=C_1+\Aconst+4\sinh T$. Collecting (i)--(iii), there is
$\eta(r)\to0$ as $r\to1$ with
\[
|W_{rT}(U_rF)-W_T(F)|\le\eta(r)\big(W_T(F)+C'\|F\|^2\big)
\]
for all $F$ in the form domain, and the right side is nonnegative. Take infima over unit $F$. This gives
$(1-\eta)\lamW-\eta C'\le\lambda^{\mathrm W}_{rT}\le(1+\eta)\lamW+\eta C'$, so
$|\lambda^{\mathrm W}_{rT}-\lamW|\le\eta(r)(|\lamW|+C')\to0$.
\end{proof}

The same proof with the comb deleted gives continuity of the bottom $\laminf$ of $\Ainf$
(\S\ref{sec:archstr}). It does not apply directly to $\lamA$, because $U_r$ does not preserve $\PT$;
the correction of the two Tate moments is the content of Lemma~\ref{lem:contA} below.

\begin{lemma}[compactness]\label{lem:compact}
For every $V\in\R$ the set $\{F\in\mathcal D_T:\|F\|=1,\ \mathfrak a_T(F)\le V\}$ is relatively compact in
$L^2(\IT)$. Consequently the form $\mathfrak a_T$ restricted to $\PT\cap\mathcal D_T$ is the form of a
self-adjoint operator $H_T$ on $\PT$ with compact resolvent, $\lamA=\min\operatorname{spec}H_T$ is attained,
and the bottoms of $\mathfrak a_T$ and $W_T$ on the unit sphere of $L^2(\IT)$, and on any closed subspace, are
attained.
\end{lemma}
\begin{proof}
By Lemma~\ref{lem:mbounds}, $\sT\ge m-\Aconst$ and $m-C_0\ge0$, so for $\|F\|=1$ and
$\mathfrak a_T(F)\le V$ one has $\frac1{2\pi}\int(m-C_0)|\Fh|^2\le V+\Aconst-C_0$, hence
$\frac1{2\pi}\int_{|\tau|>R}|\Fh|^2\le(V+\Aconst-C_0)/(m(R)-C_0)\to0$ as $R\to\infty$, uniformly. For each
$R$ the map $F\mapsto\Fh\mathbf 1_{[-R,R]}$ is Hilbert--Schmidt from $L^2(\IT)$ to $L^2(\R)$ (its kernel
$e^{i\tau t}$ is square integrable on $[-R,R]\times\IT$), hence compact. A diagonal argument gives a
subsequence along which $\Fh_n$ converges in $L^2([-R,R])$ for every $R$, and with the uniform tail bound
$\Fh_n$ is Cauchy in $L^2(\R)$. The form $\mathfrak a_T$ is closed and bounded below by $C_0-\Aconst$, since its
form norm is equivalent to $\int\log(2+|\tau|)|\Fh|^2$ (Lemma~\ref{lem:mbounds}(ii)); $\PT$ is closed, and
the polar term $2\Rel(A\bar B)$ is continuous on $L^2(\IT)$. Lower semicontinuity of closed forms gives
attainment, and relative compactness of the sublevel sets is compactness of the resolvent.
\end{proof}

On $\PT$ the operator $H_T$ commutes with $J$, since $\sT$ is even and $A(JF)=B(F)$; so its eigenspaces
split into even and odd parts.

\begin{lemma}[continuity of the primitive bottom]\label{lem:contA}
The function $T\mapsto\lamA$ is continuous on $(0,\infty)$.
\end{lemma}

\begin{proof}
Fix $T$ and consider windows $T'\in[T/2,2T]$ and factors $r\in[\frac12,2]$ with $rT'\in[T/2,2T]$; all
constants depend only on $T$. Put $M_0=\mathsf A_{2T}-C_0$, so that $\mathfrak a_S+M_0\|\cdot\|^2\ge0$
for every $S\le2T$ (Lemma~\ref{lem:mbounds}).

\emph{Step 1.} Steps (i) and (ii) of the proof of Lemma~\ref{lem:contT} do not involve the polar term.
Taking $\delta=8T|1/r-1|$ and summing the comb over $n<e^{8T}$, which covers every window in
$[T/2,2T]$ and every contact that enters or leaves, they give $\eta(r)\to0$ as $r\to1$ with
$|\mathfrak a_{rT'}(U_rF)-\mathfrak a_{T'}(F)|\le\eta(r)\big(\mathfrak a_{T'}(F)+C\|F\|^2\big)$ for all
$F$ in the form domain of $\mathfrak a_{T'}$, uniformly in $T'$.

\emph{Step 2.} If $F\in\mathcal P_{T'}$, then $A(U_rF)=r^{1/2}\int F(s)\big(e^{rs/2}-e^{s/2}\big)ds$,
so $|A(U_rF)|\le\varepsilon(r)\|F\|$ with $\varepsilon(r)=r^{1/2}\|e^{rs/2}-e^{s/2}\|_{L^2(I_{2T})}\to0$,
and likewise for $B$.

\emph{Step 3.} Fix real $\phi_1,\phi_2\in C_c^\infty(I_{T/4})$ for which the matrix
$\big(A(\phi_j),B(\phi_j)\big)_{j}$ is invertible. By zero extension (proof of Lemma~\ref{lem:monoT}),
$\mathfrak a_S(\phi_i,\phi_j)$ does not depend on $S\ge T/4$. For $u=U_rF$ with $F\in\mathcal P_{T'}$,
$\|F\|=1$, let $G=c_1\phi_1+c_2\phi_2$ have the same Tate moments as $u$. Then
$v:=u-G\in\mathcal P_{rT'}$, $\|G\|\le C_2\varepsilon(r)$ and $\mathfrak a_{rT'}(G)\le C_3\varepsilon(r)^2$.

\emph{Step 4.} By the Cauchy--Schwarz inequality for the nonnegative form
$\mathfrak a_{rT'}+M_0\|\cdot\|^2$,
$|\mathfrak a_{rT'}(u,G)|\le(\mathfrak a_{rT'}(u)+M_0)^{1/2}(\mathfrak a_{rT'}(G)+M_0\|G\|^2)^{1/2}+M_0\|G\|$.
Hence, if $\mathfrak a_{T'}(F)\le V$ for a fixed $V$, then
$\mathfrak a_{rT'}(v)\le\mathfrak a_{T'}(F)+\omega(r)$ with $\omega(r)\to0$ uniformly in such $F$. Since
also $\mathfrak a_{rT'}(v)\ge-M_0\|v\|^2$ and $\big|\|v\|-1\big|\le C_2\varepsilon(r)$, the Rayleigh
quotient gives $\lambda_{\mathrm A}(rT')\le\mathfrak a_{T'}(F)+\omega'(r)$ with
$\omega'(r)\to0$.

\emph{Step 5.} By Lemma~\ref{lem:monoT}, $\lamA$ is bounded on $[T/2,2T]$ between $C_0-\mathsf A_{2T}$
and $\lambda_{\mathrm A}(T/2)$; take $V$ one more than the latter. Applying Step 4 to
near-minimizers at $T$ gives $\limsup_{r\to1}\lambda_{\mathrm A}(rT)\le\lamA$, and applying
it to near-minimizers at $rT$ with the factor $1/r$ gives the reverse inequality for the $\liminf$.
\end{proof}


\section{Primal and dual programs}\label{sec:duality}

\subsection{Primal and dual}
Fix $T$ and the weights $v=(v_n)_{n\in\NT}$ (any nonnegative weights may be used; $v_n=2\Lambda(n)/\sqrt n$
is the case of $\zeta$).

\begin{definition}\label{def:PD}
\textbf{(P)} $\ \lamP=\inf\big\{\int_\R\sT\,d\mu:\ \mu$ an even probability measure on $\R$,
$\check\mu(u)=0$ for $|u|\ge2T$, $\int\gh0\,d\mu=\int\gh1\,d\mu=0\big\}$.

\textbf{(D)} $\ D_T=\sup\big\{\inf_{\tau\in\R}\big[\sT(\tau)+\nuh(\tau)+\beta_0\gh0(\tau)+\beta_1\gh1(\tau)\big]:\
\nu$ a finite even signed Borel measure with $\supp\nu\subseteq\{|u|\ge2T\}$, $\beta\in\R^2\big\}$.
A triple $(\nu,\beta_0,\beta_1)$ with $f:=\sT+\nuh+\beta\cdot\widehat g\ge c$ on $\R$ is a
\emph{certificate of value $c$}.
\end{definition}

\begin{proposition}[the primal is the eigenvalue problem; weak duality]\label{prop:weak}
The value of \textup{(P)} is the $\lamP$ of Theorem~\ref{thm:A}. Every certificate of value $c$
satisfies $c\le\lamP\le\lamA$.
\end{proposition}
\begin{proof}
If $F$ is a unit vector, $\mu_F=|\Fh|^2d\tau/2\pi$ is feasible ($\check\mu_F=\cF$) with
$\int\sT d\mu_F=\mathfrak a_T(F)$, and $\int\gh j\,d\mu_F=\int g_j\cF$. Conversely if $\mu$ is
feasible, $\Phi=\check\mu$ is continuous, positive definite, $\Phi(0)=1$, supported in $[-2T,2T]$,
so by the Boas--Kac theorem \cite{BoasKac1945} $\Phi=\cF$ for some $F\in L^2(\IT)$, $\|F\|=1$, and
$\Phi\in\Kp$; thus the feasible set of (P) is exactly $\Kp$ read through $\mu=|\Fh|^2d\tau/2\pi$.
Weak duality: for feasible $\mu$, $\int\nuh\,d\mu=\int\check\mu\,d\nu=0$ because $\check\mu=\cF$
vanishes identically on the closed set $\{|u|\ge2T\}$ ($\cF(\pm2T)=0$ by continuity, so an atom
of $\nu$ at $\pm2T$ is harmless), and $\int\gh j\,d\mu=0$; hence
$\int\sT d\mu=\int f\,d\mu\ge c$.
\end{proof}

\begin{lemma}[bounds on the multipliers]\label{lem:gbounds}
$\gh0(\tau)=2\int_0^{2T}\cosh\frac u2\cos\tau u\,du$ and $\gh1(\tau)=2\int_0^{2T}u\sinh\frac u2\cos\tau u\,du$
are even, real, bounded, tend to $0$ at infinity, and for $\tau>0$
\[
|\gh0(\tau)|\le k_0(\tau):=\frac{2\cosh T}{\tau}+\frac{\sinh T}{\tau^2},\qquad
|\gh1(\tau)|\le k_1(\tau):=\frac{8T\sinh T}{\tau},
\]
both bounds decreasing in $\tau$. Moreover $\gh0(\tau)=\frac{2\cosh T\sin2T\tau}{\tau}+O(\tau^{-2})$,
and for any $g=y_0g_0+y_1g_1$, $\widehat g(\tau)=\frac{2g(2T^-)\sin2T\tau}{\tau}+\frac{2g'(2T^-)\cos2T\tau}{\tau^2}+O(\tau^{-3})$.
\end{lemma}
\begin{proof}
Integrate by parts: $\gh0=\frac{2\cosh T\sin2T\tau}\tau-\frac1\tau\int_0^{2T}\sinh\frac u2\sin\tau u\,du$,
and once more, $\frac1\tau\int_0^{2T}\sinh\frac u2\sin\tau u\,du=-\frac{\sinh T\cos2T\tau}{\tau^2}+\frac{\gh0(\tau)}{4\tau^2}$;
solving, $|\gh0|(1+\frac1{4\tau^2})\le\frac{2\cosh T}\tau+\frac{\sinh T}{\tau^2}$. For $\gh1$, one
integration by parts gives $\gh1=\frac{4T\sinh T\sin2T\tau}{\tau}-\frac2\tau\int_0^{2T}\big(\sinh\frac u2+\frac u2\cosh\frac u2\big)\sin\tau u\,du$,
and the integrand is the derivative of $u\sinh\frac u2$, whose integral over $[0,2T]$ is $2T\sinh T$.
The asymptotic expansion is two integrations by parts for a function real-analytic on $[0,2T]$
with $g'(0)=0$.
\end{proof}

\subsection{Strong duality}

\begin{theorem}[strong duality; finitely atomic dual sequences]\label{thm:D1}
$D_T=\lamP$. Moreover, for every $\varepsilon>0$ there are finitely many atoms
$u_1,\dots,u_N\ge2T$, reals $a_1,\dots,a_N$ and $\beta\in\R^2$ such that
$\sT(\tau)+\sum_j2a_j\cos(u_j\tau)+\beta_0\gh0(\tau)+\beta_1\gh1(\tau)>\lamP-\varepsilon$ for all $\tau\in\R$.
In particular, for every $T$ with $\lamP>0$, $\lamA>0$ is provable by a finite certificate.
\end{theorem}
\begin{proof}
\emph{(i) Finite relaxations.} For a finite set $U=\{u_1<\dots<u_N\}\subset[2T,\infty)$ let
$P(U)=\inf\{\int\sT d\mu:\ \mu\ \text{even probability},\ \check\mu(u_i)=0\ \forall i,\ \int\gh j d\mu=0\}$
and $D(U)$ the supremum in (D) restricted to $\nu$ supported on $\pm U$. Then $P(U)\le\lamP$
(the feasible set of (P) is contained in that of $P(U)$) and $D(U)\le D_T\le\lamP$.

\emph{(ii) $P(U)=D(U)$.} Write $f_0\equiv1$, $f_i(\tau)=\cos(u_i\tau)$ ($1\le i\le N$),
$f_{N+1}=\gh0$, $f_{N+2}=\gh1$, $b=(1,0,\dots,0)\in\R^{N+3}$, and let
$C=\{\sum_kt_kf(\tau_k):\ t_k\ge0\}$ be the convex cone generated by the moment curve
$\tau\mapsto f(\tau)=(f_0,\dots,f_{N+2})(\tau)$. The value function
$w(b')=\inf\{\int\sT d\mu:\ \mu\ge0,\ \int f\,d\mu=b'\}$ is convex, $>-\infty$ everywhere
($\sT\ge-\Aconst-|C_0|$ and $\int f_0d\mu=b'_0$ bounds the mass) and finite on $C$ (finite atomic
$\mu$ have finite value). Its Legendre biconjugate at $b$ is the Lagrangian dual
$w^{**}(b)=\sup\{y_0:\ \sT-\sum_ky_kf_k\ge0\text{ on }\R\}$, which is $D(U)$ after
$a_i=-y_i/2$, $\beta_j=-y_{N+1+j}$. A proper convex function agrees with its biconjugate at every
interior point of its domain \cite[Thm.~7.4, Cor.~12.1.1]{Rockafellar}; so it suffices to show
$b\in\operatorname{int}C$.

Suppose not. Since $C$ is a convex cone containing $0$, there is $y\ne0$ with
$\langle y,z\rangle\ge\langle y,b\rangle$ for all $z\in C$; taking $z=0$ gives $y_0\le0$, and
taking $z=tf(\tau)$, $t\to\infty$, gives
$G(\tau):=y_0+\sum_iy_i\cos(u_i\tau)+y_{N+1}\gh0(\tau)+y_{N+2}\gh1(\tau)\ge0$ for all $\tau$.
The Ces\`aro mean $X^{-1}\int_0^XG\to y_0$ (each $\cos(u_i\tau)$ has mean $0$ because $u_i>0$,
and $\gh j\to0$ by Lemma~\ref{lem:gbounds}); as $G\ge0$ this forces $y_0=0$, and then
$X^{-1}\int_0^X|G|\to0$. Testing against the bounded functions $\cos(u_i\tau)$ gives
$y_i/2=\lim X^{-1}\int_0^XG\cos(u_i\tau)\,d\tau=0$ for each $i$. There remains
$\widehat g\ge0$ with $g=y_{N+1}g_0+y_{N+2}g_1$. By the expansion in Lemma~\ref{lem:gbounds},
$\widehat g\ge0$ forces $g(2T^-)=0$ and then $g'(2T^-)=0$; but
$g(2T^-)=y_{N+1}\cosh T+2T\,y_{N+2}\sinh T$ and $g'(2T^-)=\frac12y_{N+1}\sinh T+y_{N+2}(\sinh T+T\cosh T)$,
whose determinant is $\sinh T\cosh T+T>0$. So $y=0$, a contradiction; hence $b\in\operatorname{int}C$.

\emph{(iii) The limit.} Choose finite sets $U_1\subset U_2\subset\cdots$ with $\bigcup U_n$ dense in
$[2T,\infty)$; $P(U_n)$ is nondecreasing and $\le\lamP$. Pick feasible $\mu_n$ with
$\int\sT d\mu_n\le P(U_n)+1/n$. Since $\sT\to+\infty$ and is bounded below, $(\mu_n)$ is tight;
let $\mu$ be a weak-$*$ limit point, a probability measure. For $u\in U_N$ and $n\ge N$,
$\check\mu_n(u)=0$ and $\cos(u\,\cdot)$ is bounded continuous, so $\check\mu(u)=0$; $\check\mu$
is continuous, so $\check\mu\equiv0$ on $[2T,\infty)$; likewise $\int\gh jd\mu=0$. Thus $\mu$ is
feasible for (P), and by lower semicontinuity of $\mu\mapsto\int\sT d\mu$ under weak convergence
($\sT$ continuous and bounded below), $\lamP\le\int\sT d\mu\le\liminf\int\sT d\mu_n=\lim P(U_n)\le\lamP$.
Hence $D_T\ge\sup_nD(U_n)=\sup_nP(U_n)=\lamP\ge D_T$, and the atomic statement follows.
\end{proof}

\begin{remark}\label{rem:D1complete}
Theorem~\ref{thm:D1} says that, at a fixed window, $\lamP\ge0$ holds if and only if for every
$\varepsilon>0$ there are an atomic $\nu$ on $\{|u|\ge2T\}$ and $\beta\in\R^2$ with
$\sT+\nuh+\beta\cdot\widehat g\ge-\varepsilon$ on $\R$; and $\lamP\ge0$ implies $\lamA\ge0$. The same
proof, with the two $\widehat g$-constraints deleted, gives strong duality for the complete form:
$\lamW=\sup\{\inf_\tau[\sT+2\gh0+\nuh]\}$ over the same class of $\nu$ (Proposition~\ref{prop:P}
explains the constant $2$).
\end{remark}

\subsection{Non-attainment}

\begin{theorem}[non-attainment]\label{thm:O}
Let $\nu$ be any finite even signed measure with $\supp\nu\subseteq\{|u|\ge2T\}$ and let
$h_0,\dots,h_J$ be finitely many continuous functions on $[-2T,2T]$ that annihilate $\Kp$
(for instance $g_0,g_1$). Then for every $\beta\in\R^{J+1}$
\[
\inf_{\tau\in\R}\Big[\sT(\tau)+\nuh(\tau)+\sum_j\beta_j\widehat h_j(\tau)\Big]<\lamP\quad\text{strictly.}
\]
The statement holds for every finite weight family, including $v=0$.
\end{theorem}
\begin{proof}
A minimizer $\Phi_\star=c_{F_\star}\in\Kp$ of the primal exists: the sublevel sets of
$\mathfrak a_T$ on the unit sphere of $L^2(\IT)$ are compact (Lemma~\ref{lem:compact}),
and, in the language of measures, the feasible set of (P) is tight and weak-$*$ closed and $\mu\mapsto\int\sT d\mu$ is
weak-$*$ lower semicontinuous, exactly as in step (iii) of the proof of Theorem~\ref{thm:D1}. Put $\varrho=\sT+\nuh+\sum\beta_j\widehat h_j-\lamP$, a
continuous function ($\nu$ finite). By the annihilation properties,
$\frac1{2\pi}\int\varrho|\Fh_\star|^2=\mathfrak a_T(F_\star)-\lamP=0$. If $\varrho\ge0$ then $\varrho=0$ a.e.\ on
$\{\Fh_\star\ne0\}$, whose complement is the discrete zero set of a nonzero entire function of
exponential type; so $\varrho\equiv0$ on $\R$, and by Fourier inversion, as distributions in $u$,
\[
\check\sT+\nu+\sum_j\beta_jh_j=\lamP\,\delta_0 .
\]
Restrict to the open set $0<|u|<2T$, where $\nu=0$: by \eqref{eq:m}, $\check m=-\frac12K(|u|)$
there, and the prime part of $\check\sT$ is $-\sum_{n\in\NT}\frac{v_n}2(\delta_{\log n}+\delta_{-\log n})$;
hence $\sum_j\beta_jh_j(u)=\frac12K(|u|)+\sum_{n\in\NT}\frac{v_n}2(\delta_{\log n}+\delta_{-\log n})$
on $0<|u|<2T$. The left side is a bounded continuous function; the right side is unbounded as
$u\to0$ ($K(u)\sim1/u$) and carries atoms as soon as $\NT\neq\emptyset$. Contradiction.
\end{proof}

\begin{remark}\label{rem:OvsB}
Theorem~\ref{thm:O} forbids attainment, not certificates: whenever $\lamP>0$ any dual value in
$(0,\lamP)$ suffices, and Theorem~\ref{thm:D1} says such values are reached by finite atomic
$\nu$. Every maximizing sequence of the dual must therefore degenerate (its total variation is
unbounded; Corollary~\ref{cor:T2} below quantifies the cost). On the primal side the infimum is
attained \cite[Thm.~3]{Bombieri2000}. The proof shows the reason for the asymmetry: the
singularity $K(u)\sim1/u$ of the archimedean L\'evy density at $u=0$ and the prime atoms lie inside
the window, where the dual measure is not allowed to have support.
\end{remark}

\subsection{The polar multiplier}

\begin{lemma}[moment interior]\label{lem:W}
For $F\in L^2(\IT)$ real put $Q_0(F)=\int\cF\cosh\frac u2=A B$ and
$Q_1(F)=\int\cF\,u\sinh\frac u2=A'B-AB'$. The image of the set of unit vectors under
$F\mapsto(Q_0,Q_1)$ contains a neighborhood of $0$ in $\R^2$, realised by $C_c^\infty(\IT)$
functions; equivalently, $0$ is an interior point of the (convex) moment image
$\mathfrak m(\mu)=(\int\gh0d\mu,\int\gh1d\mu)$ of the feasible set of (P) without the two
constraints, and the interior is realised by measures of finite value.
\end{lemma}
\begin{proof}
The formulas for $Q_0,Q_1$ follow from expanding $\cosh\frac{t-s}2$ and $(t-s)\sinh\frac{t-s}2$
in exponentials. The functions $e^{t/2},e^{-t/2},te^{t/2},te^{-t/2}$ are linearly independent
on $\IT$, and a nontrivial combination is real-analytic, so it cannot annihilate
$C_c^\infty(\IT)$: the linear map $F\mapsto(A,B,A',B')$ is onto $\R^4$ from $C_c^\infty(\IT)$.
If $0$ were not interior, some $(t_0,t_1)\ne0$ would satisfy $t_0Q_0(F)+t_1Q_1(F)\ge0$ for all
$F$ (normalization only rescales $Q_j$ by a positive factor). Take $(A,B,A',B')=(1,0,0,\mp1)$:
$Q_0=0$, $Q_1=\pm1$, so $t_1=0$; take $(1,\pm1,0,0)$: $Q_1=0$, $Q_0=\pm1$, so $t_0=0$. A finite
convex combination of the four witnesses covers a ball, and each has finite value because
$\Fh$ is Schwartz.
\end{proof}

\begin{proposition}[the polar multiplier]\label{prop:P}
For every $F\in L^2(\IT)$ in the form domain,
$\frac1{2\pi}\int\gh0|\Fh|^2=\Rel\big(A(F)\overline{B(F)}\big)$ and hence
\[
\frac1{2\pi}\int\big(\sT+2\gh0\big)|\Fh|^2\,d\tau=W_T(F);
\]
if moreover $F$ is real and $\sum_\rho|\Phi_F(\rho)\Phi_F(1-\rho)|<\infty$ (in particular if
$F\in C_c^\infty(\IT)$ is real), then $W_T(F)=\sum_\rho\Phi_F(\rho)\Phi_F(1-\rho)$
(Lemma~\ref{lem:roughEF} below).
Consequently a certificate with $(\beta_0,\beta_1)=(2,0)$ and value $c$ proves $W_T\ge c$ on all
of $L^2(\IT)$, and $\beta_0=2$ is the unique value of $\beta_0$ with this property.
\end{proposition}
\begin{proof}
$\frac1{2\pi}\int\gh0|\Fh|^2=\int g_0\cF=\iint\cosh\frac{t-s}2F(t)\overline{F(s)}$ and
$\cosh\frac{t-s}2=\frac12e^{t/2}e^{-s/2}+\frac12e^{-t/2}e^{s/2}$ gives $\frac12(A\bar B+B\bar A)$.
The identification of $W_T$ with the zero sum is the explicit formula for the autocorrelation
$\cF$, \cite{Weil52,Bombieri2000}, \cite[Thm.~8.4]{Paper45}, whose polar term is
$2\Phi_F(0)\Phi_F(1)=2\Rel(A\bar B)$ for the real forms considered here. Uniqueness: for
$\beta_0\ne2$, $\frac1{2\pi}\int(\sT+\beta_0\gh0)|\Fh|^2=W_T(F)+(\beta_0-2)\Rel(A\bar B)$, and
$\Rel(A\bar B)$ takes both signs on the unit sphere of $L^2(\IT)$ (Lemma~\ref{lem:W}, witnesses
$(1,\pm1,0,0)$), so no lower bound of the left side bounds $W_T$; the same argument with
$\Rel(A'\bar B-A\bar B')$ forces $\beta_1=0$.
\end{proof}

\cite[Thm.~8.4]{Paper45} is stated for real $F\in C_c^\infty$. The first equality of
Proposition~\ref{prop:P} holds on the whole form domain and is the only one the certificates use;
wherever the zero sum is used for a rough test (Theorem~\ref{thm:margin}), it enters through the
following lemma.

\begin{lemma}[the explicit formula for rough tests]\label{lem:roughEF}
Let $F\in L^2(\IT)$ be real, with \emph{(a)} $\int\log(2+|\tau|)|\Fh|^2<\infty$ and \emph{(b)}
$\sum_\rho|\Phi_F(\rho)\Phi_F(1-\rho)|<\infty$. Then
\[
\sum_\rho\Phi_F(\rho)\Phi_F(1-\rho)=W_T(F)=2\Phi_F(0)\Phi_F(1)+\mathfrak a_T(F),
\]
the sum being absolutely convergent. In particular $\mathfrak a_T(F)=\sum_\rho\Phi_F(\rho)\Phi_F(1-\rho)$
for $F\in\PT$, and a zero $\rho=\frac12+i\gamma$ on the critical line contributes $|\Fh(\gamma)|^2$.
\end{lemma}
\begin{proof}
For real $F$, $\Phi_F(1)=A(F)$ and $\Phi_F(0)=B(F)$, so $2\Phi_F(0)\Phi_F(1)=2\Rel(A\bar B)$. Let
$\eta_\varepsilon\ge0$, $\eta_\varepsilon\in C_c^\infty(-\varepsilon,\varepsilon)$, $\int\eta_\varepsilon=1$,
and $F_\varepsilon=F*\eta_\varepsilon\in C_c^\infty$, supported in $[-T-\varepsilon,T+\varepsilon]$.
Apply \cite[Thm.~8.4]{Paper45} to $F_\varepsilon$ at the window $T+\varepsilon$. Zero side:
$\Phi_{F_\varepsilon}=\Phi_F\Phi_{\eta_\varepsilon}$, $|\Phi_{\eta_\varepsilon}(s)|\le e^{\varepsilon|\Rel s-\frac12|}\le e^{\varepsilon/2}$
in the critical strip and $\Phi_{\eta_\varepsilon}\to1$ pointwise, so by (b) and dominated convergence
the zero sum (which for $F_\varepsilon$ converges absolutely, so that the symmetric order is
immaterial) tends to $\sum_\rho\Phi_F(\rho)\Phi_F(1-\rho)$. Polar term:
$\Phi_{F_\varepsilon}(0)\Phi_{F_\varepsilon}(1)\to\Phi_F(0)\Phi_F(1)$. Archimedean term:
$|\Fh_\varepsilon|^2=|\Fh|^2|\widehat{\eta_\varepsilon}|^2\le|\Fh|^2$ and $|m(\tau)|\le|C_0|+\log(1+|\tau|)$
(Lemma~\ref{lem:mbounds}(i),(ii); $m\le0$ for $|\tau|\le\tau_0$), so by (a) and dominated convergence
$\frac1{2\pi}\int m|\Fh_\varepsilon|^2\to\frac1{2\pi}\int m|\Fh|^2$. Prime term: $c_{F_\varepsilon}\to\cF$
uniformly, the sum runs over the finitely many $n<e^{2T+2\varepsilon}$, and for
$e^{2T}\le n<e^{2T+2\varepsilon}$ one has $\cF(\log n)=0$; so it tends to $-\sum_{n\in\NT}v_n\cF(\log n)$.
\end{proof}


\section{Tail lemmas and an explicit certificate}\label{sec:tails}

A certificate is an inequality $f\ge c$ on all of $\R$. The lemmas of this section reduce it to a
bounded interval; they rest on Lemmas~\ref{lem:mbounds} and \ref{lem:gbounds}.

\subsection{Lattice atoms: the graded periodic tail}

\begin{lemma}[graded periodic tail]\label{lem:F1}
Let $d=2T/M$, $u_j=jd$ for $M\le j\le J$ (so $u_M=2T$), $\nu=\sum_ja_j(\delta_{u_j}+\delta_{-u_j})$,
$P=2\pi/d$, so that $\nuh(\tau)=\sum_j2a_j\cos(u_j\tau)$ is $P$-periodic. Let
$f=\sT+\nuh+\beta_0\gh0+\beta_1\gh1$. Fix $\tau_1>0$, $N\ge1$, and for $x\in[\tau_1,\tau_1+P)$ put
\begin{align*}
Y_k(x)&=m(x+kP)-\sum_{n\in\NT}v_n\cos\big((x+kP)\log n\big)-|\beta_0|k_0(x+kP)-|\beta_1|k_1(x+kP)\ \ (k<N),\\
Y_N(x)&=m(x+NP)-\Aconst-|\beta_0|k_0(x+NP)-|\beta_1|k_1(x+NP).
\end{align*}
If $\nuh(x)+Y_k(x)\ge c$ for all $x\in[\tau_1,\tau_1+P)$ and all $0\le k\le N$, then $f(\tau)\ge c$
for all $\tau\ge\tau_1$.
\end{lemma}
\begin{proof}
Every $\tau\ge\tau_1$ is uniquely $x+kP$ with $x\in[\tau_1,\tau_1+P)$, $k\ge0$, and
$\nuh(\tau)=\nuh(x)$ since $u_jP=2\pi j$. For $k<N$ the $k$-th row is $f(x+kP)$ with the exact
prime term and the two multiplier terms replaced by the lower bounds of Lemma~\ref{lem:gbounds}.
For $k\ge N$: $m(x+kP)\ge m(x+NP)$, $k_i(x+kP)\le k_i(x+NP)$, and the prime term is $\ge-\Aconst$.
\end{proof}

The prime term is never required to be periodic: it is evaluated exactly on the first $N$
copies of one period of $\nuh$ and replaced by its infimum beyond. The lemma therefore holds
for any finite set of contacts; the incommensurability of $\log2$ and $\log3$ plays no role.
It is a sufficient criterion only (the rows are conservative in $\beta$ by
$|\beta_0|k_0+|\beta_1|k_1$).

\subsection{Exponential columns: the monotone tail}

\begin{proposition}[exponential columns]\label{prop:T0}
For $s>0$ and an integer $p\ge0$ let $d\nu_{s,p}(u)=\frac12(|u|-2T)^pe^{-s(|u|-2T)}\mathbf 1_{|u|\ge2T}du$,
a positive finite measure of mass $p!/s^{p+1}$ supported outside the window. Then
\[
\nuh_{s,p}(\tau)=\Rel\Big[e^{2iT\tau}\frac{p!}{(s-i\tau)^{p+1}}\Big],\qquad
|\nuh_{s,p}(\tau)|\le\operatorname{cap}_{s,p}(\tau):=\frac{p!}{(s^2+\tau^2)^{(p+1)/2}},
\]
and $\operatorname{cap}_{s,p}$ is decreasing in $\tau\ge0$.
\end{proposition}
\begin{proof}
$\nuh_{s,p}(\tau)=\int_{2T}^\infty(u-2T)^pe^{-s(u-2T)}\cos\tau u\,du=\Rel\big[e^{2iT\tau}\int_0^\infty t^pe^{-(s-i\tau)t}dt\big]$.
\end{proof}

\begin{lemma}[monotone tail]\label{lem:T}
Let $\nu=\sum_{k=1}^Kc_k\nu_{s_k,p_k}+\sum_{j=1}^Ja_j(\delta_{u_j}+\delta_{-u_j})$ with $u_j\ge2T$, and
$f=\sT+\nuh+\beta_0\gh0+\beta_1\gh1$. Define for $\tau>0$
\[
L(\tau)=\log\frac\tau{2\pi}-\frac1\tau-\frac1{2\tau^2}-\Aconst-\sum_k|c_k|\operatorname{cap}_{s_k,p_k}(\tau)
-2\sum_j|a_j|-|\beta_0|k_0(\tau)-|\beta_1|k_1(\tau).
\]
Then $f\ge L$ on $(0,\infty)$, $L$ is nondecreasing, and $L(\tau_1)\ge c$ for one $\tau_1$ implies
$f\ge c$ on $[\tau_1,\infty)$.
\end{lemma}
\begin{proof}
Lemma~\ref{lem:mbounds}(ii),(iii), Proposition~\ref{prop:T0}, the trivial bound on the atoms and
Lemma~\ref{lem:gbounds}; every term of $L$ is nondecreasing.
\end{proof}

\begin{corollary}[cost of the tail]\label{cor:T2}
For an atom-free $\nu$ with $A:=\sum_k|c_k|p_k!$ and $A':=|\beta_0|(2\cosh T+\sinh T)+|\beta_1|8T\sinh T$,
the tail inequality $L(\tau_1)\ge c$ holds for every $\tau_1\ge\max\{A+A'+2,\ 2\pi e^{c+\Aconst+1}\}$,
so the tail costs linearly in the $\ell^1$ size of the certificate. For atoms it costs
exponentially, $\tau_1\approx2\pi e^{\Aconst+2\sum|a_j|}$, which is why Lemma~\ref{lem:F1} and not
Lemma~\ref{lem:T} is used for lattice certificates. Neither cost is uniform over the dual class:
by Theorem~\ref{thm:O}, $\|\nu\|\to\infty$ along every maximizing sequence.
\end{corollary}
\begin{proof}
For $\tau\ge1$, $\operatorname{cap}_{s,p}(\tau)\le p!/\tau$ and $|\beta_i|k_i\le A'/\tau$, so
$L(\tau)\ge\log\frac\tau{2\pi}-\Aconst-\frac{A+A'+2}{\tau}$.
\end{proof}

\begin{lemma}[hybrid tail]\label{lem:hybrid}
Let $\nu=\nu_{\mathrm{lat}}+\sum_\ell c_\ell\nu_{s_\ell,p_\ell}$ with $\nu_{\mathrm{lat}}$ as in Lemma~\ref{lem:F1}.
For $x\in[\tau_1,\tau_1+P)$ define $H_k(x)=\nuh_{\mathrm{lat}}(x)+Y_k(x)+\sum_\ell c_\ell\nuh_{s_\ell,p_\ell}(x+kP)$
for $0\le k<N$, and $H_N(x)=\nuh_{\mathrm{lat}}(x)+Y_N(x)-\sum_\ell|c_\ell|\operatorname{cap}_{s_\ell,p_\ell}(x+NP)$.
If $H_k\ge c$ on $[\tau_1,\tau_1+P)$ for all $0\le k\le N$ then $f\ge c$ on $[\tau_1,\infty)$.
\end{lemma}
\begin{proof}
As in Lemma~\ref{lem:F1}; on the rows $k<N$ the column term is evaluated exactly, and for
$k\ge N$ it is bounded by $-\sum_\ell|c_\ell|\operatorname{cap}(x+kP)\ge-\sum_\ell|c_\ell|\operatorname{cap}(x+NP)$
since $\operatorname{cap}$ decreases.
\end{proof}

\subsection{Exhaustion by exponential columns}

\begin{theorem}[exponential columns exhaust the dual]\label{thm:E}
Let $\mathcal E$ be the real span of $\{\nu_{s,p}:s>0,p\ge0\}$. Then
$\sup\{\inf_\tau[\sT+\nuh+\beta\cdot\widehat g]:\ \nu\in\mathcal E,\beta\in\R^2\}=\lamP$, and the same
holds with $\mathcal E$ replaced by the absolutely continuous measures with bounded density and
compact support in $\{|u|\ge2T\}$.
\end{theorem}
\begin{proof}
``$\le$'' is weak duality. For ``$\ge$'' fix $\varepsilon>0$ and, by Theorem~\ref{thm:D1}, an atomic
$\nu_0=\sum_{j\le J}a_j(\delta_{u_j}+\delta_{-u_j})$ and $\beta$ with
$c_0:=\inf[\sT+\nuh_0+\beta\cdot\widehat g]>\lamP-\varepsilon/2$. By Lemma~\ref{lem:T} there is
$\tau_1<\infty$ with $\sT+\nuh_0+\beta\cdot\widehat g\ge c_0-\varepsilon/8$ on $[\tau_1,\infty)$,
and the lemma uses only $\sum|a_j|$, so the same bound holds for every signed measure of the
same total variation supported in $\{|u|\ge2T\}$. Replace each atom $a_j\delta_{u_j}$ by
$a_j$ times the uniform probability density on $(u_j,u_j+\eta)$: since
$|\cos u\tau-\cos u_j\tau|\le|\tau|\eta$, the transform changes by at most $2\sum|a_j|\tau_1\eta$
on $[0,\tau_1]$, which is $<\varepsilon/8$ for small $\eta$; the new $\nu_1$ has bounded density
and compact support and value $\ge c_0-\varepsilon/4$. Finally $\{t^pe^{-t}\}_{p\ge0}$ is total in
$L^1(0,\infty)$ (if $g\in L^\infty$ annihilates them, its Laplace transform is holomorphic on
$\Rel z>0$ with all derivatives vanishing at $z=1$, hence $g=0$), so there is $\nu_2\in\mathcal E$
with $\|\nu_2-\nu_1\|_{\TV}<\varepsilon/8$, and $\|\nuh_2-\nuh_1\|_\infty\le\|\nu_2-\nu_1\|_{\TV}$ on all of $\R$.
\end{proof}

\subsection{An explicit certificate at \texorpdfstring{$T=\tfrac12\log3$}{T = (1/2) log 3}}\label{sec:cert}

At $T=\frac12\log3$ the window is $[-\log3,\log3]$, the contact set is $\NT=\{2\}$, and
$v_2=2\log2/\sqrt2=0.98025814\ldots=\Aconst$. Positivity of the complete form here is known by
\cite[Thm.~1.2]{Zhu2026}, and positivity of both forms also follows from Theorems~\ref{thm:log5A} and
\ref{thm:log5W} below. The point of the following theorem is the method: a proof by an explicit dual
feasible point, with a margin about $9.4\%$ of the numerical value $\lambda_{\mathrm A}(\frac12\log3)\approx9.448\cdot10^{-4}$.

\begin{theorem}[explicit certificate]\label{thm:C}
Let $T=\frac12\log3$, $d=2T/24=\frac{\log3}{24}$,
$u_j=jd$ for $24\le j\le460$ (so $u_{24}=\log3$ and $u_{460}=21.0567\ldots$), let
$a_{24},\dots,a_{460}$ be the $437$ rational numbers given as exact decimal strings in the
certificate data accompanying this paper, and let
$\beta_0=2.00256170925026\ldots$, $\beta_1=-0.00965443404485\ldots$. Put
$\nu=\sum_ja_j(\delta_{u_j}+\delta_{-u_j})$, so $\|\nu\|_{\TV}=2\sum|a_j|=90.75$, $\nuh(0)=1.7337$.
Then
\[
\sT(\tau)+\nuh(\tau)+\beta_0\gh0(\tau)+\beta_1\gh1(\tau)\ \ge\ c:=8.8577\cdot10^{-5}\qquad\text{for every }\tau\in\R,
\]
and therefore $\lamA\ge\lamP\ge8.8577\cdot10^{-5}>0$. The proof includes a computer verification
in interval arithmetic (Appendix~\ref{app:encl}).
\end{theorem}

\begin{proof}
The function $f$ is even, so it suffices to take $\tau\ge0$. Split $[0,\infty)$ as $[0,\tau_1]$
with $\tau_1=30$, where $f$ itself is bounded below, and $[\tau_1,\infty)$, where
Lemma~\ref{lem:F1} with $N=2$ and $P=2\pi/d=137.2608\ldots$ reduces the inequality to the three
rows $\nuh+Y_k\ge c$, $k=0,1,2$, on the single period $[30,167.2608]$.

Each of the four functions ($f$ on $[0,\tau_1]$ and the three rows) is covered by the exhaustive
second-order procedure of Proposition~\ref{prop:TAY} (Appendix~\ref{app:encl}). The second-derivative
bound is $M_2=6264.92$ for $f$, computed from the exact decimals; the three rows admit the smaller
constant $6262.99$, and the larger one is used for all four. The initial cell width is $0.02$, the
minimal width $w_{\min}=10^{-9}$; there is no cap on the number of cells and no cell is discarded.
The cover uses $2\,378\,593$ evaluations, leaves no unresolved interval, and consists of $20$
segments tiling the four intervals exactly.

The target in the cover is $8.8578\cdot10^{-5}$ minus a floating-point allowance
$\varepsilon_{\mathrm{fp}}=10^{-9}$. This allowance exceeds the a priori error of the double-precision
evaluator: argument rounding in the $437$ cosines contributes at most
$2\sum_j|a_j|\cdot u_{460}\cdot(\tau_1+NP)\cdot2^{-53}<10^{-10}$, the Taylor remainder at the minimal
width at most $M_2w_{\min}^2/8<10^{-15}$, and the summation of the $445$ terms together with the
evaluation of $\psi$ less than $10^{-12}$.

The constants $H_0,H_1$, $\psi'(\frac14)$, $k_0,k_1$ and the values of $m$ at the four minimizers
are enclosed in interval arithmetic at $40$ digits, $m$ through the recurrence
$\Rel\psi(z)=\Rel\psi(z+L)-\sum_{j<L}\Rel(z+j)^{-1}$ and the asymptotic series with Stieltjes'
remainder bound. Writing $h_k=\nuh+Y_k$ for the rows, the four minima $f(17.645)$, $h_0(37.815)$, $h_1$, $h_2(58.817)$
are strictly positive intervals with lower ends $\ge8.8578\cdot10^{-5}$; the smallest is
$h_0(37.8153)=8.85786\cdot10^{-5}$. The certificate data and the verification statistics accompany
the paper, and an independent replay from the decimal strings reproduces all
figures.
\end{proof}

\begin{remark}[how the certificate was found]
The atoms were placed on the lattice $u_j=j\cdot2T/M$ because Lemma~\ref{lem:F1} then controls the
tail on one period; the resolution near the window edge, rather than the range of the atoms, turned
out to be the decisive parameter. A linear program in $(a_j,\beta,c)$ with cutting planes on a
$\tau$-grid produced an approximate profile; since values on a grid are unreliable at the required
accuracy, the profile was re-optimized in $128$ blocks on a grid $10^3$ times finer. The result has three visible features:
\begin{enumerate}
\item $\beta_0=2.0026$ and $\beta_1\approx0$: the optimizer restores the two polar terms of the
explicit formula that the primitive class removes (Proposition~\ref{prop:P});
\item the weights $a_j$ alternate in sign at $70\%$ of consecutive lattice sites, decay smoothly away
from the window edge (mean $|a_j|$ from $0.27$ near $u=2T$ to $0.02$ at $u\approx40$), and show no
arithmetic structure (no $\log2$-spacing, no equal amplitudes);
\item the contact set $\{f-c<10^{-5}\}$ is a comb: a band inside $[0,13]$, just beyond the last
sign change of $\sT$ at $\tau=10.53$, plus one phase per tail row. This is the complementary
slackness behind Theorem~\ref{thm:O}: the contact set must grow as $c\uparrow\lamP$ and can never be
all of $\R$.
\end{enumerate}
\end{remark}

Zhu \cite[Thm.~1.2]{Zhu2026} combines a pointwise envelope of $\sT$ with a certified Cholesky
factorization of one $200\times200$ Legendre block; Theorem~\ref{thm:C} is instead an explicit dual
feasible point, whose ingredients (polar multiplier, outside-window measure, contact comb) are
interpretable, and whose existence for every $c<\lamP$ is guaranteed by Theorem~\ref{thm:D1}.


\section{The dual variable is the explicit formula}\label{sec:Z}

So far every certificate has been a list of numbers produced by an optimizer. This section
identifies these numbers. Up to an arbitrarily small uniform correction, every certificate of the
complete Weil form is the Weil explicit formula evaluated on a family of modulated test functions,
and its value at $\tau$ is a smoothing of the zero measure at $\tau$.

\subsection{Flat-top weights}
\begin{definition}\label{def:W}
$\Wcl$ is the set of even measurable $w:\R\to\R$ with $w\equiv1$ on $[-2T,2T]$ and
$|w(u)|\le Ce^{-(\frac12+\delta)|u|}$ for $|u|\ge2T$, for some $C,\delta>0$. $\Wcl^{\mathrm{reg}}\subset\Wcl$
consists of the $w\in C^\infty(\R)$ with compact support. For $w\in\Wcl$ define the even measure
\begin{equation}\label{eq:nuw}
d\nu_w=\Big[\tfrac12\big(1-w(u)\big)K(|u|)+2w(u)\cosh\tfrac u2\Big]\mathbf 1_{|u|\ge2T}\,du
-\sum_{n\notin\NT}\tfrac{v_n}2\,w(\log n)\,\big(\delta_{\log n}+\delta_{-\log n}\big),
\end{equation}
and $f_w:=\sT+\nuh_w+2\gh0$.
\end{definition}
The measure $\nu_w$ is finite ($K$ is integrable on $[2T,\infty)$, $|w|\cosh\frac u2$ is
integrable by the decay condition, and $\sum_nv_n|w(\log n)|\le2C\sum\Lambda(n)n^{-1-\delta}$) and
supported in $\{|u|\ge2T\}$. Hence $(\nu_w,2,0)$ is a dual feasible point, and by
Propositions~\ref{prop:weak} and \ref{prop:P}, $W_T\ge\inf f_w$ on $L^2(\IT)$.

\begin{theorem}\label{thm:Z}
\emph{(i)} For every $w\in\Wcl$ and $\tau\in\R$,
\begin{equation}\label{eq:fw}
\begin{split}
f_w(\tau)=C_0+\int_0^\infty\big(1-w(u)\cos\tau u\big)K(u)\,du&-\sum_{n\ge2}v_n\,w(\log n)\cos(\tau\log n)\\
&+2\int_\R w(u)\cosh\tfrac u2\cos(\tau u)\,du,
\end{split}
\end{equation}
which is the right-hand side of the Weil explicit formula for the even test function
$h_\tau(u)=w(u)\cos(\tau u)$.
\emph{(ii)} For $w\in\Wcl^{\mathrm{reg}}$, with $\wh(t)=\int w(u)e^{itu}du$ (an even entire function
of exponential type),
\begin{equation}\label{eq:Z}
f_w(\tau)=\sum_\rho\wh(\tau-\gamma_\rho),\qquad\rho=\tfrac12+i\gamma_\rho,
\end{equation}
the sum running over the non-trivial zeros of $\zeta$ with multiplicity, and converging absolutely.
A zero on the critical line contributes $\wh(\tau-\gamma_\rho)$ with $\gamma_\rho$ real, so on
the zeros of the critical line $f_w$ is the smoothing of the zero-counting measure by the kernel $\wh$.
\end{theorem}
\begin{proof}
(i) By \eqref{eq:m}, $m(\tau)+\int_{2T}^\infty(1-w)K\cos\tau u\,du=C_0+\int_0^{2T}(1-\cos\tau u)K
+\int_{2T}^\infty K-\int_{2T}^\infty wK\cos\tau u=C_0+\int_0^\infty(1-w\cos\tau u)K$, using $w=1$ on
$[0,2T]$. The prime terms inside and outside the window combine into $-\sum_{n\ge2}v_nw(\log n)\cos(\tau\log n)$,
again because $w(\log n)=1$ for $n\in\NT$, and $2\gh0+2\int_{|u|\ge2T}w\cosh\frac u2\cos\tau u=2\int_\R w\cosh\frac u2\cos\tau u$.
(ii) For even $h\in C_c^\infty(\R)$ the Weil explicit formula \cite{Weil52}, \cite[Thm.~5.12]{IK}, in the
normalization of this paper, reads
\[
\sum_\rho\widehat h(\gamma_\rho)=2\int_\R h(u)\cosh\tfrac u2\,du-\sum_{n\ge2}v_n\,h(\log n)+\frac1{2\pi}\int_\R m(t)\,\widehat h(t)\,dt .
\]
Apply it to $h_\tau=w\cos(\tau\,\cdot)$, whose transform is $\frac12[\wh(t-\tau)+\wh(t+\tau)]$. The
zero sum is $\sum_\rho\wh(\tau-\gamma_\rho)$, since $\wh$ is even and the multiset $\{\gamma_\rho\}$ is
invariant under $\gamma\mapsto-\gamma$ ($\rho\mapsto1-\rho$). The polar and prime terms are the
last two terms of \eqref{eq:fw}. For the archimedean term insert \eqref{eq:m}: by Fourier
inversion $\frac1{2\pi}\int\wh(t\mp\tau)\,dt=w(0)=1$ and
$\frac1{2\pi}\int\wh(t\mp\tau)\cos(tu)\,dt=w(u)\cos(\tau u)$, and Fubini is justified by
$\int|\wh(t-\tau)|\int_0^\infty\min(2,t^2u^2)K(u)\,du\,dt<\infty$; the result is
$C_0+\int_0^\infty(1-w(u)\cos\tau u)K(u)\,du$. Absolute convergence of the zero sum holds because $\wh$ is
Schwartz on horizontal lines of the strip $|\Iml z|\le\frac12$, uniformly, and $N(t)\ll t\log t$.
\end{proof}

\begin{theorem}[exhaustion by flat-top weights]\label{thm:Zexh}
Let $T\ge\frac12\log2$. \emph{(i)} $\sup_{w\in\Wcl^{\mathrm{reg}}}\inf_{\tau\in\R}f_w(\tau)=\lamW$.
\emph{(ii)} $\sup\{\inf_\tau[f_w+(\beta_0-2)\gh0+\beta_1\gh1]:\ w\in\Wcl^{\mathrm{reg}},\beta\in\R^2\}=\lamP$.
\end{theorem}
Thus every certificate of the complete Weil form is a flat-top explicit-formula certificate up to
an arbitrarily small uniform correction, and every certificate of the primitive form is one plus a
correction in the two primitivity multipliers.
\begin{proof}
``$\le$'' is weak duality in both cases. For ``$\ge$'' in (i): the proof of Theorem~\ref{thm:D1}
with the two $\widehat g$-constraints removed gives $\lamW=\sup\inf[\sT+2\gh0+\nuh]$ over finite
atomic $\nu$ outside the window, and the argument of Theorem~\ref{thm:E} (with $\beta=(2,0)$
fixed) replaces the atoms by a bounded even density $h$ with compact support in $\{|u|\ge2T\}$,
losing at most $\varepsilon/2$. It remains to realize $h\,du$ as $\nu_w$ up to total variation
$\varepsilon/2$, with $w\in\Wcl^{\mathrm{reg}}$; the transform then moves by at most $\varepsilon/2$
uniformly. On $u\ge2T$ the denominator $d_*(u)=2\cosh\frac u2-\frac12K(u)$ is positive and increasing
($d_*(0.3)>0$ and $d_*'>0$ on $(0,\infty)$), because $2T\ge\log2>0.3$; put
$w_h(u)=\big(h(u)-\tfrac12K(u)\big)/d_*(u)$ for $u\ge2T$, extended evenly and by $1$ on the window.
Then $|w_h|\ll e^{-|u|}$ and the absolutely continuous part of $\nu_{w_h}$ is exactly $h\,du$, while
its atomic part is $-\sum_{n\notin\NT}\frac{v_n}2w_h(\log n)(\delta_{\pm\log n})$, of total variation
$\ll\sum_n\Lambda(n)n^{-3/2}<\infty$. Now modify, in this order: (a) replace $h$ on $[2T,2T+\eta]$
by a smooth function matching $d_*(u)+\frac12K(u)$ to infinite order at $2T$ (so that $w_h\equiv1$
smoothly across the window edge); (b) mollify $h$ and cut $w_h$ off smoothly beyond a large $U_1$
(where $w_h\ll e^{-u}$; the a.c.\ part of $\nu_w$ becomes $\frac12K$ there, of total variation
$\ll e^{-U_1/2}$, and all atoms with $\log n>U_1$ disappear because $w=0$ there); (c) for the finitely
many $n\notin\NT$ with $\log n\le U_1$, replace $h$ by $\frac12K$ on disjoint intervals of small total
length around $\log n$ (so that $w_h$ vanishes there and the remaining atoms disappear). Each step
changes $\nu_{w_h}$ by an arbitrarily small total variation, and the result is in
$\Wcl^{\mathrm{reg}}$. (ii) is the same argument applied to the certificates of Theorem~\ref{thm:E}
with their own $\beta$.
\end{proof}

\subsection{Rigidity, smoothness, and finite zero sets}

\begin{lemma}[rigidity of flat-top weights]\label{lem:Rrig}
Let $w\in L^1(\R)$ be even and continuous with $w\equiv1$ on $[-2T,2T]$ for some $T>0$. Then
$\wh\ge0$ on $\R$ is impossible. Equivalently, if an even continuous $w\in L^1$ with $\wh\ge0$
satisfies $w\equiv1$ on an interval around the origin, then $w\equiv1$ on $\R$, which contradicts
integrability.
\end{lemma}
\begin{proof}
Suppose $w\in L^1$ is continuous with $\wh\ge0$. Then $\wh\in L^1$ and the inversion formula
$w(u)=\frac1{2\pi}\int\wh(t)e^{-itu}dt$ holds pointwise: this is the classical summability argument
(Gauss--Weierstrass or Fej\'er means converge to $w$ uniformly on compacta, and monotone convergence
applied to the nonnegative $\wh$ against the approximate identity gives $\int\wh=2\pi w(0)<\infty$;
see \cite[Ch.~VI]{Katznelson}). Hence $\wh\,dt/2\pi$ is a finite positive measure whose Fourier
transform is $w$, so by Bochner's theorem $w$ is positive definite, with $w(0)=1$. For a positive
definite $\phi$ the nonnegativity of the $3\times3$ Gram matrix at the points $0,-u,a$ gives the
standard inequality
\[
|\phi(u+a)-\phi(u)|^2\ \le\ 2\phi(0)\big(\phi(0)-\Rel\phi(a)\big)\qquad(u,a\in\R).
\]
Take any $a$ with $0<|a|\le2T$: then $w(a)=1=w(0)$, the right side vanishes, and $w(u+a)=w(u)$ for
every $u$. So $w$ is periodic with every period in $(0,2T]$; being continuous, it is constant, equal
to $w(0)=1$. A constant function is not in $L^1(\R)$.
\end{proof}

\begin{corollary}[positivity is never termwise]\label{cor:nontermwise}
For every even continuous $w\in\Wcl$ the kernel $\wh$ takes negative values. Consequently the
inequality $f_w=\sum_\rho\wh(\tau-\gamma_\rho)\ge0$ for $w\in\Wcl^{\mathrm{reg}}$ always expresses
cancellation across the zeros, never the termwise positivity of a kernel: there is no flat-top $w$
with $\wh\ge0$.
\end{corollary}
\begin{proof}
Lemma~\ref{lem:Rrig} applied to $w$, which lies in $L^1$ by the decay hypothesis of
Definition~\ref{def:W}, is even and equals $1$ on the window. A merely measurable $w$ with $\wh\ge0$
agrees almost everywhere with a continuous function, so nothing is lost.
\end{proof}

The following elementary lemma extends Theorem~\ref{thm:Z}(ii), stated for $w\in C_c^\infty$, to the weights $w\in C^k_c$ that arise in finite computations, and bounds the contribution of far zeros.

\begin{lemma}[smoothness and far zeros]\label{lem:smooth}
Let $w\in\Wcl$ be supported in $[-U,U]$, and let $k\ge0$ be an integer such that $w^{(k)}$ is
continuous on $\R$ (in particular across the window edge) and $w^{(k+1)}$ is of bounded variation,
with total variation $V=V(w^{(k+1)})$. Then:
\emph{(i)} for every $z\in\C$, $|\wh(z)|\le V\,e^{U|\Iml z|}|z|^{-k-2}$; in particular
$\wh(t)=O(|t|^{-k-2})$ on $\R$.
\emph{(ii)} The zero sum $\sum_\rho\wh(\tau-\gamma_\rho)$ converges absolutely and
Theorem~\ref{thm:Z}(ii) holds verbatim for such $w$; thus \eqref{eq:Z} is valid on
$\Wcl\cap C_c^k$ for every $k\ge0$, not only on $\Wcl^{\mathrm{reg}}$.
\emph{(iii)} Uniformly for $|\tau|\le H/2$,
\[
\sum_{|\gamma_\rho|>H}\big|\wh(\tau-\gamma_\rho)\big|\ \ll_{U}\ V\,H^{-k-1}\log H .
\]
\emph{(iv)} A zero $\rho=\frac12+i\gamma_\rho$ off the critical line, necessarily with
$|\Iml\gamma_\rho|\le\frac12$, contributes to the sum at most $e^{U/2}$ times the bound that
\emph{(i)} gives for a zero on the line with the same real ordinate.
\end{lemma}
\begin{proof}
(i) Integrating by parts $k+2$ times, the boundary terms vanishing because $w^{(j)}$ is continuous
and compactly supported for $j\le k$,
$\wh(z)=\int w(u)e^{izu}du=(iz)^{-(k+2)}\int e^{izu}\,dw^{(k+1)}(u)$, and $|e^{izu}|\le e^{U|\Iml z|}$
on $[-U,U]$. (ii) By (i) with $k\ge0$ we have $\wh(t)\ll V(1+|t|)^{-2}$ on $\R$ and
$|\wh(t+iy)|\ll Ve^{U/2}(1+|t|)^{-2}$ for $|y|\le\frac12$, while the number of zeros with
$|\gamma_\rho|\in[t,t+1]$ is $O(\log(2+t))$ by $N(t)\ll t\log t$; hence the sum converges absolutely,
and every step of the proof of Theorem~\ref{thm:Z}(ii) (the explicit formula for even $h$ of bounded
variation with compact support, and the Fubini justification, which used only
$\int|\wh(t-\tau)|\int_0^\infty\min(2,t^2u^2)K(u)du\,dt<\infty$) applies unchanged.
(iii) With $|\tau|\le H/2$ and $|\gamma_\rho|>H$ one has $|\tau-\gamma_\rho|\ge|\gamma_\rho|/2$,
so by (i) and partial summation against $N(t)=\frac t{2\pi}\log\frac t{2\pi e}+O(\log t)$,
$\sum_{|\gamma_\rho|>H}|\wh(\tau-\gamma_\rho)|\ll Ve^{U/2}\int_H^\infty t^{-k-2}\log t\,dt
\ll_UV H^{-k-1}\log H$. (iv) is the factor $e^{U|\Iml z|}$ of (i) at $|\Iml z|\le\frac12$.
\end{proof}

For $k=-1$, that is, a weight with a jump at the edge such as the box, the same computation gives
only $\wh(t)=O(|t|^{-1})$; the zero sum then converges conditionally, and (iii) becomes a statement
about oscillation rather than an absolute bound. For this reason the flat-top weights used in
practice are continuous across $u=\pm2T$.

\begin{lemma}[finite zero sets and positive definiteness]\label{lem:fdict}
Let $T>0$, let $Z\subset\R\setminus\{0\}$ be a finite symmetric multiset, and let $w:\R\to\R$ be
even, continuous and integrable with $w\equiv1$ on $[-2T,2T]$. Put
$g_Z(u)=\sum_{\gamma\in Z}e^{i\gamma u}=\sum_{\gamma\in Z}\cos(\gamma u)$ and $\phi:=w\,g_Z$. Then
$\phi$ is continuous and integrable, $\widehat\phi(\tau)=\sum_{\gamma\in Z}\wh(\tau-\gamma)$,
$\phi(0)=|Z|$, $\phi\equiv g_Z$ on $[-2T,2T]$, and
\[
\sum_{\gamma\in Z}\wh(\cdot-\gamma)\ \ge\ 0\ \text{ on }\R
\qquad\Longleftrightarrow\qquad
\phi\ \text{ is positive definite.}
\]
\end{lemma}
\begin{proof}
Continuity, $|\phi|\le|Z|\,|w|\in L^1$ and the values on the window are immediate from $w\equiv1$
there; $\widehat\phi=\wh*\mu_Z$ with $\mu_Z=\sum_{\gamma\in Z}\delta_\gamma$ finite is the
convolution theorem. ($\Leftarrow$) If $\phi$ is positive definite, its spectral measure is
nonnegative, and $\phi\in L^1\cap C$ forces that measure to be $\widehat\phi\,d\tau/2\pi$, so
$\widehat\phi\ge0$. ($\Rightarrow$) If $\widehat\phi\ge0$, the summability argument of
Lemma~\ref{lem:Rrig} (Gauss--Weierstrass means, then monotone convergence against the nonnegative
$\widehat\phi$) gives $\int\widehat\phi=2\pi\phi(0)<\infty$, so $\widehat\phi\,d\tau/2\pi$ is a
finite positive measure with Fourier transform $\phi$, and $\phi$ is positive definite by Bochner.
\end{proof}

\begin{theorem}[finite zero sets are determinate]\label{thm:finDet}
Let $T>0$, let $Z\subset\R\setminus\{0\}$ be a finite symmetric multiset with $|Z|=\ell\ge2$
(multiplicities allowed), and put $g_Z(u)=\sum_{\gamma\in Z}\cos(\gamma u)$. Let $\nu$ be a
finite positive symmetric measure on $\R$ with $\nuh=g_Z$ on $(-2T,2T)$. Then
$\nu=\mu_Z=\sum_{\gamma\in Z}\delta_\gamma$.
\end{theorem}
No decay hypothesis on $\nu$ and no product structure is assumed: the extension problem for a
finite zero set is determinate.
\begin{proof}
Since $\nu$ is finite, $\phi:=\nuh$ is continuous, bounded, real, even and positive definite.

\emph{Step 1 (a lattice inside the window).} Let $w_1,\dots,w_r$ be the distinct elements of
$Z$, with multiplicities $m_1,\dots,m_r\ge1$, so $\sum_jm_j=\ell$ and $r\le\ell$. Put $z_j=e^{iw_ja}$ for
a step $a>0$, so that $g_Z(ka)=\sum_{j\le r}m_jz_j^{\,k}$ for all $k\in\Z$. Since $Z$ is symmetric
with multiplicity, the list $(z_j,m_j)$ is closed under conjugation. Each condition $z_j=z_{j'}$
($j\ne j'$) excludes a discrete set of $a$, so let $E\subset(0,2T/\ell)$ be that exceptional discrete
set and fix $a\in(0,2T/\ell)\setminus E$, so that $z_1,\dots,z_r$ are pairwise distinct. As $r\le\ell$ we
have $ra\le\ell a<2T$, hence $\phi(ka)=g_Z(ka)$ for $|k|\le r$.

\emph{Step 2 (a rank-$r$ Toeplitz block).} Put $c_k:=\phi(ka)\in\R$ ($\phi$ is real and even) and
$A:=[c_{i-j}]_{0\le i,j\le r}$. By Step 1, $A=V^*DV$ with $V_{j,i}=z_j^{\,i}$ and
$D=\operatorname{diag}(m_1,\dots,m_r)$, $m_j>0$; $V$ has rank $r$ because the $z_j$ are distinct, so
$\operatorname{rank}A=r$ and
\[
\ker A=\Big\{v\in\C^{r+1}:\ \textstyle\sum_{i=0}^{r}v_iz_j^{\,i}=0\ (j\le r)\Big\}=\C\cdot v,
\]
$v$ being the coefficient vector of $P(z)=\prod_{j\le r}(z-z_j)$. The root set is closed under
conjugation, so $v\in\R^{r+1}$, $v_r=1$ and $|v_0|=\prod_j|z_j|=1\ne0$.

\emph{Step 3 (the recursion forced by positive definiteness).} Fix $m>r$. The Gram matrix of $\phi$
at $\{0,a,\dots,ra,ma\}$ is positive semidefinite; in block form it is
$M=\begin{psmallmatrix}A&b\\b^*&c_0\end{psmallmatrix}$ with $b_i=\phi(ia-ma)=c_{m-i}$. For
$v\in\ker A$ and $x=(v,t)$ one has $x^*Mx=2\Rel(\bar t\,b^*v)+|t|^2c_0\ge0$ for every $t\in\C$;
choosing the phase of $t$ adversely and letting $|t|\to0$ gives $b^*v=0$, i.e.
\begin{equation}\label{eq:frec}
\sum_{i=0}^{r}v_i\,c_{m-i}=0 .
\end{equation}
The sequence $k\mapsto g_Z(ka)$ satisfies the same relation: $\sum_iv_iz^{m-i}=z^{m-r}z^rP(1/z)$
vanishes at $z=z_j$ because $1/z_j=\bar z_j$ is again a root of $P$; multiplying by $m_j$ and
summing over $j$ gives $\sum_iv_ig_Z((m-i)a)=0$.

\emph{Step 4 (induction).} By Step 1, $c_k=g_Z(ka)$ for $0\le k\le r$. If $c_k=g_Z(ka)$ for all
$0\le k\le m-1$ with $m>r$, then in \eqref{eq:frec} every term with $i\ge1$ equals the
corresponding term for $g_Z$; subtracting the two relations and using $v_0\ne0$ gives
$c_m=g_Z(ma)$. Hence $\phi(ma)=g_Z(ma)$ for every integer $m\ge0$ and every
$a\in(0,2T/\ell)\setminus E$.

\emph{Step 5 (a variable step).} Fix an integer $m\ge1$. By Step 4 the identity
$\phi(ma)=g_Z(ma)$ holds for every $a\in(0,2T/\ell)$ outside the discrete set $E$; both sides are
continuous in $a$, so it holds for every $a\in(0,2T/\ell)$. Since
$\bigcup_{m\ge1}m\cdot(0,2T/\ell)=(0,\infty)$ and both functions are even, $\phi=g_Z$ on $\R$. The
spectral measure of a continuous positive definite function is unique (uniqueness of
Fourier--Stieltjes transforms), so $\nu=\mu_Z$.
\end{proof}

\begin{theorem}[no finite zero set carries a flat-top certificate]\label{thm:finZ}
Let $T>0$, let $Z\subset\R\setminus\{0\}$ be a finite symmetric multiset with $|Z|=\ell\ge2$
(multiplicities allowed), and let $w:\R\to\R$ be even, continuous and integrable with $w\equiv1$ on
$[-2T,2T]$ and $w(u)\to0$ as $|u|\to\infty$. Then
\[
f_{w,Z}(\tau):=\sum_{\gamma\in Z}\wh(\tau-\gamma)
\]
takes negative values on $\R$. In particular, for every $H>0$ and every $w\in\Wcl$ the
truncation $\sum_{|\gamma_\rho|\le H}\wh(\tau-\gamma_\rho)$ of a flat-top certificate is
not nonnegative.
\end{theorem}
\begin{proof}
Suppose $f_{w,Z}\ge0$. By Lemma~\ref{lem:fdict} the function $\phi=w\,g_Z$ is continuous and
positive definite with $\phi(0)=\ell$, and $\widehat\phi=f_{w,Z}\in L^1$ with $\int f_{w,Z}=2\pi\ell$, so
$\nu_w:=f_{w,Z}\,d\tau/2\pi$ is a finite positive symmetric measure with $\nuh_w=\phi$, which equals
$g_Z$ on $(-2T,2T)$ because $w\equiv1$ there. By Theorem~\ref{thm:finDet}, $\nu_w=\mu_Z$, i.e.
$w\,g_Z=g_Z$ on $\R$. But along a lattice $\{ma\}$ as in Step~1 of that proof, the $z_j$ are
distinct points of the unit circle, so the Ces\`aro averages of $(z_j\bar z_{j'})^m$ vanish for
$j\ne j'$ (a geometric sum) and
\[
\lim_{M\to\infty}\frac1M\sum_{m=1}^{M}|g_Z(ma)|^2=\sum_{j\le r}m_j^2\ \ge\ r\ \ge\ 1 ,
\]
whence $\limsup_m|g_Z(ma)|\ge1$, while $|w(ma)g_Z(ma)|\le\ell|w(ma)|\to0$, a contradiction. Hence
$f_{w,Z}\ge0$ is impossible; being continuous with $\int_\R f_{w,Z}=2\pi\ell>0$, it takes negative
values.
\end{proof}

\begin{corollary}[finite sets of zeros do not suffice]\label{cor:infinite}
Let $T>0$. Every flat-top certificate $f_w\ge0$, $w\in\Wcl^{\mathrm{reg}}$, is a statement about the complete,
infinite zero multiset of $\zeta$: no truncation of $\sum_\rho\wh(\tau-\gamma_\rho)$ to finitely many
zeros is nonnegative, for any window, any admissible weight and any truncation height. A finite
table of zeros can evaluate $f_w$ to any accuracy on a bounded range (Lemma~\ref{lem:smooth}(iii)),
but cannot establish $f_w\ge0$. Quantitatively, if $f_w\ge c>0$ for some $w\in\Wcl\cap C_c^k$
supported in $[-U,U]$ with $V=V(w^{(k+1)})$, and $H$ is large enough that
$C_UVH^{-k-1}\log H\le c/2$, then the truncation $f_{w,H}$ is $\ge c/2$ on
$|\tau|\le H/2$ while being negative somewhere; hence all of its negativity lies at
$|\tau|\gtrsim H/2$.
\end{corollary}
\begin{proof}
The first assertion is Theorem~\ref{thm:finZ} applied to the multiset
$Z=\{\gamma_\rho:|\gamma_\rho|\le H\}$, which is finite, symmetric and free of $0$. The
evaluation statement and the quantitative one follow from Lemma~\ref{lem:smooth}(iii), which bounds
$|f_w-f_{w,H}|$ by $C_UVH^{-k-1}\log H$ uniformly for $|\tau|\le H/2$.
\end{proof}

The mechanism is the following. A finite $Z$ makes the Toeplitz forms $[g_Z((i-j)a)]$ of finite rank;
beyond that rank, positive semidefiniteness of any extension propagates the linear recursion of $P$
indefinitely, so the extension is unique and does not decay (Theorem~\ref{thm:finDet}), whereas a
flat-top weight must decay. For the zero set of $\zeta$ no such $P$ exists and the rank is never finite.

\subsection{The decay tail of a flat-top certificate}

Proposition~\ref{prop:env} bounds the tail of any certificate at $2\pi e^{\Aconst+\|\nu\|_{\TV}}$,
and Lemma~\ref{lem:T} at a threshold exponential in the atomic mass. Both ignore the decay of
$\nuh_w$. For a $C^2$ flat-top weight this decay is available. The following theorem extends
Lemma~\ref{lem:T}: the single term $2\sum_j|a_j|$ is split into an atomic part, which does not decay,
and an absolutely continuous part, which decays like $\tau^{-2}$. It is also the smooth analogue of
Lemma~\ref{lem:F1}, with one essential difference recorded in Remark~\ref{rem:Dgrade}.

\begin{theorem}[decay tail bound for flat-top weights]\label{thm:Dtail}
Let $T>0$ and let $w\in\Wcl\cap C^2(\R)$ (so $w$ is even, $w\equiv1$ on $[-2T,2T]$, and $w'$, $w''$
match across $u=\pm2T$; in particular $w'(2T)=w''(2T)=0$), with $w''$ of bounded variation, and
assume either \emph{(a)} $\supp w\subset[-U,U]$, or \emph{(b)} $|w^{(j)}(u)|\le Ce^{-s|u|}$ for
$j\le2$ with $s>\frac12$. Put $\Omega(u)=w(|u|)\cosh\frac u2$ and $R=(1-w)K$ on $[2T,\infty)$, and
\[
D_c(w)=4\int_0^\infty\Big|\big(w(u)\cosh\tfrac u2\big)''\Big|\,du,\qquad
D_K(w)=\int_{2T}^\infty\big|\big((1-w)K\big)''(u)\big|\,du,
\]
\[
D(w)=D_c(w)+D_K(w),\qquad E_{\mathrm{at}}(w)=\sum_{\log n>2T}v_n\,|w(\log n)| ,
\]
all four finite. Then, with $\Psi_w(\tau)=\int_{2T}^\infty(1-w)K\cos\tau u\,du$,
\begin{equation}\label{eq:Dsplit}
\nuh_w(\tau)+2\gh0(\tau)=2\widehat\Omega(\tau)+\Psi_w(\tau)-\sum_{\log n>2T}v_n\,w(\log n)\cos(\tau\log n),
\end{equation}
and
\[
|2\widehat\Omega(\tau)|\le D_c(w)\,\tau^{-2},\qquad |\Psi_w(\tau)|\le D_K(w)\,\tau^{-2}.
\]
Consequently, for $f=f_w+(\beta_0-2)\gh0+\beta_1\gh1$ and every $\tau>0$,
\[
f(\tau)\ \ge\ L_w(\tau):=\log\frac\tau{2\pi}-\frac1\tau-\frac1{2\tau^2}-\Aconst-E_{\mathrm{at}}(w)
-|\beta_0-2|k_0(\tau)-|\beta_1|k_1(\tau)-D(w)\,\tau^{-2},
\]
$L_w$ is nondecreasing on $(0,\infty)$, and positivity of $f$ on $\R$ reduces to the compact
interval $[0,\tau_{\mathrm{stop}}(w)]$, $\tau_{\mathrm{stop}}(w)=\min\{\tau:L_w(\tau)\ge0\}$.
\end{theorem}
\begin{proof}
\emph{The identity \eqref{eq:Dsplit}, and why the $\cosh$ part must be regrouped with $2\gh0$.}
By \eqref{eq:nuw} the absolutely continuous part of $\nu_w$ contributes
$\Psi_w(\tau)+4\int_{2T}^\infty w\cosh\frac u2\cos\tau u\,du$ and its atomic part contributes the
last sum of \eqref{eq:Dsplit}. Since $\gh0(\tau)=2\int_0^{2T}\cosh\frac u2\cos\tau u\,du$ and
$w\equiv1$ on the window, $4\int_{2T}^\infty w\cosh\frac u2\cos\tau u\,du+2\gh0(\tau)
=4\int_0^\infty w\cosh\frac u2\cos\tau u\,du=2\widehat\Omega(\tau)$. The regrouping is necessary: the
density of $\nu_w$ at $u=2T^+$ equals $2\cosh T\ne0$, so $\nu_w$ has a jump at the window edge and
one integration by parts leaves the irreducible term $-4\cosh T\sin(2T\tau)/\tau$. Thus $\nuh_w$
alone is only $O(\tau^{-1})$, and no bound $|\nuh_w|\le D\tau^{-2}$ can hold. By
Lemma~\ref{lem:gbounds}, $2\gh0(\tau)=4\cosh T\sin(2T\tau)/\tau+O(\tau^{-2})$ cancels that term
exactly, with the same coefficient and the opposite sign; this is visible at once in
$2\widehat\Omega$, since $\Omega$ has no edge.

\emph{The two bounds.} $\Omega$ is even and, by the hypotheses, globally $C^2$ on $\R$ with $\Omega''$ of
bounded variation: on the window $\Omega=\cosh\frac u2$, and $w,w',w''$ match across $\pm2T$. In case
(a) $\Omega$ is compactly supported; in case (b) $|\Omega^{(j)}|\le C'e^{-(s-\frac12)|u|}$ for $j\le2$ is
integrable, which is exactly where $s>\frac12$ is needed. Two integrations by parts, all boundary
terms vanishing, give $\widehat\Omega(\tau)=-\tau^{-2}\int_\R e^{i\tau u}\,d\Omega'(u)$, so
$|2\widehat\Omega|\le2V(\Omega')\tau^{-2}=4\int_0^\infty|\Omega''|\,du\cdot\tau^{-2}=D_c(w)\tau^{-2}$ by evenness.
For $R$, extend by $0$ to $(-2T,2T)$ and evenly: $R(2T)=(1-w(2T))K(2T)=0$ and
$R'(2T)=-w'(2T)K(2T)+(1-w(2T))K'(2T)=0$ because $w\equiv1$ on the window and $w\in C^1$, so the
extension is $C^1$ on $\R$ with $R'$ of bounded variation, $V(R')=2\int_{2T}^\infty|R''|$ (in case
(a), $R=K$ beyond $U$, and $K,K',K''$ are integrable at infinity; in case (b),
$|R''|\le(|w''|+|w'|+|1-w|)\sup_{u\ge2T}(K+|K'|+|K''|)e^{-u/2}$). As $\Psi_w=\frac12\widehat R$,
two integrations by parts give $|\Psi_w|\le\frac12V(R')\tau^{-2}=D_K(w)\tau^{-2}$.

\emph{The atoms carry no decay}, and are bounded only by their mass:
$|\sum_{\log n>2T}v_nw(\log n)\cos(\tau\log n)|\le E_{\mathrm{at}}(w)$, the sum being finite in case
(a) and convergent in case (b) because $\sum_n\Lambda(n)n^{-\frac12-s}<\infty$ for $s>\frac12$,
the same threshold as above.

\emph{The lower bound.} $\sT(\tau)\ge\log\frac\tau{2\pi}-\frac1\tau-\frac1{2\tau^2}-\Aconst$ by
Lemma~\ref{lem:mbounds}(ii),(iii), the two multiplier terms are bounded by Lemma~\ref{lem:gbounds},
and the three pieces of \eqref{eq:Dsplit} by the above. Each term of $L_w$ is nondecreasing, hence so
is $L_w$, and $L_w(\tau_{\mathrm{stop}})\ge0$ gives $f\ge0$ on $[\tau_{\mathrm{stop}},\infty)$.
\end{proof}

\begin{remark}\label{rem:Dgrade}
The exponent $\tau^{-2}$ is the one Lemma~\ref{lem:smooth}(i) gives at $k=0$ after the two extra
derivatives, and it cannot be improved without $w\in C^3$: integrating by parts once more, a jump of
$w'''$ at the support endpoint contributes a genuine $\tau^{-2}$ to $\widehat\Omega$. The bound on
the atoms is sharp in the worst case by Lemma~\ref{lem:mbounds}(iii). This is where the method
differs from Lemma~\ref{lem:F1}: a lattice $\nuh$ is $P$-periodic, so the comb can be evaluated
exactly on $N$ periods and replaced by its infimum beyond, whereas a flat-top $\nuh_w$ is not
periodic and its comb can only be bounded by its mass. Theorem~\ref{thm:Dtail} handles the absolutely
continuous part by decay. We do not know how to combine the two techniques; the remaining
exponential dependence on the atomic mass comes from this gap.
\end{remark}

\begin{corollary}[comb-vanishing weights]\label{cor:comb}
Assume in addition $\supp w\subset[-U,U]$ and $w(\log n)=0$ for every prime power $n$ with
$2T<\log n\le U$ (a finite set of linear conditions). Then $E_{\mathrm{at}}(w)=0$, and for
$D(w)\le D_0$,
\[
\tau_{\mathrm{stop}}(w)\ \le\ \max\Big\{2\pi e^{\Aconst+1+|\beta_0-2|k_0(1)+|\beta_1|k_1(1)},\ \sqrt{D_0}+2\Big\},
\]
since for $\tau$ beyond the first entry one has $\log\frac\tau{2\pi}-\Aconst-|\beta_0-2|k_0-|\beta_1|k_1\ge1$
and then $L_w(\tau)\ge0$ as soon as $\tau^2\ge D_0+O(\tau)$. In particular
$\tau_{\mathrm{stop}}=O(D(w)^{1/2})$ as $D(w)\to\infty$: the verification interval is polynomial,
not exponential, in the size of the weight, compared with $2\pi e^{\Aconst+E_{\mathrm{at}}(w)}$ in
general and $2\pi e^{\Aconst+\|\nu\|_{\TV}}$ in Proposition~\ref{prop:env}.
\end{corollary}

\subsection{The Kre\u{\i}n extension problem at one window}\label{subsec:KEP}
The dual variable of Section~\ref{sec:Z} is a positive measure prescribed only on the window. This
is an extension problem for a positive definite datum given on an interval, in the sense of
Kre\u{\i}n \cite{Krein1940}, \cite[Ch.~2]{Sasvari1994}. In this subsection we make the identification
at one window. The classical theory concerns continuous positive
definite functions, whereas the datum below is a distribution with atoms; we therefore import
nothing from it, and every statement is proved from Theorems~\ref{thm:Z} and \ref{thm:Zexh} and
elementary Fourier analysis.

\begin{definition}[the window datum, and $\mathrm{KEP}(T)$]\label{def:KEP}
Let $g$ be the even distribution defined on even $h\in C_c^\infty(\R)$ by
\begin{equation}\label{eq:gdatum}
\langle g,h\rangle=C_0\,h(0)+\int_0^\infty\big(h(0)-h(u)\big)K(u)\,du-\sum_{n\ge2}v_n\,h(\log n)
+2\int_\R h(u)\cosh\tfrac u2\,du,
\end{equation}
that is, the right-hand side of the explicit formula in the normalization of the proof of
Theorem~\ref{thm:Z}(ii): the L\'evy pair $C_0\delta_0-\frac12\mathrm{Pf}\,K(|u|)$, whose Fourier
transform is the archimedean multiplier $m$ of \eqref{eq:m}; the polar term $2\cosh\frac u2$; and
the prime atoms $-\frac{v_n}2\big(\delta_{\log n}+\delta_{-\log n}\big)$, $n\ge2$. For $T>0$ put
\[
\mathfrak S(T)=\Big\{\nu\ \text{a positive symmetric tempered measure on }\R:\
\langle\nuh-g,h\rangle=0\ \text{for every even }h\in C_c^\infty(-2T,2T)\Big\}.
\]
The \emph{Kre\u{\i}n extension problem} $\mathrm{KEP}(T)$ is \emph{solvable} if
$\mathfrak S(T)\ne\emptyset$, \emph{determinate} if $\#\mathfrak S(T)=1$, and \emph{indeterminate}
if $\#\mathfrak S(T)\ge2$.
\end{definition}
Symmetry is imposed because even test functions
determine only the even part of a representing measure, so without it every problem would be
indeterminate for a trivial reason. We define ``the datum is positive definite on the window'' to
mean solvability; by the Schwartz--Bochner theorem this is equivalent to saying that the restriction
of $g$ to $(-2T,2T)$ extends to a positive definite tempered distribution on $\R$.

\begin{lemma}[one negative square in the polar term]\label{lem:polar}
At any $n\ge2$ distinct points $t_1,\dots,t_n$ the Gram matrix of the kernel
$(s,t)\mapsto2\cosh\frac{s-t}2$ is $ab^{\mathsf T}+ba^{\mathsf T}$ with $a_i=e^{t_i/2}$ and
$b_i=e^{-t_i/2}$; it has rank $2$ and signature $(1,1)$. Hence $2\cosh\frac u2$ is positive definite
on no interval, and the solvability of $\mathrm{KEP}(T)$ is automatic at no window, however small.
\end{lemma}
\begin{proof}
$2\cosh\frac{t_i-t_j}2=e^{t_i/2}e^{-t_j/2}+e^{-t_i/2}e^{t_j/2}$, which is the stated rank-two matrix.
It vanishes on $\{a,b\}^\perp$ and on $\mathrm{span}\{a,b\}$ has eigenvalues $a\cdot b\pm|a||b|$;
$a\cdot b=n$ while $|a||b|>n$ strictly by Cauchy--Schwarz, the vectors $a,b$ not being parallel when
the $t_i$ are distinct and $n\ge2$. So one eigenvalue is positive and one negative.
\end{proof}
The datum \eqref{eq:gdatum} is therefore the sum of a positive definite part ($m\ge0$), a part with
exactly one negative square (the pole), and a negative part (the prime atoms). This is the reason a
polar multiplier $\beta_0$ must be carried at all (Proposition~\ref{prop:P}).

\begin{theorem}[a certificate is an absolutely continuous solution]\label{thm:KEP1}
Let $T\ge\frac12\log2$ and let $w\in\Wcl^{\mathrm{reg}}$ satisfy $f_w\ge c$ on $\R$.
\emph{(i)} If $c\ge0$ then $\nu_w:=f_w(\tau)\,d\tau/2\pi$ belongs to $\mathfrak S(T)$; in particular
$\mathrm{KEP}(T)$ is solvable, by an absolutely continuous solution.
\emph{(ii)} If $c>0$ then $\mathfrak S(T)$ contains a continuum of pairwise distinct absolutely
continuous measures; in particular $\mathrm{KEP}(T)$ is indeterminate.
\end{theorem}
\begin{proof}
(i) For $w\in\Wcl^{\mathrm{reg}}$ the product $w\cdot g$ is a tempered distribution and
Theorem~\ref{thm:Z} says $\widehat{w\cdot g}=f_w$, a continuous function with
$f_w(\tau)=\log|\tau|+O(1)$ (Lemma~\ref{lem:mbounds}, Theorem~\ref{thm:Dtail}), hence tempered;
Fourier inversion gives $\nuh_w=w\cdot g$, which agrees with $g$ on $(-2T,2T)$ because $w\equiv1$
there. Positivity of $\nu_w$ is $f_w\ge0$, and symmetry is the evenness of $w$ and $g$.
(ii) Fix an even $\chi\in C_c^\infty(\R)$, $\chi\ge0$, $\chi\not\equiv0$, with
$\supp\chi\subset\{|u|\ge\max(2T,\log2)\}$. For $\varepsilon\in\R$ put $w_\varepsilon=w+\varepsilon\chi
\in\Wcl^{\mathrm{reg}}$; then $f_{w_\varepsilon}=f_w+\varepsilon r_\chi$ with
\[
r_\chi(\tau)=2\int\chi(u)\cosh\tfrac u2\cos\tau u\,du-\sum_{n\ge2}v_n\chi(\log n)\cos(\tau\log n)
-\tfrac12\int_\R\chi(u)K(u)\cos\tau u\,du ,
\]
a bounded continuous function, say $|r_\chi|\le R$ ($\chi$ is compactly supported away from the
singularity of $K$, and $\sum_nv_n|\chi(\log n)|<\infty$). For $|\varepsilon|<c/R$ one has
$f_{w_\varepsilon}>0$, so $\nu_{w_\varepsilon}\in\mathfrak S(T)$ by (i). They are pairwise distinct:
$\nu_{w_\varepsilon}-\nu_{w_{\varepsilon'}}=\frac{\varepsilon-\varepsilon'}{2\pi}r_\chi\,d\tau$, and
$r_\chi\equiv0$ would force $\chi\cdot g=0$, whereas on $\supp\chi\subset\{|u|\ge\log2\}$ the
absolutely continuous part of $g$ is $\chi(u)d_*(u)$ with $d_*(u)=2\cosh\frac u2-\frac12K(u)>0$
(positivity and monotonicity of $d_*$ on $[0.3,\infty)$ are established in the proof of
Theorem~\ref{thm:Zexh}), which is nonnegative and not identically zero and which the singular prime
atoms cannot cancel.
\end{proof}

\begin{theorem}[solvability implies Weil positivity at that window]\label{thm:KEP2}
If $\mathfrak S(T)\ne\emptyset$ then $\lamW\ge0$.
\end{theorem}
\begin{proof}
Let $\nu\in\mathfrak S(T)$ and $F\in C_c^\infty(\IT)$, and put $h=F\star\widetilde F$ with
$\widetilde F(u)=\overline{F(-u)}$. Then $h$ is even, $\supp h\subset(-2T,2T)$ and
$\widehat h=|\Fh|^2\ge0$ is rapidly decreasing, so the pairing against the positive tempered measure
$\nu$ is legitimate. Since $h$ vanishes off the window, only the contacts $n\in\NT$ appear in
\eqref{eq:gdatum}, and comparison with \eqref{eq:Aform} and \eqref{eq:W} gives
$\langle g,h\rangle=\mathfrak a_T(F)+2\Rel\big(A(F)\overline{B(F)}\big)=W_T(F)$. By
Definition~\ref{def:KEP}, $W_T(F)=\int\widehat h\,d\nu=\int|\Fh|^2d\nu\ge0$. Finally
$C_c^\infty(\IT)$ is a form core: for $F$ in the form domain, the dilates $F_\lambda(t)=F(t/\lambda)$,
$\lambda\uparrow1$, are supported in $\lambda\IT$ and converge to $F$ in $L^2$ with
$\Fh_\lambda(\tau)=\lambda\Fh(\lambda\tau)$, so that
$\int\log(2+|\tau|)|\Fh_\lambda-\Fh|^2\to0$ (the weight is doubling), and mollifying $F_\lambda$ at a
scale below $(1-\lambda)T$ keeps the support inside $\IT$ and converges in the same norm. Hence
$W_T\ge0$ on the whole form domain and $\lamW\ge0$.
\end{proof}

\begin{theorem}[the dichotomy at one window]\label{thm:KEPdich}
For every $T\ge\frac12\log2$,
\[
\lamW>0\ \Longrightarrow\ \mathrm{KEP}(T)\ \text{is indeterminate}\ \Longrightarrow\ \lamW\ge0 ,
\]
so certificates and indeterminacy coincide except on the set $\{T:\lamW=0\}$.
\end{theorem}
\begin{proof}
If $\lamW>0$, Theorem~\ref{thm:Zexh}(i) gives $w\in\Wcl^{\mathrm{reg}}$ with $\inf f_w>0$, and
Theorem~\ref{thm:KEP1}(ii) applies; the second implication is Theorem~\ref{thm:KEP2}.
\end{proof}

Off the set $\{T:\lamW=0\}$ the solutions of the extension problem and the certificates therefore
determine each other. A Nevanlinna-type parametrization of $\mathfrak S(T)$ in the indeterminate case
would describe the set of certificates, but it presupposes one known solution.


\section{Lattice certificates and their ceilings}\label{sec:class}

This section bounds from above what a given class of dual objects can certify.

\subsection{A class-dependent bound}
\begin{problem}[class-dependent value bound]\label{prob:A}
For the lattice class of step $d=2T/M$ and range $U$, with the tail rows of Lemma~\ref{lem:F1},
find an explicit $D(T,d,U)\to0$ as $d\to0$, $U\to\infty$, such that the best certificate value
$c(T,d,U)$ satisfies $c\ge\lamP-D(T,d,U)$. Equivalently: a quantitative one-sided approximation
theorem for the function $-\sT-\beta\cdot\widehat g$ on $[0,\tau_{\mathrm{env}}]$ by $P$-periodic
trigonometric polynomials whose low modes $|j|<M$ vanish, in the norm $L^1(|\Fh_\star|^2d\tau/2\pi)$
of the primal minimizer.
\end{problem}
A solution would decide each window by a computation of predictable size and would turn
Theorem~\ref{thm:C} into a construction. By the identity $\lamP-c=\int(f-c)\,d\mu_\star$
(Proposition~\ref{prop:weak}) the deficit is the $\mu_\star$-average of the slack, so the
approximation is needed only where $|\Fh_\star|^2$ lives. Theorem~\ref{thm:Vp} shows that $D$ must
depend on the class, not on a radius alone. If Zhu's conjectural law $-\log\lamP\asymp e^{2T}$ holds,
the effort in Problem~\ref{prob:A} grows at least like $e^{2T}$: finite for each $T$, never uniform.

The converse of Problem~\ref{prob:A}, an upper bound on the value of a class, is a theorem.

\begin{proposition}[test-kernel ceiling for lattice classes]\label{prop:TK}
Let $d=2T/M$, $P=2\pi/d$, and let $\nu$ be any finite even measure supported on the lattice
$\{\pm jd:j\ge M\}$ (any range, any number of atoms). Put $S_\beta=\sT+\beta_0\gh0+\beta_1\gh1$ and,
for $\tau_0\in\R$, let $C_M(\tau_0,\beta)$ be the $M\times M$ Hermitian Toeplitz matrix with entries
$c_{j-k}$, $c_\ell=P^{-1}\int_{\tau_0}^{\tau_0+P}S_\beta(\tau)e^{i\ell d\tau}\,d\tau$. If
$\sT+\nuh+\beta\cdot\widehat g\ge c$ on $\R$ then
\[
c\ \le\ \lambda_{\min}\big(C_M(\tau_0,\beta)\big)\ \text{ for every }\tau_0,\qquad\text{hence}\qquad
c\ \le\ \operatorname{Ceil}(M,T):=\sup_{\beta\in\R^2}\inf_{\tau_0}\lambda_{\min}\big(C_M(\tau_0,\beta)\big),
\]
and $\beta\mapsto\inf_{\tau_0}\lambda_{\min}(C_M(\tau_0,\beta))$ is concave. The ceiling depends on
the class only through $M$.
\end{proposition}
\begin{proof}
$\nuh(\tau)=\sum_{j\ge M}2a_j\cos(jd\tau)$ is $P$-periodic with Fourier modes $|j|\ge M$ only; in
particular it has mean zero over any period. Let $q(z)=\sum_{k=0}^{M-1}q_kz^k$ and
$\vartheta(\tau)=|q(e^{id\tau})|^2\ge0$, a trigonometric polynomial in $d\tau$ of degree $\le M-1$. Then
$\int_{\tau_0}^{\tau_0+P}\nuh\vartheta\,d\tau=0$, so integrating $f\ge c$ against $\vartheta$ gives
$c\int\vartheta\le\int S_\beta\vartheta$. Now $\int_{\tau_0}^{\tau_0+P}\vartheta=P\sum|q_k|^2$ and
$\int_{\tau_0}^{\tau_0+P}S_\beta\vartheta=P\sum_{j,k}\bar q_kq_jc_{j-k}$, so
$c\le\langle C_Mq,q\rangle/\|q\|^2$ for every $q\ne0$, i.e.\ $c\le\lambda_{\min}(C_M)$; by the
Fejér--Riesz theorem every nonnegative trigonometric polynomial of degree $\le M-1$ is of the form
$|q|^2$, so nothing is lost. For fixed $q$ and $\tau_0$ the Rayleigh quotient is affine in $\beta$
($S_\beta$ is), so $\inf_{q,\tau_0}$ is concave in $\beta$.
\end{proof}

The ceiling is the smallest eigenvalue of a small matrix and needs no optimizer; like
Theorem~\ref{thm:Vp}, it depends on the class and not on a radius. Numerically, at
$T=\frac12\log3$ it is negative for $M=8$ and positive for $M=12,24$, and the certificate of
Theorem~\ref{thm:C} realizes a fraction of about $15\%$ of its own class ceiling
(double-precision values, Appendix~\ref{app:num}). The negative sign at $M\le8$ survives an interval enclosure; this is
Theorem~\ref{thm:N} below. Since the lattices $2T/2^kM_0$ are nested with dense union,
Theorem~\ref{thm:D1} gives $\sup_M\operatorname{Ceil}(M,T)\ge\lamP$; that the ceilings in fact
converge to $\lamP$, so that the test-kernel family is not lossy in the limit, is
Theorem~\ref{thm:R} and Corollary~\ref{cor:R3} below. We do not know whether the family is lossy at
fixed $M$, nor at what rate the ceilings converge; this is the tractable half of
Problem~\ref{prob:A}.


\subsection{Lattices of small order at $T=\tfrac12\log3$}\label{subsec:N}
The ceiling of Proposition~\ref{prop:TK} is a supremum over $\beta\in\R^2$, and turning a numerical
ceiling into a theorem requires control of $\beta$. The following a priori bound provides it: the
two primitivity multipliers of a certificate cannot be large, with a constant determined by the
geometry of the moment image of Lemma~\ref{lem:W}.

\begin{theorem}[a priori bound on the multipliers]\label{thm:Bbeta}
Fix $T$ and nonnegative weights $v$. Suppose there are $r_0>0$, $V^\star<\infty$ and finitely many
unit vectors $F_1,\dots,F_N\in H^1_0(\IT)$ such that, writing $w_i=(Q_0(F_i),Q_1(F_i))$ in the
notation of Lemma~\ref{lem:W} and $V_i=\mathfrak a_T(F_i)$,
\[
\{w\in\R^2:|w|\le r_0\}\subseteq\operatorname{conv}\{w_1,\dots,w_N\},\qquad V_i\le V^\star\ \ (1\le i\le N).
\]
Then every certificate $(\nu,\beta)$ of value $c$ satisfies $\|\beta\|_2\le(V^\star-c)/r_0$.
\end{theorem}
\begin{proof}
For a unit $F\in H^1_0(\IT)$ the measure $\mu_F=|\Fh|^2d\tau/2\pi$ is an even probability measure with
$\check\mu_F=\cF$ supported in $[-2T,2T]$ and $\cF(\pm2T)=0$, so $\int\nuh\,d\mu_F=\int\cF\,d\nu=0$ for
every measure $\nu$ supported in $\{|u|\ge2T\}$, exactly as in the proof of
Proposition~\ref{prop:weak}; moreover $\int\gh j\,d\mu_F=\int g_j\cF=Q_j(F)$ and
$\int\sT d\mu_F=\mathfrak a_T(F)$. Integrating $f\ge c$ against $\mu_F$ gives
$\mathfrak a_T(F)+\beta\cdot w(F)\ge c$. The relaxed feasible set (even probability measures whose
cosine transform vanishes off the window, with the two moment constraints dropped) is
convex, and both $\mu\mapsto\int\sT d\mu$ and $\mu\mapsto\mathfrak m(\mu)$ are affine; hence for
$w=\sum\lambda_iw_i$ in the convex hull the mixture $\mu=\sum\lambda_i\mu_{F_i}$ is admissible, has
moment $w$ and value at most $V^\star$, so $\beta\cdot w\ge c-V^\star$ for every $|w|\le r_0$.
Applying this to $w$ and $-w$ gives $|\beta\cdot w|\le V^\star-c$ on the disc of radius $r_0$, and
taking the supremum there, $r_0\|\beta\|_2\le V^\star-c$.
\end{proof}

\begin{lemma}[value of a witness]\label{lem:val}
For a unit $F\in H^1_0(\IT)$ with $L:=\|F'\|_2$ and every $u_0>0$,
\[
\mathfrak a_T(F)\ \le\ C_0+\Aconst+\frac{L^2}2\int_0^{u_0}u^2K(u)\,du+2\int_{u_0}^\infty K(u)\,du,
\]
the right side being minimized at $u_0=2/L$; the second integral is
$\log\coth\frac{u_0}2+\pi-2\arctan e^{u_0}$ by Theorem~\ref{thm:S}(iii).
\end{lemma}
\begin{proof}
By Theorem~\ref{thm:A}, $\mathfrak a_T(F)=C_0+\Lev(\cF)-\sum v_n\cF(\log n)\le C_0+\Aconst+
\int_0^\infty(1-\Rel\cF(u))K(u)du$, since $|\cF|\le\|F\|^2=1$. Now
$1-\Rel\cF(u)=\frac1{2\pi}\int(1-\cos\tau u)|\Fh|^2d\tau\le\min\{2,\frac{u^2}2\|F'\|^2\}$, the first
bound from $|\cF|\le1$ and the second from $1-\cos x\le x^2/2$ and Plancherel; split the integral at
$u_0$. The derivative in $u_0$ is $K(u_0)(\frac{L^2}2u_0^2-2)$, which vanishes and changes sign at
$u_0=2/L$.
\end{proof}

At $T=\frac12\log3$ we take as witnesses the unit vectors of the four-dimensional span of
$\psi_k(t)=(T^2-t^2)t^k$, $k=0,1,2,3$, inside $H^1_0(\IT)$: all the quantities entering
$Q_0=AB$, $Q_1=A'B-AB'$ and Lemma~\ref{lem:val} are elementary integrals of a polynomial against
$e^{\pm t/2}$. Taking, for each of $48$ directions, the top generalized eigenvector of the
corresponding combination of the two moment forms relative to the Gram matrix, computing the convex
hull of the $48$ moment vectors and the distance from the origin to its boundary gives
\[
r_0=8.4515\cdot10^{-5},\qquad V^\star=3.817495,\qquad
\|\beta\|_\infty\le\frac{V^\star}{r_0}=45169.4
\]
\[
\text{for every certificate of value }c\ge0 .
\]
The size of this constant is structural rather than a defect of the witnesses: on the window
$g_0(u)-\frac14g_1(u)=\cosh\frac u2-\frac u4\sinh\frac u2=1-\frac{u^4}{384}+O(u^6)$, so
$Q_0(F)-\frac14Q_1(F)=|\int F|^2+O(T^4)$ and the moment image is a thin sliver lying essentially on
one side of the line $Q_0=\frac14Q_1$; hence $r_0=O(T^4)$ for any witness family, and the bound is
$O(T^{-4})$. It is far too large to restrict a search, and good enough for the qualitative
statements below, because the coefficients it multiplies are of size $10^{-9}$.

\begin{theorem}[class impossibility at $T=\frac12\log3$]\label{thm:N}
Let $T=\frac12\log3$, $\NT=\{2\}$, $M\in\{1,\dots,8\}$,
$d=2T/M$, and let $\nu$ be any finite even signed measure supported on the lattice
$\{\pm jd:j\ge M\}$ (any range, any number of atoms), and $\beta\in\R^2$ arbitrary. Then
\[
\inf_{\tau\in\R}\big[\sT(\tau)+\nuh(\tau)+\beta_0\gh0(\tau)+\beta_1\gh1(\tau)\big]\ <\ 0 .
\]
Consequently no dual certificate in the sense of Definition~\ref{def:PD} supported on a lattice of
step $2T/M$ with $M\le8$ has nonnegative value, and none proves $\lamP>0$ at this window, for any
$\beta$, any range, any atom count and whichever of the tail lemmas of \S\ref{sec:tails} is used to
verify the inequality. The proof includes a computer verification in interval arithmetic
(Appendix~\ref{app:encl}).
\end{theorem}
\begin{proof}
Suppose $f:=\sT+\nuh+\beta\cdot\widehat g\ge c$ on $\R$ with $c\ge0$.

\emph{Step 1: $\beta$ is bounded.} By Theorem~\ref{thm:Bbeta} and the constants above,
$\|\beta\|_\infty\le45169.4$. All that is used below is $\|\beta\|_\infty\le B$ for some
$B\le9.7\cdot10^5$; the computed bound is about $4.5\cdot10^4$, so the hypothesis holds with a
factor $21$ to spare, and no sharp value of $V^\star/r_0$ is needed.

\emph{Step 2: one test kernel.} Fix $M$ and a pair $(\tau_0,q)$ with $\tau_0\in\Q$ and
$q\in\C^M$, and put $\vartheta(\tau)=|\sum_{k<M}q_ke^{ikd\tau}|^2\ge0$, $\mathcal I_0=[\tau_0,\tau_0+P]$,
$P=2\pi/d=\pi M/T$. As in the proof of Proposition~\ref{prop:TK}, $\nuh$ is $P$-periodic with Fourier
modes in $\{|j|\ge M\}$ while $\vartheta$ has modes in $\{|j|\le M-1\}$, so $\int_{\mathcal I_0}\nuh\vartheta=0$;
integrating $f\ge c$ against $\vartheta\ge0$ over $\mathcal I_0$ and dividing by $\int_{\mathcal I_0}\vartheta=P\|q\|^2$ (exact,
by orthogonality over a full period) gives
\[
c\ \le\ a+b_0\beta_0+b_1\beta_1,\qquad a=\frac{\int_{\mathcal I_0}\sT\vartheta}{P\|q\|^2},\quad
b_j=\frac{\int_{\mathcal I_0}\gh j\vartheta}{P\|q\|^2},
\]
valid for every $\beta$. This is Proposition~\ref{prop:TK} at a single $(\tau_0,\vartheta)$: no
infimum over $\tau_0$, no eigenvalue and no concave maximization are needed, which is what makes the
enclosure cheap. Any pair $(\tau_0,q)$ gives a valid bound, so the way the pair was found (in
practice, as a double-precision eigenvector) plays no role in the proof.

\emph{Step 3: the enclosure at $M=8$.} With $\tau_0=-27.223$ and the vector $q$ contained in the
certificate data accompanying this paper, interval arithmetic (Appendix~\ref{app:encl}) gives
\[
a\in[-6.009422\cdot10^{-3},\,-6.006218\cdot10^{-3}],\qquad |b_0|+|b_1|\le6.18\cdot10^{-9},
\]
so $c\le-6.006218\cdot10^{-3}+6.18\cdot10^{-9}B\le-5.70\cdot10^{-3}<0$ for every $B\le9.7\cdot10^5$,
contradicting $c\ge0$.

\emph{Step 4: the remaining $M$.} The same four enclosures with their own pairs $(\tau_0,q)$ give
$c\le-9.038987\cdot10^{-2}$ $(M=3)$, $-6.234677\cdot10^{-2}$ $(M=5)$, $-2.175677\cdot10^{-2}$ $(M=6)$
and $-4.242754\cdot10^{-2}$ $(M=7)$, all negative, the allowance $(|b_0|+|b_1|)B$ already included.
For $M\in\{1,2,4\}$ no computation is needed: $M\mid8$, so every multiple of $2T/M$ that is $\ge2T$ is
a multiple of $2T/8$ that is $\ge2T$, i.e.\ the class at step $2T/M$ is contained in the class at step
$2T/8$, which Step~3 already excludes.
\end{proof}

\begin{remark}
Theorem~\ref{thm:N} is not a statement about $\lamP(\frac12\log3)$, which is positive by
Theorem~\ref{thm:C}, a certificate at $M=24$; it says that the class $M\le8$ cannot detect this
positivity.
\end{remark}

\subsection{The same impossibility at the two-prime windows}\label{subsec:Np}

The enclosure of Theorem~\ref{thm:N} is limited by the size of $|a|$, which at $\frac12\log3$ is of
order $10^{-3}$. At the two-prime windows the ceilings are four orders of magnitude smaller, and
reaching them needs a quadrature rule of high order and a test kernel designed against the
$\beta$-box rather than merely against $\sT$. Both are elementary, and they give the same kind of
statement at windows where no certificate is known.

\begin{theorem}[class impossibility at the two-prime windows]\label{thm:Np}
Let $T\in\{0.68,\log2\}$, so $\NT=\{2,3\}$ and
$\Aconst=2.24883$. Let $M\in\{1,\dots,24\}$, $d=2T/M$, let $\nu$ be any finite even signed measure
supported on the lattice $\{\pm jd:j\ge M\}$ (any range, any number of atoms, any signs), and
$\beta\in\R^2$ arbitrary. Then
\[
\inf_{\tau\in\R}\big[\sT(\tau)+\nuh(\tau)+\beta_0\gh0(\tau)+\beta_1\gh1(\tau)\big]\ <\ 0 .
\]
Consequently no dual certificate in the sense of Definition~\ref{def:PD} supported on a lattice of
step $2T/M$ with $M\le24$ has nonnegative value at either window, and none proves $\lamP>0$ there,
for any $\beta$, any range, any atom count and whichever tail lemma of \S\ref{sec:tails} is used to
verify the inequality. In particular
\[
\operatorname{Ceil}(24,0.68)\le-7.493129\cdot10^{-6},\qquad
\operatorname{Ceil}(24,\log2)\le-6.584323\cdot10^{-7}.
\]
The proof includes a computer verification in interval arithmetic (Appendix~\ref{app:encl}).
\end{theorem}
\begin{proof}
The three steps are those of Theorem~\ref{thm:N}. \emph{Step 1.} Theorem~\ref{thm:Bbeta}, with its
witnesses and moment polygon recomputed at these two windows, gives $r_0=2.4309\cdot10^{-4}$,
$V^\star=4.595697$ at $T=0.68$ and $r_0=2.6720\cdot10^{-4}$, $V^\star=4.550671$ at $T=\log2$, hence
$\|\beta\|_\infty\le B$ for every certificate of value $c\ge0$ with $B=1.89\cdot10^4$, resp.
$1.70\cdot10^4$. \emph{Step 2.} For a pair $(\tau_0,q)$ put $\vartheta(\tau)=|\sum_{k<M}q_ke^{ikd\tau}|^2$,
$\mathcal I_0=[\tau_0,\tau_0+P]$, $P=\pi M/T$; as in Proposition~\ref{prop:TK}, $\int_{\mathcal I_0}\nuh\vartheta=0$ and
$c\le a+b_0\beta_0+b_1\beta_1$ with $a=(P\|q\|^2)^{-1}\int_{\mathcal I_0}\sT\vartheta$ and
$b_j=(P\|q\|^2)^{-1}\int_{\mathcal I_0}\gh j\vartheta$, so $c\le a+(|b_0|+|b_1|)B$. \emph{Step 3.} Sixteen enclosures
per window, at $M\in\{6,7,9,11,13,14,\dots,24\}$, are all negative; the two displayed values are the
upper endpoints at $M=24$, the whole $\beta$-allowance included. For the remaining
$M\in\{1,2,3,4,5,8,10,12\}$ no computation is needed: the class at step $2T/M$ is contained in the
class at step $2T/(kM)$ whenever $k\in\N$, and every such $M$ divides $20$ or $24$, both of which are
enclosed.
\end{proof}

\begin{remark}
Theorem~\ref{thm:Np} is an upper bound on what a class can certify; the sign of $\lamP(0.68)$ and of
$\lamP(\log2)$ is not decided by the methods of this paper. It shows that the class $M=24$, which
carries the certificate of Theorem~\ref{thm:C} at $\frac12\log3$, contains no certificate of
nonnegative value at either of these windows, and neither does any lattice class of order $M\le24$.
\end{remark}

\subsection{The ceilings converge}\label{subsec:R}
Proposition~\ref{prop:TK} bounds a lattice class from above by one eigenvalue. The question left
open after it is whether this bound is asymptotically sharp, i.e.\ whether the test kernels see all
of the primal problem. They do, in the limit; the mechanism is Poisson summation used in the
direction opposite to the one that makes a lattice $\nuh$ periodic.

\begin{lemma}[the periodization of a primal measure is a test kernel]\label{lem:R1}
Let $\mu$ be feasible for \textup{(P)}, $M\in\N$, $d=2T/M$, $P=2\pi/d=\pi M/T$, and let $\rho$ be the
periodization of $\mu$ on $\R/P\Z$, $\rho(E)=\sum_{k\in\Z}\mu(E+kP)$. Then $\rho$ is a probability
measure with Fourier coefficients $\widehat\rho(j)=\check\mu(jd)$, which vanish for $|j|\ge M$. Hence
$\rho$ is absolutely continuous with a nonnegative trigonometric polynomial density of degree
$\le M-1$ in $\theta=d\tau$, and by the Fej\'er--Riesz theorem
$d\rho=P^{-1}|q(e^{id\tau})|^2d\tau$ for some polynomial $q$ of degree $\le M-1$: $\rho$ is a test
kernel of order $M$ in the sense of Proposition~\ref{prop:TK}.
\end{lemma}
\begin{proof}
$\rho$ is a finite positive measure of mass $\mu(\R)=1$, and
$\widehat\rho(j)=\int_{[\tau_0,\tau_0+P)}e^{-ijd\tau}d\rho=\int_\R e^{-ijd\tau}d\mu=\check\mu(jd)$,
because $e^{-ijd(\tau+kP)}=e^{-ijd\tau}$ ($dP=2\pi$) and $\mu$ is even. Since $|jd|\ge2T$ exactly when
$|j|\ge M$, and feasibility is the closed condition $\check\mu=0$ on $\{|u|\ge2T\}$ (so the boundary
index $j=M$ is included), only finitely many coefficients are nonzero; a measure with finitely many
nonzero Fourier coefficients is a trigonometric polynomial, here nonnegative, and Fej\'er--Riesz
factorizes it.
\end{proof}

\begin{theorem}[the ceilings converge to the primal value]\label{thm:R}
For $\beta\in\R^2$ write $\operatorname{Ceil}(M,T;\beta)=\inf_{\tau_0}\lambda_{\min}(C_M(\tau_0,\beta))$
and $S_\beta=\sT+\beta\cdot\widehat g$. Then for every $\mu$ feasible for \textup{(P)},
\[
\operatorname{Ceil}(M,T;\beta)\ \le\ \int\sT\,d\mu+\Delta(M,\beta)\,\mu\big(\{|\tau|>\pi M/2T\}\big),
\]
\[
\Delta(M,\beta)=\big[m(P/2)-C_0\big]+2\Aconst+2|\beta_0|\gh0(0)+2|\beta_1|\gh1(0),
\]
with $P/2=\pi M/2T$. Consequently $\limsup_{M\to\infty}\operatorname{Ceil}(M,T;\beta)\le\lamP$ for
every fixed $\beta$, and uniformly for $\beta$ in a bounded box.
\end{theorem}
\begin{proof}
Take $\tau_0=-P/2$, $\mathcal I_0=[-P/2,P/2]$, and let $\pi_0(\tau)=\tau-kP\in\mathcal I_0$ be the folding map. By
Lemma~\ref{lem:R1} and the Fej\'er--Riesz description of the matrix $C_M$ in the proof of
Proposition~\ref{prop:TK}, the infimum over test kernels is at most its value at $P\rho$, so
$\operatorname{Ceil}(M,T;\beta)\le\int_{\mathcal I_0}S_\beta\,d\rho=\int_\R S_\beta(\pi_0\tau)\,d\mu(\tau)$. Now
$\int S_\beta d\mu=\int\sT d\mu$ because $\mu$ has both moments zero, and $\pi_0\tau=\tau$ on
$|\tau|\le P/2$, so
\[
\int S_\beta\circ\pi_0\,d\mu-\int\sT\,d\mu=\int_{|\tau|>P/2}\big(S_\beta\circ\pi_0-S_\beta\big)d\mu
\le\big(\sup_{\mathcal I_0}S_\beta-\inf_\R S_\beta\big)\,\mu(|\tau|>P/2),
\]
and $\sup_{\mathcal I_0}S_\beta\le m(P/2)+\Aconst+|\beta_0|\gh0(0)+|\beta_1|\gh1(0)$ ($m$ increasing,
$|\gh j|\le\gh j(0)$) while $\inf_\R S_\beta\ge C_0-\Aconst-|\beta_0|\gh0(0)-|\beta_1|\gh1(0)$, which
is $\Delta$. For the limit, fix $\delta>0$ and $\mu$ with $\int\sT d\mu\le\lamP+\delta$. Then
$\int m\,d\mu<\infty$ (since $\sT\ge m-\Aconst$), so by dominated convergence
$(m(R)-C_0)\mu(|\tau|>R)\le\int_{|\tau|>R}(m-C_0)d\mu\to0$ as $R\to\infty$; as
$\Delta(M,\beta)=(m(P/2)-C_0)+O_\beta(1)$ and $P/2\to\infty$, the product tends to $0$, uniformly for
$\beta$ in a bounded box.
\end{proof}

\begin{corollary}[the test-kernel family is not lossy in the limit]\label{cor:R3}
Suppose $\lamP>0$, and let $B=V^\star/r_0$ be the a priori bound of Theorem~\ref{thm:Bbeta}, so that
every certificate of value $c\ge0$ has $\|\beta\|_\infty\le B$. Put
$\operatorname{Ceil}_B(M,T)=\sup_{\|\beta\|_\infty\le B}\operatorname{Ceil}(M,T;\beta)$. Then
$\operatorname{Ceil}_B(M,T)\to\lamP$ along the nested lattices $2T/2^kM_0$ with growing range. In
particular the subfamily of test kernels $\{|q|^2\}$ used in Proposition~\ref{prop:TK} is not lossy in
the limit: its optimal value is the same $\lamP$ as that of the full relaxed-feasible family.
\end{corollary}
\begin{proof}
$\limsup_M\operatorname{Ceil}_B\le\lamP$ is Theorem~\ref{thm:R}, $\Delta$ being bounded uniformly on
the box. Conversely, by Theorem~\ref{thm:D1} there are, for every $\varepsilon>0$, lattice
certificates of value $c>\lamP-\varepsilon$ on the nested lattices (whose union is dense in
$[2T,\infty)$); eventually $c\ge0$, so by Theorem~\ref{thm:Bbeta} their $\beta$ lies in the box, and
Proposition~\ref{prop:TK} gives $c\le\operatorname{Ceil}_B(M,T)$.
\end{proof}
We do not know whether the family is lossy at fixed $M$, nor at what rate the ceilings converge.

\subsection{Sampling a flat-top certificate}\label{subsec:lat}
Problem~\ref{prob:A} asks for the loss incurred by restricting the dual to a lattice class. For the
flat-top certificates of \S\ref{sec:Z} the loss of the sampled certificate can be computed
exactly. It is an edge effect of order $1/M$, not a smoothness effect, and it is carried entirely by
the polar block: the density of $\nu_w$ in \eqref{eq:nuw} is $\frac12(1-w)K+2w\cosh\frac u2$, whose
second summand does not vanish at $u=2T^+$ (it equals $2\cosh T$ there, the jump already met in
the proof of Theorem~\ref{thm:Dtail}), so no smoothness of $w$ can make the sampling error
$O(M^{-p})$.

\begin{theorem}[the edge cost of lattice sampling]\label{thm:latedge}
Let $p\ge2$ and let $w\in\Wcl$ be even with $w^{(p-1)}$ locally absolutely continuous and
$w^{(p)}\in \mathrm{BV}\cap L^1$, $\supp w\subseteq[-U,U]$ with $C^{p-1}$ contact with $0$ at $\pm U$,
and $w(\log n)=0$ for every $n$ with $\log n>2T$ \textup{(}comb-vanishing, finitely many linear
conditions\textup{)}. No sign condition and no bound $0\le w\le1$ is assumed. The boundary case
$\log n=2T$ is not excluded; it is treated by carrying the atom it produces \textup{(}see
$\nu_{\mathrm{lat}}$ below\textup{)}, which is legitimate because Definition~\ref{def:W} forces
$w(2T)=1$ there, so that atom is a fixed multiple of $\delta_{\pm2T}$ and can be reproduced exactly.
Put
$\Omega(u)=w(|u|)\cosh\frac u2$, $R=\frac12(1-w)K\mathbf 1_{|u|\ge2T}$ and
$G=R+2\Omega\mathbf 1_{|u|\ge2T}$, which by comb-vanishing is the whole absolutely continuous density
of $\nu_w$, the only surviving atom being the boundary one at $\pm2T$ when $2T\in\{\log n\}$; put
$V_0=\|G\|_{L^1}$, $V_p=\|R^{(p)}\|_{L^1}+2\|\Omega^{(p)}\|_{L^1}$ and $c_w=\inf_\R f_w$. For $M\in\N$ set
$d=2T/M$, $P=2\pi/d=\pi M/T$ and
\[
\nu_{\mathrm{lat}}=\sum_{|j|\ge M}d\,\kappa_j\,G(jd)\,\delta_{jd}
-\sum_{n:\,\log n=2T}\tfrac{v_n}2\big(\delta_{2T}+\delta_{-2T}\big),\qquad
\kappa_{\pm M}=\tfrac12,\quad\kappa_j=1\ (|j|>M),
\]
an even signed measure with countably many atoms and finite total variation
$\|\nu_{\mathrm{lat}}\|_{\TV}=V_0+O(d)$. Then $(\nu_{\mathrm{lat}},\beta=(2,0))$ is dual feasible and
\[
c_{\mathrm{lat}}(M)\ \ge\ \min\Big\{c_w-E_{\mathrm{edge}}(M)-\varepsilon_p(M),\
m\big(\tfrac{\pi M}{2T}\big)-\Aconst-\|\nu_{\mathrm{lat}}\|_{\TV}-2\gh0(0)\Big\},
\]
\[
E_{\mathrm{edge}}(M)=\frac{8\cosh T}{P}+\frac{2\pi^2\sinh T}{P^2}+\frac{16.9\cosh T}{P^3}
=\frac4\pi\cosh T\cdot d+O(d^2),\qquad
\varepsilon_p(M)=2\big(1-2^{-p}\big)\zeta(p)V_p\Big(\frac{\pi M}{2T}\Big)^{-p}.
\]
Moreover $\sup_\tau|f_{\mathrm{lat}}-f_w|=\frac4\pi\cosh T\cdot d\,(1+O(d))$: the rate $M^{-1}$ and
its constant are attained, and no smoothness of $w$ improves them.
\end{theorem}
\begin{proof}
Feasibility: the atoms lie at $|u|=|j|d\ge Md=2T$, $(D)$ imposes no moment equality on $\nu$
(Proposition~\ref{prop:weak}), Definition~\ref{def:PD} allows $\nu$ to be signed, the boundary atom
at $\pm2T$ is harmless because $\cF(\pm2T)=0$, and $\beta$ is free. The atoms with $|j|<M$ are absent
because $G$ vanishes on the window. Finiteness of the total variation is
$\|\nu_{\mathrm{lat}}\|_{\TV}=d\sum_{|j|\ge M}\kappa_j|G(jd)|<\infty$, since $G\in L^1\cap
\mathrm{BV}$ on $[2T,\infty)$; the number of atoms is countably infinite, as $G=\frac12K>0$ beyond
$U$.

Value: since $w(\log n)=0$ for $\log n>2T$, $\nu_w=G(u)\,du-\sum_{\log n=2T}\frac{v_n}2(\delta_{2T}+\delta_{-2T})$,
and the second sum, present exactly when $2T\in\{\log n:n\ge2\}$ (as at the window $T=\frac12\log3$,
where it is $-\frac{v_3}2(\delta_{2T}+\delta_{-2T})$), has been put verbatim into
$\nu_{\mathrm{lat}}$, so it cancels in $f_{\mathrm{lat}}-f_w$ and may be ignored below.
Theorem~\ref{thm:Dtail}
gives $\nuh_w+2\gh0=2\widehat\Omega+\Psi_w$ with $\Psi_w=\widehat R$. Because $2T=Md$ is a lattice point,
the samples of $G$ on $|j|\ge M$ split into those of $R$ and those of $2\Omega$, so
\[
\nuh_{\mathrm{lat}}(\tau)=\Big[d\sum_{|j|\ge M}\kappa_jR(jd)e^{ijd\tau}\Big]
+\Big[d\sum_{j\in\Z}2\Omega(jd)e^{ijd\tau}-T_d(\tau)\Big],
\]
where $T_d$ is the end-point-corrected trapezoidal sum, on the nodes of the window, of
$h=2\cosh\frac u2\mathbf 1_{[-2T,2T]}$, whose transform is $\widehat h=2\gh0$. Both $R$ and $2\Omega$ are
$C^{p-1}$ with $p$-th derivative in $L^1$, so Poisson summation applies absolutely to each and the
two bracketed defects are bounded on $|\tau|\le P/2$ by
$V_p\sum_{k\ne0}|\tau+kP|^{-p}\le\varepsilon_p(M)$. Hence
$f_{\mathrm{lat}}-f_w=[2\gh0-T_d]+O(\varepsilon_p)$: the whole non-smooth cost is the quadrature
error of the polar block. Poisson summation for the compactly supported function $h$ of bounded
variation, with its symmetric values at the two jumps, gives $T_d=\sum_k\widehat h(\tau+kP)$, and
since $2TP=2\pi M$ every alias carries the same phase:
\[
2\gh0-T_d=-4\Big[\tfrac12\sinh T\cos2T\tau\sum_{k\ne0}\frac1{\xi_k^2+\frac14}
+\cosh T\sin2T\tau\sum_{k\ne0}\frac{\xi_k}{\xi_k^2+\frac14}\Big],\qquad\xi_k=\tau+kP .
\]
The first sum is at most $\pi^2/P^2$. For the second write
$\xi/(\xi^2+\frac14)=\xi^{-1}-\frac14\xi^{-1}(\xi^2+\frac14)^{-1}$ and use
$\sum_{k\ne0}\xi_k^{-1}=P^{-1}\big[\pi\cot\pi z-z^{-1}\big]$ with $z=\tau/P$, a function
regular at $z=0$ (the subtracted $z^{-1}$ is exactly the omitted $k=0$ term), odd and of
modulus at most $2/P$ on $|z|\le\frac12$; the correction is at most $4.21/P^3$. This is
$E_{\mathrm{edge}}$, uniformly in $\tau$, and it is attained where $\sin2T\tau=\pm1$. For
$|\tau|>P/2$ use $\sT\ge m(P/2)-\Aconst$ (Lemma~\ref{lem:mbounds}),
$|\nuh_{\mathrm{lat}}|\le\|\nu_{\mathrm{lat}}\|_{\TV}$ and $|2\gh0|\le2\gh0(0)$.
\end{proof}

\begin{corollary}\label{cor:ALdef}
At fixed $w$ and for $M\ge M_0(w)$ (large enough that the second branch of the minimum in
Theorem~\ref{thm:latedge} does not bind, which happens for all large $M$ since $m(\pi M/2T)\to\infty$),
the lattice sampling costs
$c_w-c_{\mathrm{lat}}(M)\le\frac8\pi T\cosh T\cdot M^{-1}+O(M^{-2})+O_p(V_pM^{-p})$, and the
$M^{-1}$ is not removable: it is the alias of the jump of $2\cosh\frac u2\mathbf 1_{|u|<2T}$ at the
window edge, an object independent of $w$. Without the end-point correction
\textup{(}$\kappa_{\pm M}=1$\textup{)} the sharp constant is
$\cosh T\sqrt{4+16/\pi^2}\,d=2.3708\cosh T\cdot d$, so the boundary atom of weight $d\cosh T$ at
$\pm2T$ improves the constant by the factor $1.862$ and nothing improves the order.
\end{corollary}

Two remarks on Corollary~\ref{cor:ALdef}. First, for sampled certificates the whole lattice cost
sits at $u=\pm2T$, the same place as the non-attainment of Theorem~\ref{thm:O}; the smooth blocks $R$
and $\Omega$ contribute only $O_p(M^{-p})$. We do not know whether the lattice class, with weights
optimized rather than sampled, also loses $\Theta(M^{-1})$; this is the remaining half of
Problem~\ref{prob:A}. Second, up to $O_p(M^{-p})$ one has $f_{\mathrm{lat}}-f_w=2\gh0-T_d$, the
discrepancy between the polar multiplier and its Riemann sum over the window; $T_d$ is a
$P$-periodic trigonometric sum of degree $M$ and $2\gh0$ is not a trigonometric polynomial, so the
difference is not one either. Heuristically, this is why Proposition~\ref{prop:TK}, whose ceiling
rests on orthogonality to every mode $|j|<M$, need not bind flat-top certificates; neither
statement implies this.

\section{The archimedean form}\label{sec:archstr}

With the primes deleted, the Weil form is an explicit convolution operator. This section studies
its thresholds, its negative index, and the number of even modes that the prime form can compete
with.

\subsection{Thresholds and nodal structure}\label{subsec:arch}

Put, for $F\in L^2(\IT)$,
\[
\Ainf(F)=2A(F)B(F)+\frac1{2\pi}\int_\R m(\tau)|\Fh(\tau)|^2\,d\tau
=2A(F)B(F)+C_0\|F\|^2+\Lev(\cF),
\]
the complete Weil form with the prime comb deleted (the second equality is Theorem~\ref{thm:A} with
$\NT=\emptyset$), and write $\PiT(F)=-\sum_{n\in\NT}v_n\Rel\cF(\log n)$ for the prime form (the real part is redundant for real $F$), so that
$W_T=\Ainf+\PiT$. Let $\laminf(T)$ be the bottom of $\Ainf$ on the unit sphere of $L^2(\IT)$ and
$\laminfP(T)$ its bottom on the unit sphere of $\PT$, where the pole term vanishes identically and
$\Ainf=C_0\|F\|^2+\Lev(\cF)$.

\begin{proposition}[two archimedean thresholds]\label{prop:arch}
\emph{(i)} $\Ainf$ does not depend on $T$ and the feasible sets increase with $T$, so both
$\laminf$ and $\laminfP$ are nonincreasing (for $\laminfP$ strictly, Lemma~\ref{lem:strict}).
\emph{(ii)} $\laminf(T)\le\laminfP(T)\le C_0+\varphi\big((3\pi/2T)^2\big)$; the first inequality is
inclusion of feasible sets and the second is the proof of Corollary~\ref{cor:mono} with the comb deleted. Hence, with
$T_0(0)=0.74921$ the root of $\varphi((3\pi/2T)^2)=-C_0$, both are nonpositive at $T=T_0(0)$ and negative for $T>T_0(0)$ (at $T_0(0)$
itself, as for every $T\ge0.60$, $\laminf\le\laminfP<0$ by Corollary~\ref{cor:Tarch}), and $\laminfP(T)\downarrow C_0$ as
$T\to\infty$, since $\Lev\ge0$ and $\varphi((3\pi/2T)^2)\downarrow0$.
\emph{(iii)} For $2T<\log2$ the contact set is empty, so $W_T=\Ainf$ and $\laminf(T)\ge0$ by the
known positivity of the complete form at autocorrelation support $\le\log2$
\textup{\cite{Yoshida1992,Bombieri2000}} (and, certified, \cite[Thm.~1.2]{Zhu2026}).
\emph{(iv)} Consequently the thresholds
$T^*_\infty=\sup\{T:\laminf(T)>0\}$ and $T^*_{\infty,\mathrm P}=\sup\{T:\laminfP(T)>0\}$ exist, and
\[
T^*_\infty\ \le\ T^*_{\infty,\mathrm P}\ \le\ 0.74921 ;
\]
for $T^*_{\infty,\mathrm P}$ this is sharpened to $0.58\le T^*_{\infty,\mathrm P}\le0.60$ by
Corollary~\ref{cor:Tarch}, and for $T^*_\infty$ to $T^*_\infty\le0.4608<0.58\le T^*_{\infty,\mathrm P}$ by
Proposition~\ref{prop:archneg}.
\end{proposition}
\begin{proof}
(i) $K$, $C_0$ and the two Tate moments are $T$-free, and an infimum over a growing set is
nonincreasing. (ii) $\PT\subset L^2(\IT)$, so the restricted infimum is the larger one; the proof of
Corollary~\ref{cor:mono} bounds the bottom of $\Lev$ on the codimension-two subspace $\PT$ by
$\lambda_3\le\varphi((3\pi/2T)^2)$ and never mentions the primes, and the comb contributes $0$.
(iii) and (iv) are immediate; a nonincreasing function changes sign once, and the thresholds are
suprema. No continuity is used. For $\laminf$ it holds by the proof of Lemma~\ref{lem:contT}, and for
$\laminfP$ it is not asserted.
\end{proof}
\begin{remark}
In (iii) the strict inequality $\laminf>0$ is not re-proved here, and that is the only reason
(iv) is stated for $T^*_\infty$ as an upper bound: strict positivity at $T=\frac12\log2$,
available in \cite{Yoshida1992}, would add the lower bound $\frac12\log2\le T^*_\infty$. For
$T^*_{\infty,\mathrm P}$ a two-sided bound is proved below (Corollary~\ref{cor:Tarch}).
\end{remark}
\begin{remark}
The archimedean form by itself was studied by Connes and Consani \cite{ConnesConsani2021}, who prove
its positivity on a window by a Toeplitz/prolate-spheroidal compression and relate the archimedean
place to the trace formula. Proposition~\ref{prop:arch} looks at the same form from the other side:
it locates the widths beyond which the archimedean form alone stops being positive, first on
$L^2(I_T)$ and then on $\PT$, so that past $T^*_{\infty,\mathrm P}$ the prime form $\PiT$ has to
supply positivity rather than merely be dominated.
\end{remark}

The bottom eigenfunction of the archimedean form on $\PT$ is never of one sign, and this is forced
by the two Tate conditions alone.

\begin{proposition}[nodal structure of $\PT$]\label{prop:nodes}
Recall that $\mT(F)=C_0\|F\|^2+\Lev(\cF)$ is the archimedean form without the pole term on
$L^2(\IT)$, so that $\Ainf=\mT$ on $\PT$. Then:
\emph{(i)} the bottom of $\mT$ on all of $L^2(\IT)$ is a simple eigenvalue whose
eigenfunction can be taken a.e.\ strictly positive;
\emph{(ii)} say that $F\in L^2(\IT)$ has \emph{at most one sign change} if for some
$t_0\in[-T,T]$ one has $F\ge0$ a.e.\ on $(-T,t_0)$ and $F\le0$ a.e.\ on $(t_0,T)$, or the same
with the two inequalities exchanged (the endpoints $t_0=\pm T$ being the case of one sign).
Then every nonzero $F\in\PT$ has at least two sign changes in this sense; in particular no
nonzero element of $\PT$ is of one sign, and an even element has at least one nodal pair in
$(0,T)$;
\emph{(iii)} consequently the bottom eigenfunction of $\mT$ on $\PT$ is not of one sign,
and the Sturm index of $\Ainf|\PT$ is shifted by one nodal pair relative to $(i)$.
\end{proposition}
\begin{proof}
(i) For $F\in L^2(\IT)$ extended by zero, Theorem~\ref{thm:A} and Fubini give the pure-jump
Dirichlet form with killing
\[
\Lev(\cF)=\frac14\iint_{\IT\times\IT}K(|x-y|)\,|F(x)-F(y)|^2\,dx\,dy
+\frac12\int_{\IT}\kappa(x)|F(x)|^2\,dx,\qquad
\kappa(x)=\int_{|y|>T}K(|x-y|)\,dy<\infty,
\]
obtained by splitting $2\Lev(\cF)=\int_\R K(|u|)\big(\|F\|^2-\Rel\cF(u)\big)du
=\frac12\iint_{\R^2}K(|x-y|)|\widetilde F(x)-\widetilde F(y)|^2$ into the square $\IT\times\IT$
and the two strips on which the zero extension $\widetilde F$ vanishes; the divergence of $K$ at
the origin is absorbed by $|F(x)-F(y)|^2$ on the form domain, and no $\delta$-kernel exists,
$\int_0^{2T}K=+\infty$. The form is therefore a regular Dirichlet form of pure-jump type with
nonnegative jump density $K$ and killing $\kappa$, up to the additive constant $C_0\|F\|^2$,
which only shifts the spectrum. By the Beurling--Deny criteria the semigroup $e^{-t\mT}$ is
positivity preserving, and since $K>0$ everywhere it is irreducible, hence positivity improving;
Perron--Frobenius then gives a simple, a.e.\ strictly positive ground state, provided the bottom
is an eigenvalue. It is: $m$ is increasing (Lemma~\ref{lem:mbounds}(i)) and $m(\tau)\to\infty$, so
$\int_{|\tau|>R}|\Fh|^2\le2\pi\,(m(R)-C_0)^{-1}\Lev(\cF)$ uniformly on form-bounded sets, while all
competitors are supported in the fixed compact $\IT$; by the Riesz--Kolmogorov criterion the form
domain embeds compactly in $L^2(\IT)$ and the resolvent is compact.
(ii) The pair $\{e^{t/2},e^{-t/2}\}$ is a Chebyshev system of order $2$ on $(-T,T)$: a nonzero
combination $c_1e^{t/2}+c_2e^{-t/2}$ vanishes only where $e^{t}=-c_2/c_1$, i.e.\ at most once. Let
$0\ne F\in\PT$ and suppose $F$ had at most one sign change, at $t_0$. Choose
$G=c_1e^{t/2}+c_2e^{-t/2}$ vanishing exactly at $t_0$ (any nonzero $G$ of the family if
$t_0=\pm T$), with the sign of $G$ matched to that of $F$ on each side. Then $FG\ge0$ a.e.\ and
$FG$ is not a.e.\ zero, so $\int FG\ne0$; but $F\in\PT$ means $\langle F,e^{t/2}\rangle=\langle
F,e^{-t/2}\rangle=0$, whence $\int FG=0$, a contradiction. For even $F$ the two sign changes are
symmetric, so at least one lies in $(0,T)$.
(iii) Immediate from (ii).
\end{proof}

\emph{Parity sectors.} For even $F$ the two Tate moments coincide,
$A(F)=B(F)=\int F(t)\cosh\frac t2\,dt$, and for odd $F$ one has $A(F)=-B(F)=\int F(t)\sinh\frac t2\,dt$.
Write $\PT^{\mathrm{ev}}$ and $\PT^{\mathrm{odd}}$ for the even and the odd elements of $\PT$
($\PT^{\mathrm{ev}}$ is the space of Theorem~\ref{thm:arena} below), $\lambda_{\mathrm{ev}}(T)$ and
$\lambda_{\mathrm{odd}}(T)$ for the bottoms of $\Ainf=\mT$ on their unit spheres (so
$\lambda_{\mathrm{ev}}=\lambda^{\mathrm{ev}}_1$ in the notation of Theorem~\ref{thm:arena}),
$n_{\mathrm{ev}}(T)$ and $n_{\mathrm{odd}}(T)$ for the numbers of negative eigenvalues of $\Ainf$ on them,
and $T^*_{\mathrm{ev}}=\sup\{T:\lambda_{\mathrm{ev}}(T)>0\}$, $T^*_{\mathrm{odd}}=\sup\{T:\lambda_{\mathrm{odd}}(T)>0\}$.
By Lemma~\ref{lem:parity} below, $\laminfP=\min(\lambda_{\mathrm{ev}},\lambda_{\mathrm{odd}})$
and $n_-(\Ainf|\PT)=n_{\mathrm{ev}}+n_{\mathrm{odd}}$. For a parity label $\mathrm{par}\in\{\mathrm{ev},\mathrm{odd}\}$
we write $\PT^{\mathrm{par}}$, $\lambda_{\mathrm{par}}$, $n_{\mathrm{par}}$, $T^*_{\mathrm{par}}$.

\begin{lemma}[parity splitting]\label{lem:parity}
$\PT=\PT^{\mathrm{ev}}\oplus\PT^{\mathrm{odd}}$, the sum being orthogonal both in
$L^2(\IT)$ and for $\mT$, which equals $\Ainf$ on $\PT$. Hence
\[
n_-(\Ainf|\PT)=n_{\mathrm{ev}}(T)+n_{\mathrm{odd}}(T),\qquad
\laminfP(T)=\min\big(\lambda_{\mathrm{ev}}(T),\lambda_{\mathrm{odd}}(T)\big).
\]
Here $n_-$ counts negative eigenvalues with multiplicity; for the closed,
lower-bounded form $\mT$, which has compact resolvent on each of these closed subspaces (proof of
Proposition~\ref{prop:nodes}(i)), this is the largest dimension of a subspace of the form domain on
which the form is negative definite.
\end{lemma}
\begin{proof}
$A(JF)=B(F)$ and $B(JF)=A(F)$, so for $F\in\PT$ both $F_{\mathrm e}=\frac12(F+JF)$ and
$F_{\mathrm o}=\frac12(F-JF)$ have vanishing Tate moments: $\PT$ is $J$-invariant and splits. The
parts are $L^2$-orthogonal, and
$\mT(F_{\mathrm e},F_{\mathrm o})=\frac1{2\pi}\int m\,\widehat{F_{\mathrm e}}\,\overline{\widehat{F_{\mathrm o}}}=0$,
because $m$ is even, $\widehat{F_{\mathrm e}}$ is even and $\widehat{F_{\mathrm o}}$ is odd. Negative
indices add over sums that are orthogonal for the norm and for the form, and bottoms take the minimum.
\end{proof}
The definition of $\Ainf$ writes the pole term as $2A(F)B(F)$, i.e.\ on real $F$, while $L^2(\IT)$ is
complex. Nothing depends on the choice: for $F=u+iv$ with $u,v$ real, $\mT(F)=\mT(u)+\mT(v)$
and $F\in\PT$ if and only if $u,v\in\PT$ (as in the proof of Lemma~\ref{lem:strict}), so the counts and
the bottoms are the same over $\R$ and over $\C$.

\begin{lemma}[zero extension]\label{lem:zeroext}
For $0<T<T'$, extension by zero maps $\PT^{\mathrm{par}}$ isometrically into
$\mathcal P^{\mathrm{par}}_{T'}$ with the same values of $\mT$. Hence $n_{\mathrm{par}}$ is
nondecreasing and $\lambda_{\mathrm{par}}$ nonincreasing in $T$, with no range condition at either window.
\end{lemma}
\begin{proof}
The norm, the Tate moments, the parity and $\cF$ are unchanged, and $\mT(F)=C_0\|F\|^2+\Lev(\cF)$
depends on $F$ only through $\cF$ on $\R$ (Theorem~\ref{thm:A}); $\Ainf$ does not depend on $T$
(Proposition~\ref{prop:arch}(i)).
\end{proof}

\begin{lemma}[strict monotonicity]\label{lem:strict}
For every $T>0$ the bottoms $\lambda_{\mathrm{ev}}(T)$, $\lambda_{\mathrm{odd}}(T)$ and
$\laminfP(T)$ are attained, and all three are strictly decreasing on $(0,\infty)$. Consequently, for each
sector, $\lambda_{\mathrm{par}}(T)>0$ for $T<T^*_{\mathrm{par}}$ and $\lambda_{\mathrm{par}}(T)<0$ for
$T>T^*_{\mathrm{par}}$; $\lambda_{\mathrm{par}}(T)<0$ if and only if $n_{\mathrm{par}}(T)\ge1$;
$T^*_{\infty,\mathrm P}=\min(T^*_{\mathrm{ev}},T^*_{\mathrm{odd}})$; and $\laminfP$ is positive on
$(0,T^*_{\infty,\mathrm P})$ and negative on $(T^*_{\infty,\mathrm P},\infty)$.
\end{lemma}
\begin{proof}
\emph{Attainment.} $\mT-C_0\|\cdot\|^2=\Lev(c_\cdot)\ge0$ is a closed form whose domain embeds
compactly in $L^2(\IT)$ (proof of Proposition~\ref{prop:nodes}(i)); $\PT^{\mathrm{par}}$ and $\PT$ are
closed, $A$ and $B$ being continuous on $L^2(\IT)$. A minimizing sequence on the unit sphere has an
$L^2$-convergent subsequence whose limit has norm $1$ and lies in the subspace, and lower semicontinuity
of the closed form gives attainment. Since the symbol is real and even, $\mT(u+iv)=\mT(u)+\mT(v)$
for real $u,v$, and $u+iv\in\PT^{\mathrm{par}}$ if and only if $u,v\in\PT^{\mathrm{par}}$ ($A,B$ are real on
real functions), so the minimizer may be taken real.

\emph{Strictness.} Monotonicity is Proposition~\ref{prop:arch}(i), sector by sector: extension by zero
maps $\mathcal P^{\mathrm{par}}_{T'}$ isometrically into $\PT^{\mathrm{par}}$ for $T'<T$ with the same
values of $\mT$ (Lemma~\ref{lem:zeroext}). Suppose $T'<T$ and
$\lambda_{\mathrm{par}}(T)=\lambda_{\mathrm{par}}(T')=:\lambda$, and let $F^*$ be a real minimizer on
$\mathcal P^{\mathrm{par}}_{T'}$. Extended by zero it is a minimizer on $\PT^{\mathrm{par}}$, so
$\mT(F^*,G)=\lambda\langle F^*,G\rangle$ for every $G\in\PT^{\mathrm{par}}$ in the form domain. Take
real $G_0\in C_c^\infty(T',T)$ with $\int G_0e^{y/2}dy=\int G_0e^{-y/2}dy=0$ and $G=G_0+\epsilon JG_0$,
$\epsilon=\pm1$ the parity sign. Then $G\in\PT^{\mathrm{par}}$, $\langle F^*,G\rangle=0$ by disjoint
supports, and $\mT(F^*,G)=2\mT(F^*,G_0)$, since $\mT$ is $J$-invariant and
$JF^*=\epsilon F^*$. Polarizing the Dirichlet form of Proposition~\ref{prop:nodes}(i),
$\Lev(\cF)=\frac14\iint_{\R^2}K(|x-y|)|\widetilde F(x)-\widetilde F(y)|^2$, for real $F,G$ with
supports at positive distance gives
$\mT(F,G)=-\frac12\iint K(|x-y|)F(x)G(y)\,dx\,dy$, absolutely convergent. Hence
$\Phi(y):=\int_{I_{T'}}K(|x-y|)F^*(x)\,dx$ satisfies $\int\Phi G_0=0$ for every such $G_0$, so
$\Phi=\alpha e^{y/2}+\beta e^{-y/2}$ on $(T',T)$. For $y>T'\ge x$, Lemma~\ref{lem:K} and Tonelli give
$\Phi(y)=2\sum_{k\ge0}c_ke^{-(2k+\frac12)y}$ with $c_k=\int F^*(x)e^{(2k+\frac12)x}dx$,
$|c_k|\le\|F^*\|_1e^{(2k+\frac12)T'}$; so $e^{y/2}\Phi(y)=2\sum_kc_kw^k$, $w=e^{-2y}$, is a convergent
power series for $0<w<e^{-2T'}$, and the identity extends to all $y>T'$. As $y\to\infty$, $\Phi\to0$,
so $\alpha=0$; then $2\sum_kc_kw^k\equiv\beta$, so $c_k=0$ for $k\ge1$ (the value $c_0=\beta/2$ is not
needed). With $u=e^{2x}$, $c_k=\int f(u)u^k\,du$ for
$f(u)=\frac12F^*(\frac12\log u)u^{-3/4}\in L^1[e^{-2T'},e^{2T'}]$, so all moments of $uf(u)$ vanish and
$uf=0$ by Weierstrass, whence $F^*=0$, contradicting $\|F^*\|=1$. For $\laminfP$ use Lemma~\ref{lem:parity}.

\emph{Consequences.} If $T>T^*_{\mathrm{par}}$, pick $T'\in(T^*_{\mathrm{par}},T)$: then
$\lambda_{\mathrm{par}}(T')\le0$ and $\lambda_{\mathrm{par}}(T)<\lambda_{\mathrm{par}}(T')$. If
$T<T^*_{\mathrm{par}}$, some $T''>T$ has $\lambda_{\mathrm{par}}(T'')>0$, and monotonicity gives
$\lambda_{\mathrm{par}}(T)>0$. A negative vector forces $\lambda_{\mathrm{par}}<0$, and conversely an
attained negative bottom spans a negative line. The last two statements follow from
$\laminfP=\min(\lambda_{\mathrm{ev}},\lambda_{\mathrm{odd}})$, the sets $\{\lambda_{\mathrm{par}}>0\}$ being
intervals.
\end{proof}
The same argument, with the pole term of $\Ainf$ in place of the two multipliers, shows that $\laminf$
and the bottom of $\mT$ on $L^2(\IT)$ are strictly decreasing as well (attainment as above, the pole
term being a bounded form of rank two); this is not used below.

\begin{theorem}[the negative index of the archimedean form, $0<T\le0.90$]\label{thm:index}
In the notation of \S\ref{subsec:arch}:
\emph{(i)} $n_{\mathrm{ev}}(T)\le1$ and $n_{\mathrm{odd}}(T)\le1$ for $0<T\le0.90$;
$n_{\mathrm{ev}}(T)=n_{\mathrm{odd}}(T)=0$ and $\laminfP(T)>0$ for $T\le0.58$; $n_{\mathrm{odd}}(T)=0$
and $\lambda_{\mathrm{odd}}(T)>0$ for $T\le\log2$.
\emph{(ii)} $n_{\mathrm{ev}}(T)\ge1$ for $T\ge0.60$, and $n_{\mathrm{odd}}(T)\ge1$ for $T\ge0.81$.
\emph{(iii)} $0.58\le T^*_{\mathrm{ev}}\le0.60$, $\log2\le T^*_{\mathrm{odd}}\le0.81$,
$T^*_{\infty,\mathrm P}=T^*_{\mathrm{ev}}$, and
\[
n_-(\Ainf|\PT)=\begin{cases}0,&0<T<T^*_{\mathrm{ev}},\\ 1,&T^*_{\mathrm{ev}}<T<T^*_{\mathrm{odd}},\\
2,&T^*_{\mathrm{odd}}<T\le0.90.\end{cases}
\]
\emph{(iv)} At the windows: $n_-(\Ainf|\PT)=0$ at $T=\frac12\log3$ and $T=0.58$; $n_-(\Ainf|\PT)=1$ at
$T=0.60$ and $T=\log2$, the negative direction being even; $n_-(\Ainf|\PT)\in\{1,2\}$ at $T=0.80$,
with $n_{\mathrm{ev}}(0.80)=1$; and $n_-(\Ainf|\PT)=2$ at $T=0.81$, $0.82$ and $0.90$.

\noindent The proof includes a computer verification in interval arithmetic (Appendix~\ref{app:encl}).
\end{theorem}
\begin{proof}
(i) The upper certificates of Proposition~\ref{prop:haynsworth} (Appendix~\ref{app:encl}) give
$n_{\mathrm{par}}(T)\le r$ with $r=0$ at $T=0.58$ in both sectors and at $T=\log2$ in the odd sector,
and $r=1$ at $T=0.90$ in both sectors. Lemma~\ref{lem:zeroext} carries each bound to every smaller $T$,
and Lemma~\ref{lem:rzero} turns $r=0$ into $\lambda_{\mathrm{par}}>0$, at the certified window and
therefore, by monotonicity, below it. (ii) The explicit negative planes of
Proposition~\ref{prop:negplane} at $T=0.60$ (even) and $T=0.81$ (odd) give $n_{\mathrm{par}}\ge1$
there, and Lemma~\ref{lem:zeroext} carries each to every larger $T$, with no range condition at the
larger window. (iii) By (i), (ii) and Lemma~\ref{lem:strict}: $\lambda_{\mathrm{ev}}(0.58)>0$ gives
$T^*_{\mathrm{ev}}\ge0.58$, $\lambda_{\mathrm{ev}}(0.60)<0$ gives $T^*_{\mathrm{ev}}\le0.60$,
$\lambda_{\mathrm{odd}}(\log2)>0$ gives $T^*_{\mathrm{odd}}\ge\log2$, $\lambda_{\mathrm{odd}}(0.81)<0$
gives $T^*_{\mathrm{odd}}\le0.81$, and $T^*_{\infty,\mathrm P}=\min(T^*_{\mathrm{ev}},T^*_{\mathrm{odd}})
=T^*_{\mathrm{ev}}$ because $0.60<\log2$. By Lemma~\ref{lem:strict}, off the threshold points
$n_{\mathrm{par}}(T)\ge1$ exactly when $T>T^*_{\mathrm{par}}$, and $n_{\mathrm{par}}\le1$ on $(0,0.90]$
by (i); add over the sectors (Lemma~\ref{lem:parity}). (iv) is (i)--(iii) at the windows:
$\frac12\log3=0.5493<0.58$; at $0.60$ and $\log2$ the even count is $1$ (at least $1$ from $0.60$, at
most $1$ from $0.90$) and the odd count is $0$ (from $\log2$); at $0.80$ the even count is $1$ and the
odd count is at most $1$; at $0.81$, $0.82$, $0.90$ both counts are $1$.
\end{proof}
The values $n_{\mathrm{par}}(T^*_{\mathrm{par}})$
at the two thresholds are not decided, and nothing depends on them; nor is it decided which of the two
values $n_-$ takes at $T=0.80$ Theorem~\ref{thm:index} concerns the
archimedean form with the prime comb deleted; it says nothing about the sign of $\mathfrak a_T$ or
$W_T$.


\begin{corollary}[the archimedean threshold]\label{cor:Tarch}
\[
0.58\ \le\ T^*_{\infty,\mathrm P}\ \le\ 0.60 .
\]
The bottom $\laminfP$ is strictly decreasing, positive on $(0,T^*_{\infty,\mathrm P})$ and negative on
$(T^*_{\infty,\mathrm P},\infty)$, and $\laminfP(T)<0$ for every $T\ge0.60$. In particular
$\laminfP(\frac12\log3)>0$: $\Ainf\ge\laminfP(\frac12\log3)\,\|F\|^2$ with a positive constant, i.e.\
the archimedean form is positive definite on $\mathcal P_{\frac12\log3}$ by itself. The proof
includes a computer verification in interval arithmetic (Appendix~\ref{app:encl}).
\end{corollary}
\begin{proof}
Theorem~\ref{thm:index}(iii) gives
$T^*_{\infty,\mathrm P}=T^*_{\mathrm{ev}}\in[0.58,0.60]$, and Theorem~\ref{thm:index}(ii) gives
$n_{\mathrm{ev}}(T)\ge1$, hence $\laminfP(T)\le\lambda_{\mathrm{ev}}(T)<0$, for $T\ge0.60$
(Lemma~\ref{lem:strict}); the rest is Lemma~\ref{lem:strict}, with $\frac12\log3=0.5493<0.58$.
\end{proof}
Corollary~\ref{cor:Tarch} proves a sign, not a value. Its lower half is a proved lower bound on an
infimum, obtained from an explicit minorant of the symbol, an exact finite section and a Schur
complement over the tail controlled by the trace of Lemma~\ref{lem:KyFan}; it improves the bound
$0.74921$ of Proposition~\ref{prop:arch}(iv).

On all of $L^2(\IT)$ the pole term $2A(F)B(F)$ is present, and one explicit odd test function makes
the archimedean form negative well before $T^*_{\infty,\mathrm P}$, improving the bound

Corollary~\ref{cor:Tarch}.

\begin{proposition}[the archimedean complete form is negative from $T=0.4608$]\label{prop:archneg}
For $T>0$ let $F^{\mathrm{odd}}_T(t)=t$ on $\IT$ and put
\[
U(T):=\frac{2T^3}3\Big(C_0+\log\coth\frac T2+2\arctan e^{-T}\Big)+\frac{2T^3(3T+7)}9
-2\Big(4T\cosh\frac T2-8\sinh\frac T2\Big)^2 .
\]
Then $\Ainf(F^{\mathrm{odd}}_T)\le U(T)$, and $U(0.4608)\in[-7.0159575,-7.0159574]\cdot10^{-6}$. Hence
$\laminf(T)<0$ for every $T\ge0.4608$, and
\[
T^*_\infty\ \le\ 0.4608\ <\ 0.58\ \le\ T^*_{\infty,\mathrm P} .
\]
The single numerical evaluation in the proof is done in interval arithmetic.
\end{proposition}
\begin{proof}
Write $F=F^{\mathrm{odd}}_T$. It is real, odd and bounded, and $|\Fh(\tau)|=O(1/|\tau|)$, so it lies in the
form domain. By the parity paragraph after Proposition~\ref{prop:nodes}, $A(F)=-B(F)=s:=\int F\sinh\frac t2$,
so the pole term is $2A(F)B(F)=-2s^2$, and $s=2\int_0^Tt\sinh\frac t2\,dt=4T\cosh\frac T2-8\sinh\frac T2$.
For $0\le u\le2T$, $\cF(u)=\int_{u-T}^Tt(t-u)\,dt=(4T^3-6T^2u+u^3)/6$, so $\|F\|^2=\cF(0)=2T^3/3$ and
$\cF(0)-\cF(u)=u(6T^2-u^2)/6\ge0$; for $u\ge2T$, $\cF(u)=0$. On $u>0$ one has $K(u)\le\frac1u+\frac12$,
since $\sinh x\ge x$ and $\cosh x\ge1$. Against the nonnegative function $\cF(0)-\cF(u)$ this gives
\begin{align*}
\Lev(\cF)&\le\int_0^{2T}\frac{u(6T^2-u^2)}6\Big(\frac1u+\frac12\Big)du+\|F\|^2\int_{2T}^\infty K\\
&=\frac{2T^3(3T+7)}9+\frac{2T^3}3\Big(\log\coth\frac T2+2\arctan e^{-T}\Big),
\end{align*}
by Theorem~\ref{thm:S}(iii) for the tail, with $\pi-2\arctan e^T=2\arctan e^{-T}$.
Adding $C_0\|F\|^2$ and $-2s^2$ gives $\Ainf(F)\le U(T)$. The value of $U(0.4608)$ is obtained by an
outward-rounded interval evaluation at $40$ digits, with $C_0=-\gamma-\frac\pi2-3\log2-\log\pi$ and
$2\arctan e^{-T}$ enclosed by its alternating series with the remainder term; the expressions for
$\cF$, the integral over the window and $s$ are also checked symbolically. So
$\laminf(0.4608)\le\Ainf(F)/\|F\|^2<0$. The rest is Proposition~\ref{prop:arch}(i), $\laminf$ being
nonincreasing, together with Corollary~\ref{cor:Tarch}.
\end{proof}
The same evaluation gives $U(0.461)\in[-3.59383,-3.59382]\cdot10^{-5}$, $U(0.47)=-1.4048\cdot10^{-3}$
and $U(\frac12\log3)=-0.0200526$. Proposition~\ref{prop:archneg} separates the two thresholds
strictly. It is the empty-comb case of Corollary~\ref{cor:convfail}(i) below: on $[0.4608,0.58]$ the
archimedean form is positive definite on $\PT$ but not on $L^2(\IT)$.

For every $T\ge0.60$, in particular at $T=\log2$, one has $\laminfP<0$, so every proof of $\lamA\ge0$
has to exhibit $\PiT(F)\ge-\Ainf(F)>0$ on the vectors on which the archimedean block is negative; at
$\log2$, by Theorem~\ref{thm:index}, this set contains a line but no two-dimensional subspace. Beyond the
threshold the prime form is not a perturbation of a positive form but the source of its positivity.

\subsection{Competing even modes}\label{subsec:arena}

No fixed finite number of archimedean modes captures the obstruction, and this is a theorem with an
explicit count. The archimedean gap along the even tower grows only logarithmically, while the prime
form is bounded by $\Aconst$ and does not decay; so the number $N^*(T)$ of even archimedean modes
whose level the prime form can reach, the competing even modes with which any truncation has to
work, is large, and grows doubly exponentially in the window.

\begin{theorem}[log-slow archimedean gaps]\label{thm:arena}
Let $\PT^{\mathrm{ev}}$ be the even part of $\PT$ and let
$\lambda^{\mathrm{ev}}_1\le\lambda^{\mathrm{ev}}_2\le\cdots$ be the eigenvalues of $\Ainf$ on it.
Then, for every $k\ge0$,
\[
\lambda^{\mathrm{ev}}_{k+1}\ \le\ C_0+\varphi\Big(\big(\tfrac{(2k+3)\pi}{2T}\big)^2\Big)
=m\Big(\frac{(2k+3)\pi}{2T}\Big)=\log k+O(1).
\]
\end{theorem}
\begin{proof}
$\Ainf$ commutes with parity, and on even functions the two Tate conditions coincide: both read
$\int F(t)\cosh\frac t2\,dt=0$. Hence $\PT^{\mathrm{ev}}$ has codimension $1$ in
$L^2_{\mathrm{ev}}(\IT)$, and the pole term vanishes on it, so $\Ainf=\mT$ there
(Proposition~\ref{prop:nodes}). Take the $k+2$ even Dirichlet modes $\cos((2j+1)\pi t/2T)$,
$j=0,\dots,k+1$, which vanish at $\pm T$ and satisfy
$\|F'\|^2\le((2k+3)\pi/2T)^2\|F\|^2$ on their span; intersecting that span with the single linear
constraint leaves a subspace of $\PT^{\mathrm{ev}}$ of dimension at least $k+1$. For $F$ in that subspace (extended by zero), Plancherel gives $\frac1{2\pi}\int\tau^2|\Fh|^2=\|F'\|^2$, and
$m-C_0=\varphi(\tau^2)$ with $\varphi$ complete Bernstein (Theorem~\ref{thm:S}(i)), hence concave and
increasing; Jensen's inequality in the variable $\tau^2$ for the probability measure
$|\Fh|^2d\tau/(2\pi\|F\|^2)$ gives
$\mT(F)/\|F\|^2\le C_0+\varphi(\|F'\|^2/\|F\|^2)\le C_0+\varphi(((2k+3)\pi/2T)^2)$, and
min--max bounds $\lambda^{\mathrm{ev}}_{k+1}$ by that number. The last equality is
$C_0+\varphi(\tau^2)=m(\tau)$ together with Lemma~\ref{lem:mbounds}(ii).
\end{proof}

\begin{corollary}[the number of competing even modes]\label{cor:arena}
Put $\tau_k=(2k+3)\pi/2T$ and let
$N^*(T)=\#\{j\ge1:\lambda^{\mathrm{ev}}_j\le\Aconst\}$ be the number of even archimedean modes
whose level the prime form, of norm at most $\Aconst$, can reach. Then
\[
N^*(T)\ \ge\ \#\{k\ge0:m(\tau_k)\le\Aconst\}\ =\ \Big\lfloor\frac{T\tau_*}{\pi}-\frac12\Big\rfloor,
\qquad \tau_*=m^{-1}(\Aconst).
\]
At $T=\log2$, where $\Aconst=2.24883$, Lemma~\ref{lem:mbounds}(ii) at the exact frequencies gives
$m(\tau_{11})\le2.2168\le\Aconst$ and $m(\tau_{12})\ge2.2596>\Aconst$, so $N^*\ge12$ and the
bound yields no more. At $T=\frac12\log3$, where $\Aconst=0.98026$, it gives
$m(\tau_1)\le0.8928\le\Aconst$ and $m(\tau_2)\ge1.1075>\Aconst$, so $N^*\ge2$. In general
$\Aconst=\sum_{n<e^{2T}}2\Lambda(n)n^{-1/2}=4e^{T}(1+o(1))$ by the prime number theorem and
partial summation, and $\tau_*=2\pi e^{\Aconst}(1+o(1))$ by Lemma~\ref{lem:mbounds}(ii), so
\[
N^*(T)\ \ge\ 2T\,e^{\Aconst}\,(1+o(1))\ =\ \exp\big(4e^{T}(1+o(1))\big):
\]
the number of competing even modes grows doubly exponentially in $T$.
\end{corollary}
\begin{proof}
By Theorem~\ref{thm:arena}, $\lambda^{\mathrm{ev}}_{k+1}\le m(\tau_k)$, so every $k$ with
$m(\tau_k)\le\Aconst$ contributes a mode. Since $m$ is increasing, $m(\tau_k)\le\Aconst$ is
$\tau_k\le\tau_*$, i.e.\ $k\le T\tau_*/\pi-\frac32$, and the number of integers $k\ge0$ with that
property is $\lfloor T\tau_*/\pi-\frac32\rfloor+1=\lfloor T\tau_*/\pi-\frac12\rfloor$. The two
window evaluations are Lemma~\ref{lem:mbounds}(ii) at $\tau_{11},\tau_{12}$ and at
$\tau_1,\tau_2$ respectively, with $\tau_k=(2k+3)\pi/2T$. The asymptotics are
$\log\Aconst$-free: $m(\tau)=\log(\tau/2\pi)+O(1/\tau)$ inverts to
$\tau_*=2\pi e^{\Aconst}(1+o(1))$, and $T\tau_*/\pi=2Te^{\Aconst}(1+o(1))$.
\end{proof}

Consequently no fixed finite truncation in the archimedean basis can decide the sign of the even
part: any argument that controls the prime form by the archimedean gap must carry at least $N^*(T)$
even modes, because up to that index the gap, growing like $\log k$, is smaller than the norm of a
perturbation that neither decays nor is compact.

\section{Obstructions to general methods}\label{sec:negative}

\subsection{Mass against value}

\begin{theorem}[logarithmic exchange]\label{thm:Vp}
Let $B_*=\Aconst-C_0$ ($=7.621$ at $T\in(\frac12\log3,\log2]$).
\emph{(i)} For every feasible $\mu$ of (P), with $\varepsilon:=\mu(\{|\tau|\ge R\})$,
$\int\sT d\mu\ge-B_*+\varepsilon\,(m(R)-m(0))$, and $m(R)-m(0)=\log R+O(1)$.
\emph{(ii)} For every $\delta>0$ there is $C_1<\infty$ such that for all large $R$ and all
$\varepsilon\in[0,1]$ there is a feasible $\mu'_{\varepsilon,R}$ with
$\mu'_{\varepsilon,R}(\{|\tau|\ge R-\sqrt R\})\ge\varepsilon(1-o(1))$ and
$\int\sT d\mu'_{\varepsilon,R}\le\lamP+\delta+\varepsilon(\log R+C_1)+o(\varepsilon)$.
\emph{(iii)} Consequently no inequality $\int\sT d\mu\ge-B_*+\varepsilon\,\omega(R)$ with
$\omega(R)/\log R\to\infty$ holds on the feasible set: the exchange rate between escaped mass and
value is $\Theta(\log R)$, and converting the depth $B_*$ into a nonnegative value through escaped
mass alone requires $R\gtrsim e^{B_*/\varepsilon}$.
\end{theorem}
\begin{proof}
(i) $\sT\ge-B_*$ everywhere and $\sT\ge m(R)-\Aconst$ on $|\tau|\ge R$; split the integral.
(ii) Pick $F_0\in C_c^\infty(\IT)$ with $c_{F_0}\in\Kp$ and $\mathfrak a_T(F_0)\le\lamP+\delta$
(such $F_0$ are dense in the primitive class), $\mu_0=\mu_{F_0}$, which has a rapidly decaying
density. Set $\mu_{\varepsilon,R}=(1-\varepsilon)\mu_0+\frac\varepsilon2(\mu_0(\cdot-R)+\mu_0(\cdot+R))$;
then $\check\mu_{\varepsilon,R}=[(1-\varepsilon)+\varepsilon\cos Ru]\check\mu_0$ is still supported in
$[-2T,2T]$, so $\mu_{\varepsilon,R}$ is feasible for the relaxation without the two moment
constraints; its moment defect is $\varepsilon\,d(R)$ with $d(R)=(\int g_0\check\mu_0\cos Ru,\int g_1\check\mu_0\cos Ru)\to0$
by Riemann--Lebesgue. By Lemma~\ref{lem:W} there are $r_0>0$, $V^*<\infty$ and, for each $w$
with $|w|\le r_0$, a relaxed-feasible $\mu_w$ of value $\le V^*$ and moment $w$. Put
$\lambda=\varepsilon|d(R)|/(r_0+\varepsilon|d(R)|)$ and choose $\mu_w$ with
$w=-\frac{1-\lambda}\lambda\varepsilon d(R)$ (of norm $r_0$); then
$\mu'_{\varepsilon,R}=(1-\lambda)\mu_{\varepsilon,R}+\lambda\mu_w$ has moment $0$, is feasible for
(P), and its value exceeds that of $\mu_{\varepsilon,R}$ by at most
$\lambda(B_*+V^*)=o(\varepsilon)$. The value of $\mu_{\varepsilon,R}$ is
$(1-\varepsilon)\mathfrak a_T(F_0)+\varepsilon\int\sT(\tau\pm R)d\mu_0$, and by
Lemma~\ref{lem:mbounds}(ii) $\int m(\tau+R)d\mu_0\le\log R+\int|\tau|d\mu_0/R+O(1)$ while the prime
part is bounded by $\Aconst$. (iii) Apply (i) to $\mu'_{\varepsilon,R}$ and let $\varepsilon=1$.
\end{proof}

\begin{remark}
Theorem~\ref{thm:Vp} excludes a rate lemma depending on $R$ alone, uniformly over the feasible
set. It does not exclude rate statements that depend on the data of a certificate class (the
lattice step $d$, the number of atoms, the $\ell^1$ size); compare \S\ref{sec:class}.
\end{remark}

\subsection{The envelope threshold}

\begin{proposition}[envelope threshold]\label{prop:env}
\emph{(i)} Every certificate $(\nu,\beta)$ satisfies $f(\tau)\ge\log\frac\tau{2\pi}-\frac1\tau-\frac1{2\tau^2}-\Aconst-\|\nu\|_{\TV}-|\beta_0|k_0(\tau)-|\beta_1|k_1(\tau)$,
so $f\ge0$ holds automatically for $\tau\ge\tau_{\mathrm{env}}:=2\pi e^{\Aconst+\|\nu\|_{\TV}+o(1)}$.
\emph{(ii)} No bound of the form $f\ge m(\tau)-E(\tau)$ with $E$ depending on $\tau$ only through
an upper bound of the prime comb can replace $\Aconst$ by a smaller constant, because
$\sup_\tau\sum_{n\in\NT}v_n\cos(\tau\log n)=\Aconst$ (Lemma~\ref{lem:mbounds}(iii)).
\emph{(iii)} $\Aconst=\sum_{n<e^{2T}}2\Lambda(n)/\sqrt n=(4+o(1))e^{T}$ by the prime number theorem,
so $\tau_{\mathrm{env}}(T)$ grows doubly exponentially in $T$. The multiplier-free part alone gives
the threshold $2\pi e^{\Aconst}$, equal to $16.7$ and $59.5$ at $T=\frac12\log3$ and $\log2$ (the last
sign change of $\sT$ is at $10.53$ and $46.13$); for a certificate of total variation $\|\nu\|_{\TV}$
the threshold is $2\pi e^{\Aconst+\|\nu\|_{\TV}}$, about $10^{40}$ for the certificate of
\S\ref{sec:cert}, which is why its tail is handled by the periodic Lemma~\ref{lem:F1} and not by this
bound.
\end{proposition}
\begin{proof}
(i) Lemma~\ref{lem:mbounds} and $|\nuh|\le\|\nu\|_{\TV}$. (ii),(iii) as stated.
\end{proof}

Thus a pointwise argument can stop at $\tau_{\mathrm{env}}(T)$ and at no earlier point, and
$2\pi e^{\Aconst}$, doubly exponential in $T$, is the shortest interval such an argument can
reach; this is the dual counterpart of \cite[Thm.~1.4]{Zhu2026}. We do not know whether every
certificate must do work on all of $[0,2\pi e^{\Aconst}]$; the periodic tail lemmas show that a
lattice certificate is checked on $[0,\tau_1+NP]$, with $\tau_1+NP\approx442$ at $T=\frac12\log3$.

\subsection{Partial prime combs}
\begin{theorem}[sub-comb negativity]\label{thm:SC}
Let $\mathcal N\subseteq\{n\ge2:\log n<2T\}$ be any set of contacts and $(v_n)_{n\in\mathcal N}$ any
nonnegative weights of total mass $A=\sum_{n\in\mathcal N}v_n$. Then the bottom of the corresponding
form on $\PT$ is at most $C_0+A+\varphi((3\pi/2T)^2)$. Since $\varphi$ is increasing with
$\varphi(0^+)=0$ and $T\mapsto\varphi((3\pi/2T)^2)$ decreases to $0$, if
\[
A<-C_0=5.3721834\ldots
\]
then the form is negative on $\PT$ for every $T>T_0(A)$, and nonpositive at $T=T_0(A)$, where $T_0(A)=\frac{3\pi}{2\sqrt{x_A}}$ and
$x_A$ is the unique solution of $\varphi(x_A)=-C_0-A$. For instance $T_0(0)=0.74921$,
$T_0(\frac12)=1.23295$, $T_0(1)=2.01848$, $T_0(2)=4.63567$, $T_0(4)=13.19461$, $T_0(5)=29.60905$.
\end{theorem}
\begin{proof}
The displayed bound is Corollary~\ref{cor:mono} with the weights $(v_n)_{n\in\mathcal N}$: the L\'evy
form $\Lev$ does not depend on the weights, its bottom on the codimension-two subspace $\PT$ is at
most $\varphi((3\pi/2T)^2)$, the comb contributes at most $A\|F\|^2$ and the constant is $C_0\|F\|^2$.
Monotonicity and the limits of $\varphi$ are immediate from the series of Theorem~\ref{thm:S}.
\end{proof}
\begin{remark}
The quantitative ancestors of Theorem~\ref{thm:SC} are the observations that the prime terms are
absent for $2T<\log2$ and that beyond that width they enter with a definite cost: Yoshida's prime-free
window \cite{Yoshida1992} and Bombieri's Theorem~12 \cite{Bombieri2000}. Theorem~\ref{thm:SC} adds
that no sub-collection of primes of total mass below $-C_0$ can carry positivity for large $T$.
\end{remark}


\begin{corollary}\label{cor:SC1}
Let the contacts be the ten primes $12049$, $124909$, $149333$, $322573$, $324011$, $373963$,
$417491$, $630229$, $782209$, $908861$, with their natural weights $v_p=2\log p/\sqrt p$ and all other
weights $0$. Then $A=0.563667\ldots$ and $T_0(A)=1.313456\ldots$. At
$T=\frac12\log908861+0.01=6.869974\ldots$, the smallest window in this family containing all ten
contacts, Theorem~\ref{thm:SC} gives
\[
\lambda_{\min}\big(A_T(v)\big|_{\PT}\big)\ \le\ -5.3721834+0.563667+2.685686=-2.122831 .
\]
\end{corollary}

Ten primes of size $10^5$--$10^6$, in a window that contains them, give a form that is negative by
$2.12$, about two fifths of the archimedean constant $|C_0|$. The complete comb at that window has
$\Aconst\sim4e^T\approx3.8\cdot10^3$, and its form is conjecturally nonnegative. Positivity, if it
holds, is thus a property of the complete initial segment $\{n<e^{2T}\}$: an argument that
establishes the inequality for a proper sub-comb and then adds contacts one at a time cannot work,
because every sub-comb of mass $<-C_0$ is negative in a large enough window. The mechanism is
explicit: each contact adds at most $v_n$ to the bound, the archimedean deficit $-C_0$ is a
constant, and the eigenvalue $\varphi((3\pi/2T)^2)$ of the window decays to $0$; only a comb whose
mass grows with $T$ keeps up. This concerns the operators $A_T(v)$, not $\zeta$. For comparison, the
window whose contact set is exactly the prime powers below $31$ (the sixteen numbers
$2,3,4,5,7,8,9,11,13,16,17,19,23,25,27,29$, among them the first ten primes) is
$T=\frac12\log31=1.71673$, where $\Aconst=16.640459$ exceeds $-C_0$ and Theorem~\ref{thm:SC} is
vacuous; the multiplier-free verification threshold there is
$2\pi e^{\Aconst}=1.06\cdot10^8$, against $16.7$ at $\frac12\log3$ and $59.5$ at $\log2$.



\subsection{A fixed window is a finite condition}\label{subsec:finite}
Fix $T$. For a nonnegative arithmetic function $a$ on $n\ge2$ write $v_n(a)=2a(n)/\sqrt n$ and
$\lambda_{\mathrm A}(T;a)$ for the bottom on $\PT$ of the form \eqref{eq:form} built with these weights
(as in Corollary~\ref{cor:SC1}). We say that \emph{Weil positivity holds at $T$ for $a$}, written
$\mathrm{WP}_T(a)$, if $\lambda_{\mathrm A}(T;a)\ge0$, and we write $\mathrm{WP}_T:=\mathrm{WP}_T(\Lambda)$, which is
$\lamA\ge0$. ($\mathrm{WP}_T$ is not the dual value $D_T$ of Definition~\ref{def:PD}.) For a real vector
$v=(v_n)_{n\in\NT}$, signed or not, $\mathfrak a^{(v)}_T$ denotes the form \eqref{eq:form} with the comb
$\sum_{n\in\NT}v_n\cos(\tau\log n)$. With $v(a):=(v_n(a))_{n\in\NT}$, the \emph{comb-level complete form} is
\[
W_T(a)(F):=\mathfrak a^{(v(a))}_T(F)+2\Rel\big(A(F)\overline{B(F)}\big),
\]
the form \eqref{eq:W} built with the weights of $a$, and $\lambda^{\mathrm W}_T(a)$ is its bottom on the unit sphere of
$L^2(\IT)$. So $W_T(\Lambda)=W_T$, and on $\PT$ the form $W_T(a)$ agrees with $\mathfrak a^{(v(a))}_T$. For
$a\ne\Lambda$, $W_T(a)$ is not the zero sum of $\zeta_a$: for the modifications of Theorem~\ref{thm:finmod}
below, $\zeta_a=\zeta E$ has exactly the zeros of $\zeta$, whose zero sum is $W_T$ (Lemma~\ref{lem:roughEF}),
while $W_T(a)(F)-W_T(F)=\sum_{n\in\NT}(v_n(\Lambda)-v_n(a))\cF(\log n)$, the contribution of $E'/E$, is in
general nonzero. Statements about $\lambda^{\mathrm W}_T(a)$ and $\mathrm{WP}_T(a)$ for $a\ne\Lambda$ are therefore
statements about forms, not about the zeros of $\zeta$ or of any $\zeta_a$.

\begin{proposition}[finite dependence]\label{prop:finite}
Put $v_T(a):=(v_n(a))_{n\in\NT}\in\R^{\NT}$, $|\NT|=\lceil e^{2T}\rceil-2$. Then:
\emph{(i)} $\mathrm{WP}_T(a)$ depends on $a$ only through $v_T(a)$;
\emph{(ii)} the set $\mathcal V_T:=\{v\in\R^{\NT}:\ \mathfrak a^{(v)}_T\ge0$ on $\PT\cap$ form domain$\}$ is closed
and convex, namely
\[
\mathcal V_T=\bigcap_{F}\Big\{v:\ \sum_{n\in\NT}v_n\,\cF(\log n)\le C_0+\Lev(\cF)\Big\},
\]
the intersection over the real unit vectors $F\in\PT$ in the form domain;
\emph{(iii)} $0\notin\mathcal V_T$ for $T>T_0(0)=0.74921$, indeed for $T\ge0.60$, and a vector $v\ge0$ with
$\sum_nv_n<-C_0-\varphi((3\pi/2T)^2)$ is not in $\mathcal V_T$ (for signed $v$ the same holds with
$\sum_n|v_n|$);
\emph{(iv)} for $a=\Lambda$ the numbers entering $\mathrm{WP}_T$ are exactly $\Lambda(n)$, $2\le n<e^{2T}$,
together with the one explicit function $m$ (equivalently $C_0$ and $K$), and nothing else.
\end{proposition}
\begin{proof}
(i) and (iv) are the definition of $\sT$. (ii) Real tests suffice: for $F=F_1+iF_2$ with $F_j$ real,
$|\Fh|^2=|\Fh_1|^2+|\Fh_2|^2+2\Iml(\Fh_1\overline{\Fh_2})$, and the last term is odd and integrates to $0$
against the even symbol, so the form splits as the sum of its values at $F_1,F_2$; and $A(F)=A(F_1)+iA(F_2)$
with $A(F_j)$ real, likewise for $B$, so $F\in\PT$ if and only if $F_1,F_2\in\PT$. For real unit $F$ in the
form domain, \eqref{eq:Aform} reads $\mathfrak a^{(v)}_T(F)=C_0+\Lev(\cF)-\sum v_n\cF(\log n)$ with
$\Lev(\cF)<\infty$, a closed half-space condition on $v$. (iii) The first statement is
Proposition~\ref{prop:arch}(ii) ($\laminfP<0$ for $T>T_0(0)$), and Corollary~\ref{cor:Tarch} for $T\ge0.60$; the second Theorem~\ref{thm:SC} with these
weights; for signed $v$ the comb term of that proof is bounded by $\sum|v_n|\,\|F\|^2$.
\end{proof}

\begin{lemma}[a contact of either sign]\label{lem:cpos}
Let $0<u_0<2T$. There are real $F_\pm\in\PT\cap C_c^\infty(\IT)$ with $\pm c_{F_\pm}(u_0)>0$.
\end{lemma}
\begin{proof}
Choose $0<\ell<\min(u_0,2T-u_0)$ and $J=(-T,-T+\ell)$: then $J$, $J+u_0$, $J+2u_0$ are pairwise disjoint
and $J\cup(J+u_0)\subset\IT$. For real $\phi\in C_c^\infty(J)$ put $F^\pm_\phi=\phi\pm\phi(\cdot-u_0)$. Then
$c_{F^\pm_\phi}(u_0)=\langle\phi\pm\phi_{u_0},\phi_{u_0}\pm\phi_{2u_0}\rangle=\pm\|\phi\|^2$, the three cross terms
vanishing by disjointness. The map $\phi\mapsto F^\pm_\phi$ is injective, so its image is
infinite-dimensional, and its intersection with $\PT$ has codimension at most $2$ in it and contains some
$F^\pm_\phi$ with $\phi\ne0$.
\end{proof}

\begin{theorem}[finite modifications]\label{thm:finmod}
Let $\psi_a(x):=\sum_{n\le x}a(n)$ (so $\psi_\Lambda$ is Chebyshev's function) and, for $\Rel s>1$,
$\zeta_a(s):=\exp\big(\sum_{n\ge2}a(n)n^{-s}/\log n\big)$.
\emph{(i) Deletion.} Let $N\ge5$ and $a:=\Lambda\mathbf 1_{n\ge N}$. For $T\le\frac12\log N$ the comb of $a$ is
empty, so $\mathrm{WP}_T(a)\iff\laminfP(T)\ge0$; hence $\mathrm{WP}_T(a)$ holds for every $T\le0.58$
and fails for every $T$ with $0.60\le T\le\frac12\log N$ (Theorem~\ref{thm:index}(i), Corollary~\ref{cor:Tarch}; for $T<\frac12\log2$
and for $T>T_0(0)$ already by Proposition~\ref{prop:arch}(iii) and (ii)); on $(0.58,0.60)$ it holds below
$T^*_{\infty,\mathrm P}$ and fails above it (Lemma~\ref{lem:strict}), the threshold not being located more
precisely here.
\emph{(ii) Amplification.} Let $n_0\ge2$, $c>0$ and $a:=\Lambda+c\,\mathbf 1_{n=n_0}$. For every $T$ with
$2T>\log n_0$ there is $c_*(T,n_0)<\infty$ such that $\mathrm{WP}_T(a)$ fails for all $c>c_*$. For $n_0=2$
this holds at every $T>\frac12\log2$, including $(\frac12\log2,\frac12\log5]$, where $\mathrm{WP}_T(\Lambda)$ holds
(Theorem~\ref{thm:log5A}).
\emph{(iii) What is preserved.} Both modifications keep $a\ge0$ and change $\psi_a$ by a constant for
$x\ge N$ (resp.\ $x\ge n_0$), so they preserve every partial-sum statement $\psi_a(x)=x+O(\varepsilon(x))$
with $\inf\varepsilon>0$; and $\zeta_a=\zeta\cdot E$ with
$E(s)=\exp\big(-\sum_{n<N}\Lambda(n)n^{-s}/\log n\big)$ (resp.\ $\exp(c\,n_0^{-s}/\log n_0)$) entire and zero-free,
so $\zeta_a$ continues meromorphically with \emph{exactly the zeros and the pole of $\zeta$}. Deletion also
preserves $a\le\Lambda$. Neither preserves the functional equation, $E(s)\ne E(1-s)$.
\end{theorem}
\begin{proof}
(i) $n\in\NT$ means $n<e^{2T}\le N$, so $a=0$ on $\NT$ and the form is the archimedean form on $\PT$, whose
bottom is $\laminfP(T)$ (\S\ref{subsec:arch}); Proposition~\ref{prop:arch}(iii) gives
$\laminfP\ge\laminf\ge0$ for $T<\frac12\log2$, and Proposition~\ref{prop:arch}(ii) gives
$\laminfP(T)\le C_0+\varphi((3\pi/2T)^2)<0$ for $T>T_0(0)$; Theorem~\ref{thm:index}(i) and Corollary~\ref{cor:Tarch} give $\laminfP>0$ on
$(0,0.58]$ and $\laminfP<0$ on $[0.60,\infty)$; $N\ge5$ makes $\frac12\log N\ge0.8047>T_0(0)$.
(ii) Take $F=F_+$ of Lemma~\ref{lem:cpos} with $u_0=\log n_0$. Being smooth, it lies in the form domain,
and $\mathfrak a^{(a)}_T(F)=\mathfrak a_T(F)-(2c/\sqrt{n_0})\,\cF(\log n_0)$, which is negative for
$c>c_*:=\sqrt{n_0}\max(\mathfrak a_T(F),0)/(2\cF(\log n_0))$. (iii) The exponent of $\zeta_a$ differs from that
of $\zeta$ by a finite Dirichlet polynomial, an entire function.
\end{proof}

\begin{proposition}[finitely supported combs]\label{prop:finsupp}
Let $a\ge0$ have finite support $S$ and put $A_S=\sum_{n\in S}v_n(a)$. Then
$\limsup_{T\to\infty}\lambda_{\mathrm A}(T;a)\le C_0-A_S<0$; in particular $\mathrm{WP}_T(a)$ fails for all
large $T$. This applies to every truncation $\Lambda\mathbf 1_{n<N}$, and for the empty comb it recovers
a bound consistent with $\laminfP\downarrow C_0$ (Proposition~\ref{prop:arch}(ii)).
\end{proposition}
\begin{proof}
For real $g\in C_c^\infty(-1,1)$, $g\ne0$, put $G_T(t)=g(t/T)$ and $F_T=(-\partial^2+\frac14)G_T\in C_c^\infty(\IT)$.
Two integrations by parts give $A(F_T)=\int G_T(-\partial^2+\frac14)e^{t/2}dt=0$ and likewise $B(F_T)=0$, so
$F_T\in\PT$, and $\Fh_T(\tau)=(\tau^2+\frac14)T\widehat g(T\tau)$. Let $2T>\max_S\log n$, so that
$\sigma:=m-\sum_Sv_n(a)\cos(\tau\log n)$ does not depend on $T$. Substituting $s=T\tau$,
\[
\frac{\mathfrak a^{(a)}_T(F_T)}{\|F_T\|^2}=\frac{\int\sigma(s/T)(s^2/T^2+\frac14)^2|\widehat g(s)|^2ds}{\int(s^2/T^2+\frac14)^2|\widehat g(s)|^2ds}
\ \longrightarrow\ \sigma(0)=C_0-A_S\qquad(T\to\infty)
\]
by dominated convergence: for $T\ge1$, $|\sigma(s/T)|\le C(1+\log(2+|s|))$ (Lemma~\ref{lem:mbounds} and
$|\text{comb}|\le A_S$), $(s^2/T^2+\frac14)^2\le(s^2+\frac14)^2$, $\widehat g$ is Schwartz, and $\sigma$ is
continuous at $0$.
\end{proof}

\begin{corollary}[classical prime-side hypotheses decide no window]\label{cor:void}
Let $\mathcal C$ be a class of nonnegative arithmetic functions containing $\Lambda$. If $\mathcal C$ is
closed under deletion of finitely many terms, then for every $T>T_0(0)$ some $a\in\mathcal C$ fails
$\mathrm{WP}_T(a)$, namely $a=\Lambda\mathbf 1_{n\ge N}$ with $N\ge\max(5,e^{2T})$. If $\mathcal C$ is closed under
amplification of one term, then for every $T>\frac12\log2$ some $a\in\mathcal C$ fails $\mathrm{WP}_T(a)$, namely
$a=\Lambda+c\,\mathbf 1_{n=2}$ with $c>c_*(T,2)$. The following classes contain $\Lambda$ and have these closure
properties:
\begin{itemize}
\item the class $\mathcal C_{\mathrm{PNT}}(\varepsilon)=\{a\ge0:\psi_a(x)=x+O(\varepsilon(x))\}$, for any error term
$\varepsilon$ with $\inf\varepsilon>0$ that is known for $\Lambda$ (e.g.\ de la Vall\'ee Poussin's), is closed under
both operations;
\item the class $\mathcal C_{\mathrm{zeros}}=\{a\ge0:\ \zeta_a=\zeta E,\ E$ entire and zero-free$\}$ of all $a$ with
the zeros and the pole of $\zeta$, and hence with every zero-free region and every zero-density estimate
of $\zeta$, is closed under both operations;
\item the class $\mathcal C_{\mathrm{BT}}=\{a:0\le a\le\Lambda\}$, which satisfies every upper-bound sieve
inequality satisfied by $\Lambda$, is closed under deletion.
\end{itemize}
Consequently no argument that uses about the primes only membership of $\Lambda$ in such a class (for
instance a prime number theorem with error term, a zero-free region, a zero-density estimate, or an
upper-bound sieve inequality) can prove $\mathrm{WP}_T$ at a window $T>T_0(0)$, since it would equally prove
$\mathrm{WP}_T(a)$ for the member $a$ above, which is false; for $\mathcal C_{\mathrm{PNT}}(\varepsilon)$ and
$\mathcal C_{\mathrm{zeros}}$ the same holds at every window $T>\frac12\log2$.
\end{corollary}
\begin{proof}
Theorem~\ref{thm:finmod}(i),(ii) give the failing members; the closure properties are
Theorem~\ref{thm:finmod}(iii).
\end{proof}

\subsection{Bounded remainders}
\begin{theorem}[infinite negative index]\label{thm:inertia-nogo}
Let $T>\frac12\log2$ and let $k\in L^\infty(-2T,2T)$ be real and even, with $C_k$ the operator with kernel
$k(x-y)$ on $L^2(\IT)$. For every $N\ge1$ there is an $N$-dimensional subspace of
$C_c^\infty(\IT)\cap\PT$ on which $-\langle C_kF,F\rangle+\PiT(F)<0$ for $F\ne0$. Hence
$n_-\big(-C_k+\PiT\,\big|\,\PT\cap\{\int F=0\}\big)=\infty$. Consequently, if $\mathfrak a_T=\mathfrak p+\mathfrak r$ with
$\mathfrak p\ge0$ and $\mathfrak r=-C_k+\PiT$ of this form, then $\mathfrak r$ is not bounded below by any form
with finitely many negative directions on $\PT$. By \cite[Thm.~4.7]{ConnesConsani2021} this applies to the
Sonin trace $\mathfrak p$ of Connes and Consani, whose remainder kernel is continuous.
\end{theorem}
\begin{proof}
Fix $c_0\in C_c^\infty(-1,1)$ real, $c_w=c_0(\cdot/w)$, $\beta=c_w''-\frac14c_w$. Two
integrations by parts give $\int\beta e^{\pm t/2}=0$, so translates of $\beta$ lie in $\PT$. Put
$h_j=\beta(\cdot-t_j)+\beta(\cdot-t_j-a)$, $a=\log2$, with $t_j$ chosen so that all mutual differences avoid
$\{\pm\log n\pm\{0,a\}:n\in\NT\}$ by more than $4w$. For $x=\sum c_jh_j$ the only nonzero value of $c_x(\log n)$,
$n\in\NT$, is the inner one at $\log2$, so $\PiT(x)=-v_2\|\beta\|_2^2\sum|c_j|^2=-\sqrt2\log2\,\|\beta\|_2^2\sum|c_j|^2$.
Also $|\langle C_kx,x\rangle|\le\|k\|_\infty\|x\|_1^2\le4N\|k\|_\infty\|\beta\|_1^2\sum|c_j|^2$ and
$\|\beta\|_1^2/\|\beta\|_2^2=O(w)$. For small $w$ the form is $\le-\frac12\sqrt2\log2\,\|\beta\|_2^2\sum|c_j|^2$.
Imposing $\int x=0$ costs one dimension. A bound with $k'$ negative directions fails on a $(k'+1)$-dimensional
such subspace by dimension count.
\end{proof}

\subsection{The margin decays}\label{subsec:margin}
The trivial half of Zhu's conjectural law $-\log\lamW\asymp e^{2T}$ \cite[Conj.~12.1]{Zhu2026}, namely that
the margin tends to $0$ at an explicit rate, can be proved from three facts about the zeros of $\zeta$.
They are the only numerical facts about zeros used in a proof in this paper, and nothing outside this
subsection depends on them.
\begin{itemize}
\item[(X19)] $\rho_1=\frac12+i\gamma_1$, $\gamma_1=14.134725141734\ldots$, is a simple zero on the critical line,
and $\zeta$ has no non-trivial zero $\rho\ne\rho_1,\bar\rho_1$ with $|\Iml\rho|\le21$
\cite{RosserYoheSchoenfeld1968,PlattTrudgian2021}. Only $\gamma_1\in[14.13,14.14]$ and the gap $(\gamma_1,21]$
are used.
\item[(X19$'$)] Every non-trivial zero $\rho$ with $|\Iml\rho|\le H_{\mathrm{PT}}:=3\cdot10^{12}$ lies on the
critical line \cite{PlattTrudgian2021}.
\item[(X-N)] For $t\ge2$, $\big|N(t)-\frac t{2\pi}\log\frac t{2\pi e}-\frac78\big|\le0.137\log t+0.443\log\log t+1.588$,
where $N(t)$ counts the zeros with $0<\Iml\rho\le t$ with multiplicity \cite[Thm.~19]{Rosser1941}.
\end{itemize}
Besides these, we use that the non-trivial zeros lie in $0<\Rel s<1$, and Lemma~\ref{lem:roughEF}.
\begin{lemma}[a crude zero count]\label{lem:Ncount}
$N(t)\le\frac t{2\pi}\log t$ for $t\ge21$.
\end{lemma}
\begin{proof}
$\frac t{2\pi}\log\frac t{2\pi e}=\frac t{2\pi}\log t-\frac{1+\log2\pi}{2\pi}t$ and
$\frac{1+\log2\pi}{2\pi}=0.45166\ldots$, so by (X-N)
$\frac t{2\pi}\log t-N(t)\ge0.45166\,t-\frac78-1.588-0.137\log t-0.443\log\log t$, which is $6.11\ldots$ at
$t=21$ and has derivative $0.45166-0.137/t-0.443/(t\log t)>0$ for $t\ge21$.
\end{proof}

\begin{lemma}[zero sums by counting]\label{lem:zerosum}
For $p\ge2$ and $H\ge21$, with the sum over the ordinates $\gamma=\Iml\rho>0$ of the non-trivial zeros
counted with multiplicity,
\[
\Sigma_p(H):=\sum_{\gamma>H}\gamma^{-p}\ \le\ \frac p{2\pi}\,H^{1-p}\Big[\frac{\log H}{p-1}+\frac1{(p-1)^2}\Big].
\]
\end{lemma}
\begin{proof}
$\sum_{\gamma>H}\gamma^{-p}=\int_{H^+}^\infty t^{-p}dN(t)=\big[t^{-p}N(t)\big]_H^\infty+p\int_H^\infty t^{-p-1}N(t)\,dt$;
the term at $\infty$ vanishes and the one at $H$ is $-H^{-p}N(H)\le0$ by Lemma~\ref{lem:Ncount}, which
also bounds the integral by $\frac p{2\pi}\int_H^\infty t^{-p}\log t\,dt=\frac p{2\pi}H^{1-p}\big[\frac{\log H}{p-1}+\frac1{(p-1)^2}\big]$.
\end{proof}

\begin{lemma}[the test family]\label{lem:testfam}
Fix $T>0$ and an integer $k\ge5$; let $b_k=\frac k2\mathbf 1_{[-1/k,1/k]}$ and $\omega_k=b_k^{*k}$ (the
$k$-fold convolution).
\emph{(i)} $\omega_k\ge0$, $\int\omega_k=1$, $\supp\omega_k\subseteq[-1,1]$, $\omega_k\in C^{k-2}(\R)$ with
$\omega_k^{(j)}(\pm1)=0$ for $j\le k-2$, and $\|\omega_k\|_2^2\ge\frac12$.
\emph{(ii)} $\widehat\omega_k(x)=\big(\sin(x/k)/(x/k)\big)^k$; $|\widehat\omega_k(x)|\le\min\{1,(k/|x|)^k\}$ for real
$x$, and $|\widehat\omega_k(z)|\le(k/|\Rel z|)^ke^{|\Iml z|}$ for complex $z$.
\emph{(iii)} Put $G(t)=\omega_k(t/T)$, $\varpi(c)=(c^2-\frac14)(c^2+\gamma_1^2)$, $\partial=d/dt$, and
\[
F:=\varpi(\partial)G=G''''+(\gamma_1^2-\tfrac14)G''-\tfrac{\gamma_1^2}4G .
\]
Then $F$ is real, supported in $[-T,T]$, bounded, and $F\in\PT$: four integrations by parts, the
boundary terms carrying $G^{(j)}(\pm T)=0$ for $j\le3$, give $\int Fe^{ct}dt=\varpi(c)\int Ge^{ct}dt$
for every $c\in\C$, and $\varpi(\pm\frac12)=0$.
\emph{(iv)} $\Fh(\tau)=(\tau^2+\frac14)(\tau^2-\gamma_1^2)\,T\,\widehat\omega_k(T\tau)$,
$\Phi_F(s)=\varpi(s-\frac12)\Phi_G(s)$ and $\Phi_G(\frac12+\beta+i\gamma)=T\,\widehat\omega_k\big(T(\gamma-i\beta)\big)$;
in particular $\Phi_F(\rho_1)=\Phi_F(\bar\rho_1)=0$.
\emph{(v)} With $Y:=T\gamma_1/\sqrt2$ and $\varepsilon_k(Y):=\frac1\pi k^{2k}Y^{1-2k}/(2k-1)$,
$\|F\|^2\ge\frac{\gamma_1^4}{64}\,T\,\big(\frac12-\varepsilon_k(Y)\big)$ whenever $\varepsilon_k(Y)<\frac12$.
\end{lemma}
\begin{proof}
(i) $\omega_k$ is a spline of degree $k-1$ supported in $[-1,1]$ and vanishing there to order $k-1$;
$1=(\int_{-1}^1\omega_k)^2\le2\int\omega_k^2$. (ii) $\widehat b_k(x)=\sin(x/k)/(x/k)$; $|\sin y/y|\le\min(1,1/|y|)$
on $\R$, and $|\sin(a+ib)|\le\cosh b\le e^{|b|}$. (iii) $G\in C^3$ with $G'''$ absolutely continuous
and $G''''\in L^\infty$ ($k=5$), and $G^{(j)}(\pm T)=0$ for $j\le3$; $\int G''e^{ct}=c^2\int Ge^{ct}$.
(iv) is (iii) with $c=i\tau$, resp.\ $c=s-\frac12$, and $\varpi(\pm i\gamma_1)=0$. (v) On
$|\tau|\le\gamma_1/\sqrt2$ one has $(\tau^2+\frac14)^2\ge\frac1{16}$ and $(\tau^2-\gamma_1^2)^2\ge\gamma_1^4/4$,
so $\|F\|^2=\frac1{2\pi}\int|\Fh|^2\ge\frac{\gamma_1^4}{64}\,T\,\frac1{2\pi}\int_{|x|\le Y}|\widehat\omega_k|^2$, and
$\frac1{2\pi}\int_{|x|\le Y}|\widehat\omega_k|^2\ge\|\omega_k\|_2^2-\frac1\pi\int_Y^\infty(k/x)^{2k}dx\ge\frac12-\varepsilon_k(Y)$.
\end{proof}

\begin{theorem}[the margin decays]\label{thm:margin}
Let $T>0$, $k\ge5$ an integer, $p:=2k-8$, $Y:=T\gamma_1/\sqrt2$, and suppose $\varepsilon_k(Y)<\frac12$. Then
\begin{equation}\label{eq:Bmargin}
\lamW\ \le\ \lamP\ \le\ \lamA\ \le\ \mathcal B(T,k):=
\frac{2\cdot1.0012\cdot T^2(k/T)^{2k}\big[\Sigma_p(21)+e^T\,\Sigma_p(H_{\mathrm{PT}})\big]}
{(14.13^4\,T/64)\,\big(\frac12-\varepsilon_k(Y)\big)},
\end{equation}
with $\Sigma_p(H)$ bounded by Lemma~\ref{lem:zerosum}. Consequently $\limsup_{T\to\infty}\lamA\le0$, and the
same for $\lamP$ and $\lamW$.
\end{theorem}
\begin{proof}
Take $F$ of Lemma~\ref{lem:testfam}. (1) $F\in\PT$, so $\lamA\le\mathfrak a_T(F)/\|F\|^2$, and
$\lamW\le\lamP\le\lamA$ by Theorem~\ref{thm:A} and Remark~\ref{rem:Kprim}.
(2) For $|\tau|\ge21$, $(\tau^2+\frac14)^2\le1.0012\,\tau^4$ and $(\tau^2-\gamma_1^2)^2\le\tau^4$, so
$|\Fh(\tau)|^2\le1.0012\,\tau^8T^2(k/(T|\tau|))^{2k}=O(|\tau|^{8-2k})$ with $8-2k\le-2$: $F$ is in the
form domain, which is hypothesis (a) of Lemma~\ref{lem:roughEF}; hypothesis (b) is verified in step (6).
Hence $\mathfrak a_T(F)=\sum_\rho\Phi_F(\rho)\Phi_F(1-\rho)$.
(3) $\rho_1,\bar\rho_1$ contribute $0$ (Lemma~\ref{lem:testfam}(iv)); by (X19) no other zero has $|\Iml\rho|\le21$.
(4) A zero with $21<|\gamma|\le H_{\mathrm{PT}}$ lies on the line by (X19$'$) and contributes
$|\Fh(\gamma)|^2\le1.0012\,T^2(k/T)^{2k}|\gamma|^{8-2k}$.
(5) A zero $\rho=\frac12+\beta+i\gamma$ with $|\gamma|>H_{\mathrm{PT}}$ has $|\beta|<\frac12$;
$|\varpi(\beta+i\gamma)|\le(\gamma^2+\beta^2+\frac14)(\gamma^2+\beta^2+\gamma_1^2)\le(1+10^{-22})\gamma^4$, $\varpi$ is even,
and $|\Phi_G(\rho)|,|\Phi_G(1-\rho)|\le T(k/(T|\gamma|))^ke^{T|\beta|}$ by Lemma~\ref{lem:testfam}(ii),(iv);
so $|\Phi_F(\rho)\Phi_F(1-\rho)|\le1.0012\,T^2(k/T)^{2k}|\gamma|^{8-2k}e^T$.
(6) Summing (4)--(5) over $\pm\gamma$ with Lemma~\ref{lem:zerosum} ($p=2k-8\ge2$) gives
$\sum_\rho|\Phi_F(\rho)\Phi_F(1-\rho)|\le2\cdot1.0012\,T^2(k/T)^{2k}\big[\Sigma_p(21)+e^T\Sigma_p(H_{\mathrm{PT}})\big]<\infty$,
which is hypothesis (b) of Lemma~\ref{lem:roughEF} and bounds $\mathfrak a_T(F)$.
(7) Divide by Lemma~\ref{lem:testfam}(v) with $\gamma_1^4\ge14.13^4$. The limit statement is
Corollary~\ref{cor:margin1}.
\end{proof}

\begin{corollary}[closed form]\label{cor:margin1}
For every $T\ge1$,
\begin{equation}\label{eq:closedmargin}
\lamW\ \le\ \lamP\ \le\ \lamA\ \le\ 4\cdot10^9\,T\,e^{-15.449\,T},\qquad 15.449<\tfrac{42}e=15.4494\ldots
\end{equation}
\end{corollary}
\begin{proof}
Take $k:=\lceil21T/e\rceil\ge8$ and write $k=21T/e+\theta$, $\theta\in[0,1)$. Then
$(k/21T)^{2k}=\exp\big(2k[\log(1+\theta e/21T)-1]\big)\le\exp(-42T/e+2\theta^2e/21T)\le e^{-15.4494\,T+0.259}$.
Next $(k/T)^{2k}\Sigma_p(21)\le21^9(k/21T)^{2k}\cdot\frac1{2\pi}\frac{2k-8}{2k-9}\big[\log21+\frac1{2k-9}\big]\le21^9\cdot0.5797\,(k/21T)^{2k}$
for $k\ge8$. The bound for the $H_{\mathrm{PT}}$-term is at most $10\,e^T(21/H_{\mathrm{PT}})^{2k-9}$ times that for the first, with
$21/H_{\mathrm{PT}}=7\cdot10^{-12}$ and $2k-9\ge15.449\,T-9$, hence less than $10^{-70}$ of it for $T\ge1$.
Finally $\varepsilon_k(Y)\le0.01566$ for all $T\ge1$: for fixed $k$, $\varepsilon_k$ decreases in $T$, so it
suffices to evaluate it at the left end of each $k$-range, namely $T=1$ for $k=8$ ($\varepsilon_8=0.0060$)
and $T=(k-1)e/21$ for $k\ge9$, where the largest value is $\varepsilon_9=0.01566$ at $T=8e/21$, the values
decreasing thereafter. So $\frac12-\varepsilon_k\ge0.484$, and
$\mathcal B(T,k)\le\frac{2\cdot1.0012\cdot0.5797\cdot21^9e^{0.259}}{(14.13^4/64)\cdot0.484}\,T\,e^{-15.4494\,T}(1+10^{-70})
\le3.97\cdot10^9\,T\,e^{-15.4494\,T}$.
\end{proof}

\begin{corollary}[the structure of the margin]\label{cor:margin2}
\emph{(i)} There is no $c>0$ with $\lamA\ge c$ for all $T$, nor for $\lamP$ or $\lamW$; for $T\ge1$ every
lower bound on any of the three lies below $4\cdot10^9Te^{-15.449T}$.
\emph{(ii)} The limits $L^{\mathrm A}:=\lim_{T\to\infty}\lamA$, $L^{\mathrm{prim}}:=\lim_{T\to\infty}\lamP$ and
$L^{\mathrm W}:=\lim_{T\to\infty}\lamW$ exist in $[-\infty,0]$, and $L^{\mathrm W}\le L^{\mathrm{prim}}\le L^{\mathrm A}$.
In particular \eqref{eq:target} is equivalent to $L^{\mathrm A}=0$.
\end{corollary}
\begin{proof}
(i) is Corollary~\ref{cor:margin1}. (ii) Monotone limits (Lemma~\ref{lem:monoT}) bounded above by
Theorem~\ref{thm:margin}; the last sentence follows from (ii) and monotonicity.
\end{proof}

\section{The complete form at a single window}\label{sec:complete}

\begin{lemma}[non-attainment for the complete form]\label{lem:Onat}
Let $T\ge\frac12\log2$ and $w\in\Wcl$. Then $\inf_\R f_w<\lamW$.
\end{lemma}
\begin{proof}
Weak duality gives $c:=\inf f_w\le\lamW$ (the paragraph after Definition~\ref{def:W}). Suppose $c=\lamW$. The
bottom of $W_T$ on the unit sphere of $L^2(\IT)$ is attained at some $F_\star$ (\cite[Thm.~3]{Bombieri2000},
or the compactness argument in the proof of Theorem~\ref{thm:O}, which uses only that the symbol tends to
$+\infty$). Then $\varrho:=\sT+\nuh_w+2\gh0-\lamW\ge0$ is continuous, and
$\frac1{2\pi}\int\varrho|\Fh_\star|^2=W_T(F_\star)-\lamW=0$ by Proposition~\ref{prop:P}, since
$\frac1{2\pi}\int\nuh_w|\Fh_\star|^2=\int c_{F_\star}d\nu_w=0$. As $\Fh_\star$ is entire and nonzero,
$\varrho\equiv0$ on $\R$. Taking inverse Fourier transforms on $0<|u|<2T$, where $\nu_w=0$ and $g_0=\cosh\frac u2$,
and where the constant contributes nothing: $2\cosh\frac u2=\frac12K(|u|)+\sum_{n\in\NT}\frac{v_n}2(\delta_{\log n}+\delta_{-\log n})$
there, impossible since $K(u)\sim1/u$ as $u\downarrow0$ while $\cosh$ is bounded.
\end{proof}

\begin{theorem}[flat-top certificates and the complete bottom]\label{thm:Al}
\emph{(a)} If $f_w\ge0$ on $\R$ for some $w\in\Wcl$, regular or not, then $W_T\ge0$ on the form domain;
hence $\lamW\ge0$ and $\lamA\ge0$.
\emph{(b)} Let $T\ge\frac12\log2$. Then there is $w\in\Wcl$ with $f_w\ge0$ on $\R$ if and only if
$\lamW>0$, and $w$ can then be taken in $\Wcl^{\mathrm{reg}}$ with $\inf f_w>0$. Likewise, if $\lamP>0$
there are $w\in\Wcl^{\mathrm{reg}}$ and $\beta\in\R^2$ with $\inf_\tau[f_w+(\beta_0-2)\gh0+\beta_1\gh1]>0$.
\end{theorem}
\begin{proof}
(a) $(\nu_w,2,0)$ is dual feasible and $W_T\ge\inf f_w$ on $L^2(\IT)$ (the paragraph after
Definition~\ref{def:W}), and $W_T=\mathfrak a_T$ on $\PT$. (b) ``$\Leftarrow$'' and the regular choice:
Theorem~\ref{thm:Zexh}(i). ``$\Rightarrow$'': $\lamW\ge\inf f_w\ge0$ by weak duality, and $\lamW>\inf f_w$ by
Lemma~\ref{lem:Onat}. The last sentence is Theorem~\ref{thm:Zexh}(ii).
\end{proof}

\emph{The per-window converse.} At a single window, the converse of Theorem~\ref{thm:Al}(a) is the
implication $\mathrm{WP}_T\Rightarrow\lamW\ge0$. It reduces to the signs of two numbers. Fix $T$. We work
with real functions. This suffices for bottoms and negative indices, as in the remark after
Lemma~\ref{lem:parity}: for $F=u+iv$ with $u,v$ real, $W_T(F)=W_T(u)+W_T(v)$ and
$\mathfrak a_T(F)=\mathfrak a_T(u)+\mathfrak a_T(v)$. Let $\mathcal D_T$ be the form domain of \eqref{eq:form}
and $\mathcal D^{\mathrm{par}}_T$ its elements of parity $\mathrm{par}\in\{\mathrm{ev},\mathrm{odd}\}$. With
$\chi_{\mathrm{par}}$ as in \S\ref{subsec:conv}, put
$\ell_{\mathrm{par}}(F)=\langle F,\chi_{\mathrm{par}}\rangle$, $\varepsilon_{\mathrm{ev}}=+1$ and
$\varepsilon_{\mathrm{odd}}=-1$. Then $A=\ell_{\mathrm{ev}}+\ell_{\mathrm{odd}}$ and
$B=\ell_{\mathrm{ev}}-\ell_{\mathrm{odd}}$ on real $F$ (the parity paragraph after
Proposition~\ref{prop:nodes}). Let $\mathcal Q^{\mathrm{par}}_T:=\PT^{\mathrm{par}}\cap\mathcal D_T
=\ker(\ell_{\mathrm{par}}|\mathcal D^{\mathrm{par}}_T)$, and write $\lambda^{\mathrm{par}}_{\mathrm A}(T)$ for the
bottom of $\mathfrak a_T$ on its unit sphere. This is the bottom with the primes, not the prime-free
$\lambda_{\mathrm{par}}$ of \S\ref{subsec:arch}. Write $\lambda^{\mathrm W}_{T,\mathrm{par}}$ for the bottom of
$W_T$ on the unit sphere of $\mathcal D^{\mathrm{par}}_T$. As in Lemma~\ref{lem:parity}, $n_-$ is the largest
dimension of a subspace of the form domain on which the form is negative definite. The letters $h$,
$\mathfrak g$, $\mathfrak G$, $S^{\mathrm{Sch}}$, $\mathcal N$ are local to this section.

\begin{theorem}[the parity Schur complement of the complete form]\label{thm:parityschur}
Fix $T>0$.
\emph{(o) Parity splitting.} On real $F=F_{\mathrm e}+F_{\mathrm o}\in\mathcal D_T$,
\[
W_T(F)=\big[\mathfrak a_T(F_{\mathrm e})+2\ell_{\mathrm{ev}}(F_{\mathrm e})^2\big]
+\big[\mathfrak a_T(F_{\mathrm o})-2\ell_{\mathrm{odd}}(F_{\mathrm o})^2\big],
\]
the sum being orthogonal for $L^2$, for $\mathfrak a_T$ and for $W_T$. Hence
$\lamA=\min_{\mathrm{par}}\lambda^{\mathrm{par}}_{\mathrm A}(T)$ and $\lamW=\min_{\mathrm{par}}\lambda^{\mathrm W}_{T,\mathrm{par}}$.

\emph{(a) Non-degenerate sector.} Suppose $0$ is not in the spectrum of the self-adjoint operator of the
closed form $\mathfrak a_T|\mathcal Q^{\mathrm{par}}_T$. This holds in particular if
$\lambda^{\mathrm{par}}_{\mathrm A}(T)>0$. Then there is a unique $h_{\mathrm{par}}\in\mathcal D^{\mathrm{par}}_T$
with $\ell_{\mathrm{par}}(h_{\mathrm{par}})=1$ and $\mathfrak a_T(h_{\mathrm{par}},G)=0$ for all
$G\in\mathcal Q^{\mathrm{par}}_T$. Put $\mathfrak g_{\mathrm{par}}(T):=W_T(h_{\mathrm{par}})=\mathfrak a_T(h_{\mathrm{par}})+2\varepsilon_{\mathrm{par}}$.
Then:
\begin{itemize}
\item[(a1)] $\mathcal D^{\mathrm{par}}_T=\mathcal Q^{\mathrm{par}}_T\oplus\R h_{\mathrm{par}}$, orthogonally for
$\mathfrak a_T$ and for $W_T$, and $W_T(G+xh_{\mathrm{par}})=\mathfrak a_T(G)+\mathfrak g_{\mathrm{par}}x^2$ for
$G\in\mathcal Q^{\mathrm{par}}_T$, $x\in\R$;
\item[(a2)] $n_-(W_T|\mathcal D^{\mathrm{par}}_T)=n_-(\mathfrak a_T|\mathcal Q^{\mathrm{par}}_T)+[\mathfrak g_{\mathrm{par}}<0]$.
The radical of $W_T$ on $\mathcal D^{\mathrm{par}}_T$ is $\R h_{\mathrm{par}}$ if $\mathfrak g_{\mathrm{par}}=0$ and $0$
otherwise;
\item[(a3)] if $\lambda^{\mathrm{par}}_{\mathrm A}(T)>0$, then
$\mathfrak a_T(h_{\mathrm{par}})=\min\{\mathfrak a_T(F):F\in\mathcal D^{\mathrm{par}}_T,\ \ell_{\mathrm{par}}(F)=1\}$ and
$\mathfrak g_{\mathrm{par}}=\min\{W_T(F):F\in\mathcal D^{\mathrm{par}}_T,\ \ell_{\mathrm{par}}(F)=1\}$;
\item[(a4)] $\lambda^{\mathrm W}_{T,\mathrm{par}}\le\mathfrak g_{\mathrm{par}}/\|h_{\mathrm{par}}\|^2$ and
$\|h_{\mathrm{par}}\|\ge1/\|\chi_{\mathrm{par}}\|$. If $\lambda^{\mathrm{par}}_{\mathrm A}(T)>0$ \emph{and}
$\mathfrak g_{\mathrm{par}}>0$, then
$\lambda^{\mathrm W}_{T,\mathrm{par}}\ge\frac12\min\big(\lambda^{\mathrm{par}}_{\mathrm A}(T),\mathfrak g_{\mathrm{par}}/\|h_{\mathrm{par}}\|^2\big)$.
If in addition $2\lambda^{\mathrm W}_{T,\mathrm{par}}<\lambda^{\mathrm{par}}_{\mathrm A}(T)$, then
$\lambda^{\mathrm W}_{T,\mathrm{par}}\le\mathfrak g_{\mathrm{par}}/\|h_{\mathrm{par}}\|^2\le2\lambda^{\mathrm W}_{T,\mathrm{par}}$.
\end{itemize}

\emph{(b) Degenerate sector.} Suppose $\lambda^{\mathrm{par}}_{\mathrm A}(T)=0$, and let
$\mathcal N=\ker(\mathfrak a_T|\mathcal Q^{\mathrm{par}}_T)$. It is nonzero and finite-dimensional. For each
$F_0\in\mathcal N$ there is a unique $\alpha(F_0)\in\R$ with $\mathfrak a_T(F_0,G)=\alpha(F_0)\ell_{\mathrm{par}}(G)$
for all $G\in\mathcal D^{\mathrm{par}}_T$. If $\alpha\not\equiv0$ on $\mathcal N$, then $W_T$ is unbounded below on
$\mathcal D^{\mathrm{par}}_T$, so $\lamW<0$. If $\alpha\equiv0$ on $\mathcal N$ (the \emph{range condition}), then
$\mathcal N$ lies in the radical of $W_T$ on $\mathcal D^{\mathrm{par}}_T$ and (a) applies on the $L^2$-orthogonal
complement of $\mathcal N$ in $\mathcal D^{\mathrm{par}}_T$. With the $h_{\mathrm{par}},\mathfrak g_{\mathrm{par}}$ so obtained,
$W_T\ge0$ on $\mathcal D^{\mathrm{par}}_T$ if and only if $\mathfrak g_{\mathrm{par}}\ge0$.

\emph{(c) Assembled.} Put $\mathfrak G_T:=\mathrm{diag}(\mathfrak g_{\mathrm{ev}}(T),\mathfrak g_{\mathrm{odd}}(T))$. It is
defined whenever each sector is as in (a), or as in (b) with the range condition. If $0$ is not in the
spectrum of the operator of $\mathfrak a_T|(\PT\cap\mathcal D_T)$, then
$n_-(W_T)=n_-(\mathfrak a_T|\PT)+n_-(\mathfrak G_T)$, and the radical of $W_T$ has dimension $n_0(\mathfrak G_T)$. At every window,
without further hypotheses,
\[
\lamW>0\iff\big[\lamA>0\ \text{and}\ \mathfrak G_T\succ0\big],
\]
\[
\lamW\ge0\iff\big[\mathrm{WP}_T,\ \text{the range condition in every sector with }
\lambda^{\mathrm{par}}_{\mathrm A}(T)=0,\ \text{and}\ \mathfrak G_T\succeq0\big].
\]
In particular, if $\lamA>0$, then $\lamW\ge0$ if and only if $\mathfrak g_{\mathrm{ev}}(T)\ge0$ and
$\mathfrak g_{\mathrm{odd}}(T)\ge0$.
\end{theorem}
\begin{proof}
(o) The symbol $\sT$ is even. For real even $F_{\mathrm e}$ and real odd $F_{\mathrm o}$, $\widehat{F_{\mathrm e}}$ is even and
$\widehat{F_{\mathrm o}}$ is odd, so $\mathfrak a_T(F_{\mathrm e},F_{\mathrm o})=0$ as in the proof of
Lemma~\ref{lem:parity}. Also $\ell_{\mathrm{ev}}$ vanishes on odd functions and $\ell_{\mathrm{odd}}$ on even ones, so
$A(F)B(F)=(\ell_{\mathrm{ev}}(F_{\mathrm e})+\ell_{\mathrm{odd}}(F_{\mathrm o}))(\ell_{\mathrm{ev}}(F_{\mathrm e})-\ell_{\mathrm{odd}}(F_{\mathrm o}))
=\ell_{\mathrm{ev}}(F_{\mathrm e})^2-\ell_{\mathrm{odd}}(F_{\mathrm o})^2$, with no cross term. Polarizing gives the
orthogonality. $F\in\PT$ if and only if $\ell_{\mathrm{ev}}(F_{\mathrm e})=\ell_{\mathrm{odd}}(F_{\mathrm o})=0$, and bottoms over an
orthogonal sum take the minimum.

(a) Write $\mathcal Q=\mathcal Q^{\mathrm{par}}_T$, $\ell=\ell_{\mathrm{par}}$ and $\varepsilon=\varepsilon_{\mathrm{par}}$. On
$\mathcal D^{\mathrm{par}}_T$, $W_T(F,G)=\mathfrak a_T(F,G)+2\varepsilon\,\ell(F)\ell(G)$ by (o).

\emph{Solvability.} First, $\mathfrak a_T=\mT+\PiT$ with $\|\PiT\|\le\Aconst$, since $|\cF(\log n)|\le\|F\|^2$ for each $n$ (this is Proposition~\ref{prop:coer-prime}(i) below). The form
$\mT\ge C_0\|\cdot\|^2$ is closed, with domain $\mathcal D_T$ (Theorem~\ref{thm:A}), and this domain embeds
compactly in $L^2(\IT)$ (proof of Proposition~\ref{prop:nodes}(i)). So $\mathfrak a_T$ is closed and bounded
below by $-M$, where $M:=\Aconst-C_0$. Second, $\mathcal Q$ is closed in the form norm, because $\ell$ and the
parity projection are $L^2$-continuous. Third, $\mathcal Q$ is dense in
$H:=\ker\ell\cap L^2_{\mathrm{par}}(\IT)$. Indeed, let $x\in H$, let $d_n\in\mathcal D^{\mathrm{par}}_T$ with $d_n\to x$ in
$L^2$ (smooth functions of parity $\mathrm{par}$ are dense), and let $e\in\mathcal D^{\mathrm{par}}_T$ with $\ell(e)=1$.
Then $d_n-\ell(d_n)e\in\mathcal Q$ tends to $x$. Hence $\mathfrak a_T|\mathcal Q$ is a closed, semibounded, densely
defined form in $H$ with compact resolvent. Its spectrum is discrete, and its negative spectral subspace
$\mathcal Q_-$ is finite-dimensional and $\mathfrak a_T$-orthogonal to $\mathcal Q_+:=\mathcal Q\cap\mathcal Q_-^\perp$. If
$0$ is not in the spectrum, then $\mathfrak a_T\ge\delta\|\cdot\|^2$ on $\mathcal Q_+$ for some $\delta>0$. Hence
$\mathfrak a_T(x)\ge\frac{\delta}{\delta+M+1}\big(\mathfrak a_T(x)+(M+1)\|x\|^2\big)$ there, so $\mathfrak a_T$ is coercive
on $\mathcal Q_+$ in the form norm and the Lax--Milgram lemma applies on $\mathcal Q_+$. On $\mathcal Q_-$ the form is a
negative definite finite matrix. So for every form-bounded linear functional $f$ on $\mathcal Q$ there is a unique
$Y\in\mathcal Q$ with $\mathfrak a_T(Y,\cdot)=f$ on $\mathcal Q$. Uniqueness holds because the operator has trivial
kernel. Now fix $H_0\in\mathcal D^{\mathrm{par}}_T$ with $\ell(H_0)=1$, for instance a smooth bump of parity
$\mathrm{par}$. The functional $f=\mathfrak a_T(H_0,\cdot)$ is form-bounded, by Cauchy--Schwarz for the nonnegative
form $\mathfrak a_T+M\|\cdot\|^2$. Put $h:=H_0-Y$. Then $\ell(h)=1$ and $\mathfrak a_T(h,\cdot)=0$ on $\mathcal Q$. Two such
$h$ differ by an element of $\mathcal Q$ that is $\mathfrak a_T$-orthogonal to $\mathcal Q$, which is $0$.

(a1) Any $F\in\mathcal D^{\mathrm{par}}_T$ decomposes as $F=(F-\ell(F)h)+\ell(F)h$ with $F-\ell(F)h\in\mathcal Q$. For
$G\in\mathcal Q$, $W_T(h,G)=\mathfrak a_T(h,G)+2\varepsilon\ell(h)\ell(G)=0$. Expand, using $\ell(G)=0$ and $\ell(h)=1$.

(a2) By (a1), $\mathcal D^{\mathrm{par}}_T=(\mathcal Q_-\oplus[\R h\text{ if }\mathfrak g<0])\oplus(\mathcal Q_+\oplus[\R h\text{ if }\mathfrak g\ge0])$,
orthogonally for $W_T$. The first summand is negative definite and has the stated dimension. The second
summand carries $W_T\ge0$. A negative definite subspace meets the second summand only in $0$, so its
projection onto the first summand along the second is injective. For the radical: if
$W_T(G+xh,\cdot)=0$, then $\mathfrak a_T(G,\cdot)=0$ on $\mathcal Q$, so $G=0$, and then $\mathfrak g x=0$.

(a3) If the bottom is positive, then $\mathfrak a_T(G+h)=\mathfrak a_T(G)+\mathfrak a_T(h)\ge\mathfrak a_T(h)$ and
$W_T(G+h)=\mathfrak a_T(G)+\mathfrak g\ge\mathfrak g$.

(a4) The upper bound is the Rayleigh quotient at $h$, and $1=\ell(h)\le\|h\|\|\chi_{\mathrm{par}}\|$. Put
$\lambda:=\lambda^{\mathrm{par}}_{\mathrm A}(T)$ and assume $\lambda>0$ and $\mathfrak g>0$. With
$\mu:=\min(\lambda,\mathfrak g/\|h\|^2)$,
\[
W_T(G+xh)=\mathfrak a_T(G)+\mathfrak gx^2\ge\lambda\|G\|^2+\frac{\mathfrak g}{\|h\|^2}\|xh\|^2
\ge\mu\big(\|G\|^2+\|xh\|^2\big)\ge\tfrac12\mu\|G+xh\|^2 .
\]
If also $2\lambda^{\mathrm W}_{T,\mathrm{par}}<\lambda$, then $\frac12\mu\le\lambda^{\mathrm W}_{T,\mathrm{par}}<\frac12\lambda$ forces
$\mu=\mathfrak g/\|h\|^2$. The hypothesis $\mathfrak g>0$ is used in the second inequality and cannot be dropped.

(b) The bottom $0$ is attained (compact resolvent), so $\mathcal N\ne0$, and $\mathcal N$ is finite-dimensional. For
$F_0\in\mathcal N$ and $G\in\mathcal Q$, Cauchy--Schwarz for the nonnegative form $\mathfrak a_T|\mathcal Q$ gives
$\mathfrak a_T(F_0,G)=0$. So $\mathfrak a_T(F_0,\cdot)$ is a linear functional on $\mathcal D^{\mathrm{par}}_T$ vanishing on
$\ker\ell$, hence a multiple $\alpha(F_0)\ell$. Since $\ell(F_0)=0$, we get $W_T(F_0,H)=\alpha(F_0)\ell(H)$ and
$W_T(F_0)=0$. So $W_T(xF_0+H)=W_T(H)+2x\alpha(F_0)\ell(H)$, which is unbounded below in $x$ if
$\alpha(F_0)\ne0$ and $\ell(H)=1$. Now suppose $\alpha\equiv0$. Then $\mathcal N$ is $W_T$-orthogonal to
$\mathcal D^{\mathrm{par}}_T$. Put $\mathcal D':=\mathcal D^{\mathrm{par}}_T\cap\mathcal N^\perp$ and $\mathcal Q':=\mathcal Q\cap\mathcal N^\perp$.
Then $\mathcal Q'=\ker(\ell|\mathcal D')$, and $\ell\ne0$ on $\mathcal D'$: subtract from $H_0$ its $L^2$-projection onto
$\mathcal N\subset\mathcal Q$, which $\ell$ annihilates. On $\mathcal Q'$ the form is bounded below by the next eigenvalue,
which is $>0$. So (a) applies with $\mathcal D',\mathcal Q'$ in place of $\mathcal D^{\mathrm{par}}_T,\mathcal Q$. Moreover
$\mathcal D^{\mathrm{par}}_T=\mathcal N\oplus\mathcal D'$, $W_T$-orthogonally. Conclude by (a1).

(c) By (o), everything adds over the two sectors, and the inertia formula is (a2) in each sector. If
$\lamW>0$, then $\lamA\ge\lamW>0$ (recall $\lamW\le\lamA$), so both sectors are as in (a), and
$\mathfrak g_{\mathrm{par}}=W_T(h_{\mathrm{par}})\ge\lamW\|h_{\mathrm{par}}\|^2>0$. The converse is (a4) in each sector. If $\lamW\ge0$,
then $\lamA\ge\lamW\ge0$, which is $\mathrm{WP}_T$. In a sector with bottom $0$, (b) excludes $\alpha\not\equiv0$, and
in every sector $\mathfrak g_{\mathrm{par}}=W_T(h_{\mathrm{par}})\ge0$. Conversely, (a1) and $\mathfrak a_T\ge0$ on
$\mathcal Q^{\mathrm{par}}_T$ give $W_T\ge0$ on a sector with positive bottom when $\mathfrak g_{\mathrm{par}}\ge0$, and (b) does the
same on a sector with bottom $0$.
\end{proof}

The range condition says that the column $\rho(G):=\mathfrak a_T(H_0,G)$ of the bordered form, with
$\ell_{\mathrm{par}}(H_0)=1$, vanishes on $\ker(\mathfrak a_T|\mathcal Q^{\mathrm{par}}_T)$, since $\rho(F_0)=\alpha(F_0)$. This is the
range condition of a bordered Schur complement. Because $\mathfrak a_T|\mathcal Q^{\mathrm{par}}_T$ has a spectral gap above
$0$, it is equivalent to $\rho$ lying in the range of the square root of the operator. The theorem takes
the complement over $\PT$ and not over $L^2(\IT)$. So only $\mathfrak a_T|\PT$ has to be invertible, which is
the natural hypothesis for a converse whose premise is on the primitive side.

\begin{proposition}[$\mathfrak G_T$ is the polar capacity]\label{prop:polarcap}
Assume $\mathrm{WP}_T$, and the range condition of Theorem~\ref{thm:parityschur}(b) in every
sector with $\lambda^{\mathrm{par}}_{\mathrm A}(T)=0$. Let
$S^{\mathrm{Sch}}_T:=\mathrm{diag}\big(\mathfrak a_T(h_{\mathrm{ev}}),\mathfrak a_T(h_{\mathrm{odd}})\big)$. This is the compression of
$\mathfrak a_T$ to $\mathrm{span}(h_{\mathrm{ev}},h_{\mathrm{odd}})$, the $\mathfrak a_T$-orthogonal complement of the primitive part,
i.e.\ the Schur complement of $\mathfrak a_T$ over $\PT$. Let $D_{\mathrm{pol}}:=\mathrm{diag}(2,-2)$. Then
\[
\mathfrak G_T=S^{\mathrm{Sch}}_T+D_{\mathrm{pol}} .
\]
For $(\beta,\alpha)\in\R^2$ put
$E_T(\beta,\alpha):=\inf\{W_T(F):\ F\in\mathcal D_T\ \text{real},\ \Phi_F(0)=B(F)=\beta,\ \Phi_F(1)=A(F)=\alpha\}$. Then
\[
E_T(\beta,\alpha)=(\alpha,\beta)\,\mathfrak G^{AB}_T\,(\alpha,\beta)^{\mathsf T},\qquad
\mathfrak G^{AB}_T=\frac14\begin{pmatrix}\mathfrak g_{\mathrm{ev}}+\mathfrak g_{\mathrm{odd}}&\mathfrak g_{\mathrm{ev}}-\mathfrak g_{\mathrm{odd}}\\
\mathfrak g_{\mathrm{ev}}-\mathfrak g_{\mathrm{odd}}&\mathfrak g_{\mathrm{ev}}+\mathfrak g_{\mathrm{odd}}\end{pmatrix},
\]
and the infimum is attained. In these coordinates $D_{\mathrm{pol}}$ becomes $\big(\begin{smallmatrix}0&1\\1&0\end{smallmatrix}\big)$,
the form $2\alpha\beta=2\Phi_F(0)\Phi_F(1)$, with the one negative square of Lemma~\ref{lem:polar}. Hence
$\lamW\ge0$ if and only if $E_T\ge0$ on $\R^2$. Under $\mathrm{WP}_T$, the complete form is nonnegative
exactly when the least complete Weil energy needed to prescribe the values of $\Phi_F$ at the pole
$s\in\{0,1\}$ is nonnegative.
\end{proposition}
\begin{proof}
$\mathfrak G_T=S^{\mathrm{Sch}}_T+D_{\mathrm{pol}}$ is the definition $\mathfrak g_{\mathrm{par}}=\mathfrak a_T(h_{\mathrm{par}})+2\varepsilon_{\mathrm{par}}$. In the
degenerate case with the range condition, $\mathfrak a_T(h_{\mathrm{par}})$ does not depend on the choice of $h_{\mathrm{par}}$
modulo $\mathcal N$. For real $F$, $\Phi_F(1)=A(F)$ and $\Phi_F(0)=B(F)$ (proof of Lemma~\ref{lem:roughEF}). The
constraints are $\ell_{\mathrm{ev}}(F_{\mathrm e})=c:=\frac{\alpha+\beta}2$ and $\ell_{\mathrm{odd}}(F_{\mathrm o})=s:=\frac{\alpha-\beta}2$,
and they split over the sectors, as does $W_T$ by Theorem~\ref{thm:parityschur}(o). In each sector,
Theorem~\ref{thm:parityschur}(a1) (or (b)) gives $W_T(G+xh_{\mathrm{par}})=\mathfrak a_T(G)+\mathfrak g_{\mathrm{par}}x^2$ with $x$ the
prescribed value of $\ell_{\mathrm{par}}$ and $\mathfrak a_T(G)\ge0$ by $\mathrm{WP}_T$. So the sector minimum is
$\mathfrak g_{\mathrm{par}}x^2$, attained at $G=0$. Hence $E_T=\mathfrak g_{\mathrm{ev}}c^2+\mathfrak g_{\mathrm{odd}}s^2$, which is the displayed
matrix. The form $2c^2-2s^2=2\alpha\beta$ gives $D_{\mathrm{pol}}$ in the new coordinates. The last statement is
Theorem~\ref{thm:parityschur}(c), since $E_T\ge0$ if and only if $\mathfrak G_T\succeq0$.
\end{proof}

The converse is therefore the pair of inequalities $\mathfrak a_T(h_{\mathrm{odd}})\ge2$ and
$\mathfrak a_T(h_{\mathrm{ev}})\ge-2$. In the odd direction the primitive Schur complement must beat the negative
polar square. In the even direction the positive polar square may absorb a negative Schur value down
to $-2$. Neither is automatic; Corollary~\ref{cor:convfail} below shows that the pair is not implied by
$\mathrm{WP}_T$ at the level of combs.


\section{Certified positivity up to \texorpdfstring{$T=\tfrac12\log5$}{T = (1/2) log 5}}\label{sec:log5}

Put $T_5:=\frac12\log5=0.804718\ldots$. Since $\NT=\{n\ge2:\log n<2T\}$, the contact set is
$\{2,3,4\}$ for $\frac12\log4\le T\le T_5$, and $5$ enters only for $T>T_5$; at $T_5$,
$\Aconst=v_2+v_3+v_4=2.94186\ldots$. The window $T_5$ is therefore the largest one with this contact set,
and it lies slightly beyond the range $T\le0.8$ of \cite[Thm.~1.2]{Zhu2026}. The notation
$\chi_{\mathrm{par}},\ell_{\mathrm{par}},\varepsilon_{\mathrm{par}},\mathcal Q^{\mathrm{par}}_T,
\mathfrak g_{\mathrm{par}},\mathfrak G_T$ is that of \S\ref{sec:complete}, and
$\lambda^{\mathrm W}_{T,\mathrm{par}}$ is the bottom of $W_T$ on the unit sphere of
$\mathcal D^{\mathrm{par}}_T$ (Theorem~\ref{thm:parityschur}(o)).

\begin{theorem}[the primitive form at $\frac12\log5$]\label{thm:log5A}
$\lambda^{\mathrm{ev}}_{\mathrm A}(T_5)\ge2.5\cdot10^{-12}$ and
$\lambda^{\mathrm{odd}}_{\mathrm A}(T_5)\ge3.0\cdot10^{-10}$. Hence
$\lamA\ge2.5\cdot10^{-12}$ for every $0<T\le\frac12\log5$.
\end{theorem}

\noindent The proof includes a computer verification in ball arithmetic (Appendix~\ref{subsec:encllog5}).

\begin{proof}
By Theorem~\ref{thm:parityschur}(o), $\lambda_{\mathrm A}(T_5)$ is the minimum of the two
parity bottoms, which are bounded below by Proposition~\ref{prop:lowercert} with $c=0$ and the
certificate data of Appendix~\ref{subsec:encllog5}. The bound for $T<T_5$ follows from
Lemma~\ref{lem:monoT}.
\end{proof}

\begin{theorem}[the complete form at $\frac12\log5$]\label{thm:log5W}
$W_{T_5}\ge10^{-18}\|\cdot\|^2$ on $\mathcal D^{\mathrm{ev}}_{T_5}$ and $W_{T_5}\ge10^{-16}\|\cdot\|^2$ on
$\mathcal D^{\mathrm{odd}}_{T_5}$. Hence $\lamW\ge10^{-18}$ for every $0<T\le\frac12\log5$, and
$\mathfrak G_T\succ0$ there.
\end{theorem}

\noindent The proof includes a computer verification in ball arithmetic (Appendix~\ref{subsec:encllog5}).

\begin{proof}
By Theorem~\ref{thm:parityschur}(o), $\lambda^{\mathrm W}_{T_5}=\min_{\mathrm{par}}\lambda^{\mathrm W}_{T_5,\mathrm{par}}$,
and each parity is bounded below by Proposition~\ref{prop:lowercert} with $c=-\frac12\varepsilon_{\mathrm{par}}$.
Lemma~\ref{lem:monoT} gives the bound for $T<T_5$. Since $\lamW>0$, Theorem~\ref{thm:parityschur}(c)
gives $\mathfrak G_T\succ0$.
\end{proof}

\begin{corollary}[indeterminacy for $T\le\frac12\log5$]\label{cor:KEPlog5}
$\mathrm{KEP}(T)$ is indeterminate for every $0<T\le\frac12\log5$.
\end{corollary}
\begin{proof}
By Theorem~\ref{thm:log5W}, $\lamW>0$ for $T\le T_5$, so Theorem~\ref{thm:KEPdich} applies for
$\frac12\log2\le T\le T_5$. For $0<T<\frac12\log2$, note that $\mathcal W_{T'}\subseteq\Wcl$ for $T'>T$ (a weight
flat on $[-2T',2T']$ is flat on $[-2T,2T]$, and the decay condition does not involve $T$), so a weight
$w$ with $\inf f_w>0$ at $T'=\frac12\log2$ is admissible at $T$, and $f_w$ does not depend on the window
by \eqref{eq:fw}; the perturbation of Theorem~\ref{thm:KEP1}(ii) applies, its $\chi$ being supported in
$\{|u|\ge\log2\}$, outside the window.
\end{proof}

\begin{remark}[scope]\label{rem:log5scope}
The gain over \cite[Thm.~1.2]{Zhu2026} is small, from $0.8$ to $0.80472$, and it is obtained by the
window variable alone: on $[0.8,T_5]$ the contact set is still $\{2,3,4\}$. The bound
$10^{-18}$ is consistent with Zhu's $8.9\cdot10^{-18}$ at $T=0.8$ and with the monotonicity of
$\lamW$. The complete bottom is about six orders of magnitude below the primitive one, which is the
size of the polar capacity $\mathfrak g_{\mathrm{ev}}/\|h_{\mathrm{ev}}\|^2$ of
Theorem~\ref{thm:parityschur}(a4). No certified upper bound is given; the values at which the
verification fails (Appendix~\ref{subsec:encllog5}) are not claims of negativity. Both theorems are
certificates, not a mechanism. By Theorem~\ref{thm:margin} and the doubly exponential growth of the
verification threshold (Proposition~\ref{prop:env}), the same method becomes rapidly more expensive as $T$
grows; the first window with the new contact $5$ is not treated.
\end{remark}

\emph{The converse across windows.} Theorem~\ref{thm:parityschur}(c) is a statement at one window. The
following lemma lets a single window control all smaller ones.

\begin{lemma}[monotonicity of the polar values]\label{lem:gmono}
Let $0<T<T'$ and suppose $\lambda^{\mathrm{par}}_{\mathrm A}(S)>0$ for $S=T$ and $S=T'$. Then
$\mathfrak g_{\mathrm{par}}(T')\le\mathfrak g_{\mathrm{par}}(T)$.
\end{lemma}
\begin{proof}
By Theorem~\ref{thm:parityschur}(a3),
$\mathfrak g_{\mathrm{par}}(S)=\min\{W_S(F):F\in\mathcal D^{\mathrm{par}}_S,\ \ell_{\mathrm{par}}(F)=1\}$. Extension by
zero maps $\mathcal D^{\mathrm{par}}_T$ into $\mathcal D^{\mathrm{par}}_{T'}$; it preserves the parity, the value of
$\ell_{\mathrm{par}}$ (the function $\chi_{\mathrm{par}}$ on $I_{T'}$ restricts to $\chi_{\mathrm{par}}$ on $\IT$),
and the value of $W$, by the proof of Lemma~\ref{lem:monoT}. So the feasible set at $T'$ contains that at
$T$, and the minimum does not increase.
\end{proof}

\begin{proposition}[one-window converse]\label{prop:onewindow}
Suppose $\lamA>0$ for $0<T\le T_1$. Then
\emph{(i)} $\lamW\ge0$ for all $0<T\le T_1$ if and only if $\mathfrak g_{\mathrm{ev}}(T_1)\ge0$ and
$\mathfrak g_{\mathrm{odd}}(T_1)\ge0$;
\emph{(ii)} $\lamW>0$ for all $0<T\le T_1$ if and only if $\mathfrak G_{T_1}\succ0$.
In either case $\mathfrak g_{\mathrm{par}}(T)\ge\mathfrak g_{\mathrm{par}}(T_1)$ for $0<T\le T_1$.
\end{proposition}
\begin{proof}
By Lemma~\ref{lem:monoT}, the conditions on $(0,T_1]$ are equivalent to the same conditions at $T_1$. By
Theorem~\ref{thm:parityschur}(c), since $\lamA>0$ at $T_1$, $\lambda^{\mathrm W}_{T_1}\ge0$ if and only if
$\mathfrak g_{\mathrm{ev}}(T_1),\mathfrak g_{\mathrm{odd}}(T_1)\ge0$, and $\lambda^{\mathrm W}_{T_1}>0$ if and only if
$\mathfrak G_{T_1}\succ0$. Since $\lamA>0$ implies that both parity bottoms are positive
(Theorem~\ref{thm:parityschur}(o)), Lemma~\ref{lem:gmono} gives the last statement.
\end{proof}

By Theorem~\ref{thm:log5A}, Proposition~\ref{prop:onewindow} applies with $T_1=T_5$, and
Theorem~\ref{thm:log5W} is the statement that its condition (ii) holds there.


\subsection{The per-window converse fails for modified weights}
Put $F^{\mathrm{odd}}_T(t)=t$ on $\IT$, as in Proposition~\ref{prop:archneg}, and use the comb-level forms of \S\ref{subsec:finite}.

\begin{corollary}[the converse fails for modified weights]\label{cor:convfail}
The proof uses Theorem~\ref{thm:index}, Corollary~\ref{cor:Tarch} and Theorem~\ref{thm:C}, whose proofs include
computer verifications, two further evaluations in interval arithmetic, and Theorem~\ref{thm:log5W}.
\emph{(i) Deletion.} Let $a=\Lambda\mathbf 1_{n\ge5}$. For every $T\in[0.4608,0.58]$, $\mathrm{WP}_T(a)$ holds
strictly while $\lambda^{\mathrm W}_T(a)<0$:
\[
\lambda_{\mathrm A}(T;a)=\laminfP(T)>0,\qquad\lambda^{\mathrm W}_T(a)=\laminf(T)<0 .
\]
On $[0.4608,\log2]$ the odd-sector implication fails. There the bottom of $\mathfrak a^{(v(a))}_T$ on $\PT^{\mathrm{odd}}$ is
$\lambda_{\mathrm{odd}}(T)>0$, while $W_T(a)$ is negative on an odd vector. On $(T^*_{\infty,\mathrm P},\log2]$,
$\mathrm{WP}_T(a)$ fails, so the full implication holds there vacuously and only its odd-sector form is refuted.
\emph{(ii) One contact.} Let $T=\frac12\log3$, so $\NT=\{2\}$, and for $\theta\in[0,1]$ let
$a_\theta:=\Lambda-(1-\theta)\Lambda(2)\mathbf 1_{n=2}$. Its comb on $\NT$ is $\theta v_2$. Then
$\lambda_{\mathrm A}(T;a_\theta)>0$ for every $\theta\in[0,1]$, and $\lambda^{\mathrm W}_T(a_\theta)<0$ for every
$\theta\in[0,0.8148]$. The set of $\theta\in[0,1]$ with $\lambda^{\mathrm W}_T(a_\theta)\ge0$ is an interval
$[\theta_*,1]$ with $0.8148<\theta_*<1$.
\emph{(iii) Consequences.} Let $\mathcal C$ be a class of nonnegative arithmetic functions containing $\Lambda$.
If $\mathcal C$ is closed under deletion of finitely many terms, it contains $\Lambda\mathbf 1_{n\ge5}$, for which
the implication $\mathrm{WP}_T(a)\Rightarrow\lambda^{\mathrm W}_T(a)\ge0$ fails at every $T\in[0.4608,0.58]$ by (i). If
$\mathcal C$ contains $a_\theta$ for some $\theta\le0.8148$ (for instance, if $\mathcal C$ is closed under lowering
one term within $[0,\Lambda(n)]$), the implication fails for that member at $T=\frac12\log3$ by (ii). The classes
$\mathcal C_{\mathrm{PNT}}(\varepsilon)$, $\mathcal C_{\mathrm{zeros}}$ and $\mathcal C_{\mathrm{BT}}$ of
Corollary~\ref{cor:void} are closed under deletion and under lowering one term. Consequently no argument
that uses about the primes only a property preserved by these modifications (for instance a prime number
theorem with error term, a zero-free region, a zero-density estimate, or an upper-bound sieve inequality)
can prove $\mathrm{WP}_T\Rightarrow\lamW\ge0$ at any $T\in[0.4608,0.58]$ or at $T=\frac12\log3$. At $T=\frac12\log3$
such an argument must distinguish $\Lambda(2)$ from every weight $\theta\Lambda(2)$ with $\theta\le0.8148$, that is,
from weights at least $18.5\%$ below it; nothing is claimed about weights above $\Lambda(2)$.
\end{corollary}
\begin{proof}
(i) For $T\le\log2<\frac12\log5$ one has $\NT\subseteq\{2,3,4\}$, where $a=0$. So the comb of $a$ is empty,
$W_T(a)=\Ainf$ and $\mathfrak a^{(v(a))}_T=\mT$ on $\PT$, as in the proof of Theorem~\ref{thm:finmod}(i). For
$T\le0.58$, $\laminfP(T)>0$ by Theorem~\ref{thm:index}(i), and for $T\ge0.4608$, $\laminf(T)<0$ by
Proposition~\ref{prop:archneg}. On $[0.4608,\log2]$, $\lambda_{\mathrm{odd}}(T)>0$ by Theorem~\ref{thm:index}(i). The
odd vector $F^{\mathrm{odd}}_{0.4608}$, extended by zero, lies in $L^2(\IT)$ with
$\Ainf(F^{\mathrm{odd}}_{0.4608})\le U(0.4608)<0$ ($\Ainf$ is $T$-free). $\mathrm{WP}_T(a)$ fails on
$(T^*_{\infty,\mathrm P},\log2]$ by Lemma~\ref{lem:strict}.

(ii) The comb: $\NT=\{n:\log n<\log3\}=\{2\}$ and $v_2(a_\theta)=\theta v_2$. So
$W_T(a_\theta)=\Ainf+\theta\PiT$, and $\mathfrak a^{(v(a_\theta))}_T=\mT+\theta\PiT$ on $\PT$.

\emph{The primitive side.} $\theta\mapsto\lambda_{\mathrm A}(T;a_\theta)=\inf_{F\in\PT,\|F\|=1}[\mT(F)+\theta\PiT(F)]$
is an infimum of affine functions of $\theta$. It is finite, since $\|\PiT\|\le v_2$, and hence concave on
$\R$. At $\theta=0$ it is $\laminfP(\frac12\log3)>0$ (Corollary~\ref{cor:Tarch}). At $\theta=1$ it is
$\lamA\ge8.8577\cdot10^{-5}$ (Theorem~\ref{thm:C}, which is stated for the actual comb $v_2$ at this window). A
concave function on $[0,1]$ is at least the smaller of its endpoint values.

\emph{The complete side.} Let
$F_5(t)=t-\frac{2634299}{500000}t^3+\frac{1919331}{250000}t^5$ on $\IT$, with $T=\frac12\log3$. As in
Proposition~\ref{prop:archneg}, $\Ainf(F_5)\le U_0$. Here $U_0$ is obtained by bounding $K\le\frac1u+\frac12$
against $c_{F_5}(0)-c_{F_5}(u)\ge0$ (Cauchy--Schwarz) on $(0,2T)$, by Theorem~\ref{thm:S}(iii) on $[2T,\infty)$,
and by $2A(F_5)B(F_5)=-2s(F_5)^2$. The autocorrelation $c_{F_5}$ is an exact polynomial in $(u,T)$ with rational
coefficients. Evaluating in outward-rounded interval arithmetic at $50$ digits with $T=\frac12\log3$ gives
\begin{gather*}
U_0\in[-0.00613778443569004,\,-0.00613778443569003],\\
\PiT(F_5)=-v_2c_{F_5}(\log2)\in[0.00753283781432786,\,0.00753283781432787],
\end{gather*}
and $U_0+0.8148\,\PiT(F_5)\in[-2.81845757,-2.81845756]\cdot10^{-8}$. Since $\PiT(F_5)>0$, we get
$W_T(a_\theta)(F_5)=\Ainf(F_5)+\theta\PiT(F_5)\le U_0+0.8148\,\PiT(F_5)<0$ for every $\theta\le0.8148$. (The bare
witness $F^{\mathrm{odd}}_T$ gives the same for $\theta\le0.4741$.) Finally, $\theta\mapsto\lambda^{\mathrm W}_T(a_\theta)$
is concave and finite on $\R$ for the same reason, hence continuous. At $\theta=1$ it is
$\lamW>0$ by Theorem~\ref{thm:log5W}. So
$\{\theta\in[0,1]:\lambda^{\mathrm W}_T(a_\theta)\ge0\}$ is a closed interval that contains $1$ and is disjoint from
$[0,0.8148]$. It has the form $[\theta_*,1]$, and $\theta_*<1$ by continuity.

(iii) By (i) and (ii), the implication is false for $\Lambda\mathbf 1_{n\ge5}$ (delete $n=2,3,4$) and for $a_\theta$.
For the closure properties: lowering one term keeps $0\le a\le\Lambda$; it changes $\psi_a$ by a constant for
$x\ge n$; and $\zeta_a=\zeta\exp\big(-(1-\theta)\Lambda(n)n^{-s}/\log n\big)$ with an entire zero-free factor, as in
Theorem~\ref{thm:finmod}(iii).
\end{proof}

As in Corollary~\ref{cor:void}, the form $W_T(a)$ is not the zero sum of $\zeta_a$ (\S\ref{subsec:finite}), so
$\lambda^{\mathrm W}_T(a)<0$ says nothing about the zeros of $\zeta_a$, which are those of $\zeta$; moreover $\zeta_a$
has no functional equation (Theorem~\ref{thm:finmod}(iii)). We have not examined a model with the
functional equation. For the empty comb, (i) is exactly the strict separation
$T^*_\infty<T^*_{\infty,\mathrm P}$, proved by Proposition~\ref{prop:archneg} with Corollary~\ref{cor:Tarch}

\section{A finite reduction}\label{sec:coer}

This section localizes the positivity question for $\mathfrak a_T$ on $\PT$ to an explicit subspace of
dimension $2k-2$, the \emph{piecewise Tate exponentials}, by proving that the archimedean form is coercive,
with explicit constants, on the complementary subspace of functions that satisfy the two Tate conditions on
each of $k$ consecutive blocks. The coercivity is a high-pass phenomenon ($m(\tau)\sim\log|\tau|$); the prime
form is handled by a norm bound. All forms are on $L^2(\IT)$ with its metric, $\mathfrak a_T=\mT+\PiT$ on
$\mathcal D_T$, and $\mT(F)=C_0\|F\|^2+\Lev(\cF)$ does not depend on $T$ except through its domain. In this
section $r=T/k$ denotes the half-length of a block.

\subsection{Blocks and the Tate-orthogonal subspace}\label{subsec:coer-blocks}

\begin{definition}\label{def:coer-V}
Let $k\ge1$ and let $B_1,\dots,B_k$ be the $k$ consecutive open intervals of length $2r$,
$r:=T/k$, that partition $\IT$ up to endpoints, $B_j=(-T+2(j-1)r,\,-T+2jr)$. Put
\[
\mathcal V_k:=\Big\{F\in L^2(\IT):\ \int_{B_j}F(t)e^{t/2}dt=\int_{B_j}F(t)e^{-t/2}dt=0\ \ (1\le j\le k)\Big\},
\qquad E_k:=\PT\ominus\mathcal V_k ,
\]
and $e_j^\pm:=e^{\pm t/2}\mathbf 1_{B_j}$. (Lemma~\ref{lem:coer-decomp} holds for any partition into
intervals; equal blocks are used from Lemma~\ref{lem:coer-lowpass} on.)
\end{definition}

\begin{lemma}[the decomposition $\PT=\mathcal V_k\oplus E_k$]\label{lem:coer-decomp}
\emph{(i)} $\mathcal V_k$ is a closed subspace of $\PT$. \emph{(ii)} The orthogonal complement of
$\mathcal V_k$ in $L^2(\IT)$ is $\operatorname{span}\{e_j^\pm\}$, of dimension $2k$.
\emph{(iii)} $E_k=\operatorname{span}\{e_j^\pm\}\cap\PT=\{\sum_jc_j^+e_j^++c_j^-e_j^-:\ A=B=0\}$ has
dimension $2k-2$, and $\PT=\mathcal V_k\oplus E_k$ (orthogonal in $L^2(\IT)$).
\emph{(iv)} $E_k\subset\mathcal D_T$, and $\PT\cap\mathcal D_T=(\mathcal V_k\cap\mathcal D_T)\oplus E_k$.
\end{lemma}
\begin{proof}
(i) The $2k$ conditions are continuous linear functionals; summing them over $j$ gives $A(F)=B(F)=0$.
(ii) $\mathcal V_k$ is the joint kernel of the functionals $F\mapsto\langle F,e_j^\pm\rangle$, and the $2k$
functions $e_j^\pm$ are linearly independent (disjoint supports; $e^{t/2},e^{-t/2}$ independent on an
interval). (iii) $E_k=\mathcal V_k^\perp\cap\PT$ since $\mathcal V_k\subset\PT$. The two functionals
$A,B$ restricted to $\operatorname{span}\{e_j^\pm\}$ are independent: on the two-dimensional space
$\operatorname{span}\{e_1^+,e_1^-\}$ their matrix is
$\begin{psmallmatrix}\int_{B_1}e^t&|B_1|\\|B_1|&\int_{B_1}e^{-t}\end{psmallmatrix}$, of determinant
$\int_{B_1}e^t\int_{B_1}e^{-t}-|B_1|^2>0$ (Cauchy--Schwarz, strict since $e^t$ is not constant). Hence
$\dim E_k=2k-2$. If $F\in\PT$ and $F=v+w$ with $v\in\mathcal V_k$, $w\in\mathcal V_k^\perp$, then
$w=F-v\in\PT$, so $w\in E_k$. (iv) $\widehat{e_j^\pm}(s)=\big[e^{(\pm\frac12+is)t}/(\pm\frac12+is)\big]_{\partial B_j}$,
so $|\widehat{e_j^\pm}(s)|\le2e^{T/2}(\frac14+s^2)^{-1/2}$ and $\int\log(2+|s|)|\widehat{e_j^\pm}|^2<\infty$.
Since $\mathcal D_T$ is a linear space containing $E_k$, the $\mathcal V_k$-component of an $F\in\PT\cap\mathcal D_T$
lies in $\mathcal D_T$.
\end{proof}

\subsection{Coercivity of the archimedean form on $\mathcal V_k$}\label{subsec:coer-B}
For $R>0$ let $P_R$ be the orthogonal projection of $L^2(\R)$ onto the band $\{|s|\le R\}$,
$\widehat{P_RF}=\Fh\mathbf 1_{|s|\le R}$.

\begin{lemma}[low-pass leakage of $\mathcal V_k$]\label{lem:coer-lowpass}
Let $F\in\mathcal V_k$ with equal blocks of half-length $r$, and $R>0$. Then
\[
\|P_RF\|^2\le\varepsilon_{\mathrm{lp}}(Rr)\,\|F\|^2,\qquad \varepsilon_{\mathrm{lp}}(\theta):=\frac{\cosh^2r}{6}\Big(\theta^2+\frac{r^2}4\Big)^2 .
\]
\end{lemma}
\begin{proof}
$G:=P_RF$ is entire. Fix a block $B$ with centre $c$ and put, on $B$,
$u(x):=G(c)\cosh\frac{x-c}2+2G'(c)\sinh\frac{x-c}2\in\operatorname{span}\{e^{x/2},e^{-x/2}\}$, so that
$u(c)=G(c)$, $u'(c)=G'(c)$, $u''=u/4$. Then $\varrho:=G-u$ satisfies $\varrho(c)=\varrho'(c)=0$ and
$\varrho''-\frac\varrho4=G''-\frac G4=:H$, whence $\varrho(x)=\int_c^x2\sinh\frac{x-t}2\,H(t)\,dt$. By Cauchy--Schwarz and
$2\sinh\frac z2\le z\cosh\frac z2\le z\cosh r$ ($0\le z\le r$),
$|\varrho(x)|^2\le\cosh^2r\,\frac{|x-c|^3}3\,\|H\|_{L^2(B)}^2$, so $\int_B\varrho^2\le\cosh^2r\,\frac{r^4}6\|H\|_{L^2(B)}^2$.
Since $F_B:=F\mathbf 1_B$ is orthogonal to $e^{\pm x/2}\mathbf 1_B$, $\langle F_B,G\rangle=\langle F_B,\varrho\rangle\le\|F_B\|\,\|\varrho\|_{L^2(B)}$.
Summing over the blocks and using Cauchy--Schwarz once more,
\begin{align*}
\|G\|^2=\langle F,G\rangle&\le\|F\|\Big(\sum_B\|\varrho\|_{L^2(B)}^2\Big)^{1/2}\le\|F\|\cosh r\,\frac{r^2}{\sqrt6}\,\|G''-\tfrac14G\|_{L^2(\R)}\\
&\le\|F\|\cosh r\,\frac{r^2}{\sqrt6}\Big(R^2+\frac14\Big)\|G\|,
\end{align*}
the first equality because $P_R$ is an orthogonal projection, the last inequality because
$\widehat G$ is supported in $|s|\le R$. Squaring gives the claim, as $r^4(R^2+\frac14)^2=((Rr)^2+\frac{r^2}4)^2$.
\end{proof}

\begin{theorem}[coercivity on $\mathcal V_k$]\label{thm:coer-B}
Let $k\ge1$, $r=T/k\le1$, $\kappa:=\cosh^2r$, $a:=r^2/4$, $\varepsilon_{\mathrm{lp}}$ as in
Lemma~\ref{lem:coer-lowpass}, and $\theta_*:=(\sqrt{6/\kappa}-a)^{1/2}$ (the root of $\varepsilon_{\mathrm{lp}}(\theta_*)=1$); $m_{\mathrm{lo}}$ is
the lower bound of Lemma~\ref{lem:mbounds}(ii).
Then for every $F\in\mathcal V_k\cap\mathcal D_T$,
\begin{equation}\label{eq:coer-B}
\mT(F)\ \ge\ \Big(\log\frac1r-C_B(r)\Big)\|F\|^2,\qquad
C_B(r):=\log\frac1r-C_0\,\varepsilon_{\mathrm{lp}}(r)-\frac{2\kappa}3\big[\Phi_r(\theta_*)-\Phi_r(r)\big],
\end{equation}
where $\Phi_r$ is the elementary primitive
\[
\Phi_r(\theta):=\Big(\frac{\theta^4}4+\frac{a\theta^2}2\Big)\big(\log\theta-\log(2\pi r)\big)-\frac{\theta^4}{16}-\frac{a\theta^2}4
-r\Big(\frac{\theta^3}3+a\theta\Big)-\frac{r^2\theta^2}4-\frac{ar^2}2\log\theta
\]
of $(\theta^3+a\theta)\,m_{\mathrm{lo}}(\theta/r)$. Moreover $C_B(r)\to\frac14-\log\frac{6^{1/4}}{2\pi}=1.6399$ as
$r\to0$, and $C_B(0.02)=1.6573$, $C_B(0.0619)=1.6957$, $C_B(0.1)=1.7331$ (evaluations of the closed form).
\end{theorem}
\begin{proof}
Let $\|F\|=1$ and $\mu:=\frac1{2\pi}|\Fh|^2ds$, a probability measure, so that $\mT(F)=\int m\,d\mu$ by
\eqref{eq:form} with $\NT=\emptyset$. By Lemma~\ref{lem:mbounds}(i) and \eqref{eq:m}, $m$ is even, real-analytic,
increasing on $[0,\infty)$ with $m(0)=C_0$; hence $m(s)=C_0+\int_0^{|s|}m'(\rho)d\rho$ with $m'\ge0$ and, by Tonelli,
\[
\int m\,d\mu=C_0+\int_0^\infty m'(\rho)\,\mu(|s|\ge\rho)\,d\rho\ \ge\ C_0+\int_0^{\rho_*}m'(\rho)\big(1-\varepsilon_{\mathrm{lp}}(\rho r)\big)d\rho,
\qquad\rho_*:=\theta_*/r,
\]
since $\mu(|s|\ge\rho)=1-\|P_{\rho}F\|^2+\mu(|s|=\rho)\ge1-\varepsilon_{\mathrm{lp}}(\rho r)$ by Lemma~\ref{lem:coer-lowpass} and the
integrand is nonnegative on $[0,\rho_*]$, where $\varepsilon_{\mathrm{lp}}(\rho r)\le1$. Integrating by parts on $[0,\rho_*]$, with
vanishing boundary terms ($1-\varepsilon_{\mathrm{lp}}(\theta_*)=0$ and $m(0)-C_0=0$),
\[
\int_0^{\rho_*}m'(\rho)(1-\varepsilon_{\mathrm{lp}}(\rho r))d\rho=\int_0^{\theta_*}\big(m(\theta/r)-C_0\big)\varepsilon_{\mathrm{lp}}'(\theta)\,d\theta,
\]
so that
\[
\int m\,d\mu\ \ge\ C_0\,\varepsilon_{\mathrm{lp}}(0)+\int_0^{\theta_*}m(\theta/r)\,\varepsilon_{\mathrm{lp}}'(\theta)\,d\theta ,
\]
using $\int_0^{\theta_*}\varepsilon_{\mathrm{lp}}'=1-\varepsilon_{\mathrm{lp}}(0)$. Here $\varepsilon_{\mathrm{lp}}'(\theta)=\frac{2\kappa}3\theta(\theta^2+a)\ge0$. By
Lemma~\ref{lem:mbounds}(i),(ii), $m\ge\max\{C_0,m_{\mathrm{lo}}\}$ on $(0,\infty)$; we use $m\ge C_0$ on $[0,r]$ (i.e.\ $\rho\le1$)
and $m\ge m_{\mathrm{lo}}$ on $[r,\theta_*]$ (note $m_{\mathrm{lo}}(1)=-3.34>C_0$, and $m_{\mathrm{lo}}$ is increasing, so no further loss occurs
there). Hence
\[
\int m\,d\mu\ \ge\ C_0\,\varepsilon_{\mathrm{lp}}(0)+C_0\big(\varepsilon_{\mathrm{lp}}(r)-\varepsilon_{\mathrm{lp}}(0)\big)+\int_r^{\theta_*}m_{\mathrm{lo}}(\theta/r)\,\varepsilon_{\mathrm{lp}}'(\theta)\,d\theta
=C_0\,\varepsilon_{\mathrm{lp}}(r)+\frac{2\kappa}3\big[\Phi_r(\theta_*)-\Phi_r(r)\big],
\]
which is \eqref{eq:coer-B}; the primitive is checked by differentiation. As $r\to0$: $\kappa\to1$, $a\to0$,
$\theta_*\to6^{1/4}$, $\varepsilon_{\mathrm{lp}}(r)\to0$, $\frac{2\kappa}3\Phi_r(r)\to0$, and
$\frac{2\kappa}3\Phi_r(\theta_*)=\log\theta_*-\log(2\pi r)-\frac14+o(1)$ (the $a$- and $r$-terms are $o(1)$),
whence $C_B(r)\to\log2\pi-\log6^{1/4}+\frac14$. The three values are evaluations of the closed form (double precision).
\end{proof}


\begin{proposition}[the norm of the prime form]\label{prop:coer-prime}
\emph{(i)} $|\PiT(F)|\le\Aconst\|F\|^2$ for all $F\in L^2(\IT)$ and all $T$.
\emph{(ii)} For $\ell>0$ let $L(T,\ell):=\lceil2T/\ell\rceil$. Then $|\Rel\cF(\ell)|\le\cos\frac{\pi}{L(T,\ell)+1}\,\|F\|^2$,
hence
\[
\|\PiT\|\ \le\ \sum_{n\in\NT}v_n\cos\frac\pi{L(T,\log n)+1}
=\begin{cases}
\Aconst/2, & T\le\log2,\\[2pt]
\dfrac{v_2}{\sqrt2}+\dfrac{v_3+v_4}2=1.6740, & \log2<T\le\tfrac12\log5 .
\end{cases}
\]
Here $\frac{v_2}{\sqrt2}+\frac{v_3+v_4}2=\log2+\frac{\log3}{\sqrt3}+\frac{\log2}2$; in the first case the bound is
$\frac{v_2}2=0.4901$ for $T\le\frac12\log3$ and $\frac{v_2+v_3}2=1.1244$ for $\frac12\log3<T\le\log2$. The bounds hold on all of $L^2(\IT)$, in particular on $\mathcal V_k$.
\end{proposition}
\begin{proof}
(i) is $|\cF(\log n)|\le\|F\|^2$. (ii) Let $V_\ell$ be the shift $F\mapsto F(\cdot-\ell)\mathbf 1_{\IT}$ on
$L^2(\IT)$, so that $\Rel\cF(\ell)=\frac12\langle(V_\ell+V_\ell^*)F,F\rangle$. The map
$F\mapsto\big(F(-T+y+j\ell)\big)_{0\le j<n(y)}$, $y\in[0,\ell)$, $n(y):=\#\{j\ge0:-T+y+j\ell<T\}\le L(T,\ell)$, is unitary
from $L^2(\IT)$ onto $\int^\oplus_{[0,\ell)}\C^{n(y)}dy$, and it carries $V_\ell$ to the fibrewise shift $j\mapsto j+1$.
Hence $V_\ell+V_\ell^*$ acts on the fibre $\C^{n}$ as the adjacency matrix of the path with $n$ vertices, whose
eigenvalues are $2\cos\frac{j\pi}{n+1}$, $1\le j\le n$; its norm is $2\cos\frac\pi{n+1}\le2\cos\frac\pi{L+1}$. For
$T\le\log2$ one has $2T/\log n\le2$ for all $n\ge2$, so $L=2$ and $\cos\frac\pi3=\frac12$. For $\log2<T\le\frac12\log5$:
$\NT=\{2,3,4\}$, $L(T,\log2)=3$ since $2<2T/\log2\le\log5/\log2=2.32$, and $L(T,\log3)=L(T,\log4)=2$ since
$2T\le\log5<2\log3$; $\cos\frac\pi4=\frac1{\sqrt2}$. The values use $v_2=\sqrt2\log2$, $v_3=2\log3/\sqrt3$, $v_4=\log2$.
\end{proof}
For $\frac12\log2<T<\frac12\log3$ the bound $\Aconst/2$ is the exact norm: on a set of fibres of positive
measure the fibre has length $2$ and the prime form acts there as $-v_2\begin{psmallmatrix}0&1/2\\1/2&0\end{psmallmatrix}$.
Once $T>\log2$ the orbit of $\log2$ has length $3$ and the path norm exceeds $1$, so the bound
$\|\PiT\|\le\Aconst/2$ is specific to $T\le\log2$.

\begin{corollary}[coercivity of the truncated form on $\mathcal V_k$]\label{cor:coer-L1}
Let $k\ge1$ with $r=T/k\le1$, and put $c(T,k):=\inf\{\mathfrak a_T(F):F\in\mathcal V_k\cap\mathcal D_T,\ \|F\|=1\}$. Then
\begin{equation}\label{eq:coer-L1}
c(T,k)\ \ge\ \log\frac kT-C_B\Big(\frac Tk\Big)-\|\PiT\| ,
\end{equation}
with $\|\PiT\|$ bounded by Proposition~\ref{prop:coer-prime}. In particular (evaluations of the closed form):
\begin{center}
\begin{tabular}{@{}lllll@{}}
\toprule
$T$ & $k$ & $\dim E_k$ & bound for $\|\PiT\|$ & lower bound for $c(T,k)$\\
\midrule
$\frac12\log3$ & $8$ & $14$ & $0.4901$ & $0.486$\\
$\log2$ & $12$ & $22$ & $1.1244$ & $0.035$\\
$0.69$ & $14$ & $26$ & $1.1244$ & $0.202$\\
$\frac12\log5$ & $25$ & $48$ & $1.6740$ & $0.094$\\
$\frac12\log5$ & $26$ & $50$ & $1.6740$ & $0.134$\\
$\frac12\log5$ & $30$ & $58$ & $1.6740$ & $0.281$\\
$\frac12\log5$ & $85$ & $168$ & $\Aconst=2.9420$ & $0.070$\\
\bottomrule
\end{tabular}
\end{center}
Moreover, for every $T>0$ there is $k_0(T)$ such that $c(T,k)>0$ for all $k\ge k_0(T)$; indeed the right-hand side
of \eqref{eq:coer-L1} tends to $+\infty$ with $k$, because $C_B(T/k)\to1.6399$ (Theorem~\ref{thm:coer-B}).
\end{corollary}
\begin{proof}
$\mathfrak a_T=\mT+\PiT$ on $\mathcal D_T$, Theorem~\ref{thm:coer-B} and Proposition~\ref{prop:coer-prime}.
\end{proof}


\subsection{The finite reduction}\label{subsec:coer-schur}

\begin{theorem}[reduction to the piecewise Tate exponentials]\label{thm:coer-schur}
Fix $T>0$ and $k\ge1$ with $c:=c(T,k)>0$ (Corollary~\ref{cor:coer-L1}: every $k\ge k_0(T)$). Put
$\mathcal V:=\mathcal V_k\cap\mathcal D_T$ and $M:=\Aconst-C_0$.
\emph{(i)} $\mathfrak a_T$ is an inner product on $\mathcal V$ for which $\mathcal V$ is complete, and for every
$e\in E_k$ the functional $v\mapsto\mathfrak a_T(e,v)$ is bounded on $(\mathcal V,\mathfrak a_T)$; hence there is a unique
$\mathsf Ke\in\mathcal V$ with $\mathfrak a_T(\mathsf Ke,v)=\mathfrak a_T(e,v)$ for all $v\in\mathcal V$. The map $\mathsf K:E_k\to\mathcal V$ is linear,
and $\|\mathsf K\|:=\sup\{\|\mathsf Ke\|:\|e\|=1\}<\infty$.
\emph{(ii)} For $v\in\mathcal V$ and $e\in E_k$,
\begin{equation}\label{eq:coer-schur}
\begin{split}
\mathfrak a_T(v+e)&=\mathfrak a_T(v+\mathsf Ke)+\mathfrak s_{T,k}(e),\\
\mathfrak s_{T,k}(e)&:=\mathfrak a_T(e)-\mathfrak a_T(\mathsf Ke)=\mathfrak a_T(e)-\sup_{u\in\mathcal V}\big[2\Rel\mathfrak a_T(e,u)-\mathfrak a_T(u)\big].
\end{split}
\end{equation}
Let $S_{T,k}$ be the matrix of the real quadratic form $\mathfrak s_{T,k}$ in an orthonormal basis of $E_k$ (a
real symmetric $(2k-2)\times(2k-2)$ matrix) and $\lambda_S$ its least eigenvalue. Then
\begin{equation}\label{eq:coer-sandwich}
\frac{\min\{c,\lambda_S\}}{(1+\|\mathsf K\|)^2}\ \le\ \lamA\ \le\ \min\{c,\lambda_S\}.
\end{equation}
\emph{(iii)} Consequently $\lamA>0$ if and only if $S_{T,k}\succ0$, and $\lamA\ge0$ if and only if $S_{T,k}\succeq0$.
\end{theorem}
\begin{proof}
(i) By Lemma~\ref{lem:mbounds}, $\sigma_T\ge C_0-\Aconst$, so $\mathfrak a_T+M\|\cdot\|^2\ge0$ on $\mathcal D_T$, and
by Lemma~\ref{lem:mbounds}(ii) and $m\ge C_0$ there are constants $c_1,c_2>0$ (depending on $T$) with
$c_1\int\log(2+|s|)|\Fh|^2\le\mathfrak a_T(F)+(M+1)\|F\|^2\le c_2\int\log(2+|s|)|\Fh|^2$ on $\mathcal D_T$. On $\mathcal V$,
$\mathfrak a_T\ge c\|\cdot\|^2$ gives $\mathfrak a_T(v)\ge\frac c{c+M+1}\big(\mathfrak a_T(v)+(M+1)\|v\|^2\big)$, so
$\mathfrak a_T$ is equivalent on $\mathcal V$ to the weighted norm $\int\log(2+|s|)|\widehat v|^2$, for which
$\mathcal D_T$ is complete and $\mathcal V_k\cap\mathcal D_T$ is closed ($\mathcal V_k$ is $L^2$-closed and the weighted norm
dominates the $L^2$ norm). For $e\in E_k\subset\mathcal D_T$ (Lemma~\ref{lem:coer-decomp}(iv)), the Cauchy--Schwarz
inequality for the nonnegative form $\mathfrak a_T+M\langle\cdot,\cdot\rangle$ gives
$|\mathfrak a_T(e,v)|\le(\mathfrak a_T(e)+M\|e\|^2)^{1/2}(\mathfrak a_T(v)+M\|v\|^2)^{1/2}+M\|e\|\|v\|\le C_e\,\mathfrak a_T(v)^{1/2}$,
so $\mathsf Ke$ exists by the Riesz representation theorem; $K$ is linear, and $\|\mathsf K\|<\infty$ since $\dim E_k<\infty$.
(ii) $\mathfrak a_T(v+\mathsf Ke)=\mathfrak a_T(v)+2\Rel\mathfrak a_T(\mathsf Ke,v)+\mathfrak a_T(\mathsf Ke)=\mathfrak a_T(v)+2\Rel\mathfrak a_T(e,v)+\mathfrak a_T(\mathsf Ke)$,
which rearranges to \eqref{eq:coer-schur}; the variational expression is $\mathfrak a_T(\mathsf Ke)=\sup_u[2\Rel\mathfrak a_T(\mathsf Ke,u)-\mathfrak a_T(u)]$
(maximum at $u=\mathsf Ke$). For \eqref{eq:coer-sandwich}: every $F\in\PT\cap\mathcal D_T$ is $v+e$ (Lemma~\ref{lem:coer-decomp}(iv)),
$\|F\|^2=\|v\|^2+\|e\|^2$, and $v\mapsto v+\mathsf Ke$ is a bijection of $\mathcal V$. Taking $e=0$ gives $\lamA\le c$; taking
$v=-\mathsf Ke$ gives $\mathfrak a_T(F)=\mathfrak s_{T,k}(e)$ with $\|F\|^2\ge\|e\|^2$, so $\lamA\le\lambda_S$. Conversely, if
$\lambda_S>0$, then $\mathfrak a_T(v+e)\ge c\|v+\mathsf Ke\|^2+\lambda_S\|e\|^2$, and
$\|v\|^2+\|e\|^2\le(\|v+\mathsf Ke\|+\|\mathsf K\|\|e\|)^2+\|e\|^2\le(1+\|\mathsf K\|)^2(\|v+\mathsf Ke\|^2+\|e\|^2)$, which gives the lower bound; if
$\lambda_S\le0$ the lower bound is implied by $\lamA\le\lambda_S$. (iii) follows from \eqref{eq:coer-sandwich}
and, for the nonstrict case, from $\mathfrak a_T(v+e)\ge\mathfrak s_{T,k}(e)$ together with $v=-\mathsf Ke$.
\end{proof}



\begin{remark}[what the reduction is and is not]\label{rem:coer-honest}
(1) Theorem~\ref{thm:coer-schur} is the Schur complement of $\mathfrak a_T$ with respect to
$\PT=\mathcal V_k\oplus E_k$, in the form-domain setting required because $\mathfrak a_T$ is unbounded on $L^2$.
Its content is a localization: the sign of $\lamA$ is the sign of an explicit real symmetric matrix of size
$2k-2$, with $k=25$ at $T=\frac12\log5$, indexed by the piecewise Tate exponentials $e^{\pm t/2}\mathbf 1_{B_j}$.
(2) There is no gain in precision. By \eqref{eq:coer-sandwich}, $\lamA\le\lambda_S\le(1+\|\mathsf K\|)^2\lamA$
whenever $\lambda_S\le c$, so $S_{T,k}$ inherits the decay of Theorem~\ref{thm:margin}; at $T=\frac12\log5$ its
least eigenvalue is of the size of $\lambda_{\mathrm A}(\frac12\log5)\approx3.10\cdot10^{-12}$ (numerical). The
entries of $S_{T,k}$ are not in closed form: the correction $\mathfrak a_T(\mathsf Ke)$ is a supremum over $\mathcal V$,
so a Galerkin computation in $\mathcal V$ gives upper bounds for $S_{T,k}$, and a certified lower bound needs a
near-maximizer together with a dual estimate, which is legitimate because $c>0$ is proved.
(3) Finite reductions of the window problem in sufficient (one-directional) form exist in Legendre bases
\cite[Thms.~1.1, 1.4]{Zhu2026}, \cite{Groskin2026}; Theorem~\ref{thm:coer-schur} is an exact reduction in an
explicit basis of explicit dimension, and it does not improve any certified value.
\end{remark}

\section{First contact}\label{sec:firstcontact}

This section studies a hypothetical first window at which positivity of the primitive form fails. We
prove that such a window carries a primitive null vector of definite parity, a finite operator identity in
which additive Poisson summation and the convolution identity $\Lambda*1=\log$ occur together
(Theorem~\ref{thm:PCfinite}), a faithful quadratic observation of a test function by the positive odd poles
of the Gamma factor \eqref{eq:Gammaenergy}, exact identities for the exterior collars produced by compressing
the prime shifts to the window \eqref{eq:collar}, an exact balance identity on the kernel \eqref{eq:Qbalance},
and a strict contraction of the Euler channel (Proposition~\ref{prop:GEgap}). The contraction does not by
itself give vanishing: on the kernel, the natural reverse inequality is equivalent to the vanishing of the
detector. What is missing is the annihilation statement of Problem~\ref{prob:annihilate}. Nothing in this
section extends the certified range.

\subsection{The first-contact reduction}
Throughout, $\mathfrak a_T(F,G)$ is the sesquilinear form of \eqref{eq:form}, $H_T$ the operator of
Lemma~\ref{lem:compact}, and $\mathrm{pr}_T$ the orthogonal projection of $L^2(\IT)$ onto $\PT$; $J$ preserves
$\PT$ and commutes with $\mathrm{pr}_T$. The orthogonal complement of $\PT$ in $L^2(\IT)$ is the
two-dimensional \emph{Tate space} $\mathrm{span}\{e^{t/2},e^{-t/2}\}$ (restricted to $\IT$).

\begin{lemma}[first-contact reduction]\label{lem:firstcontact}
If $\lamA<0$ for some $T$, then there is a window $T_*>\frac12\log5$ and a nonzero $F\in\mathcal P_{T_*}$ such
that $\lambda_{\mathrm A}(T)\ge0$ for $T\le T_*$, $\lambda_{\mathrm A}(T)<0$ for $T>T_*$, and
\[
 H_{T_*}\ge0,\qquad H_{T_*}F=0,\qquad JF=\varepsilon F,\quad \varepsilon\in\{\pm1\}.
\]
\end{lemma}
\begin{proof}
Let $T_*=\inf\{T:\lamA<0\}$. By Lemma~\ref{lem:monoT}, $\lamA<0$ for every $T>T_*$ and $\lamA\ge0$ for every
$T<T_*$; by Lemma~\ref{lem:contA}, $\lambda_{\mathrm A}(T_*)=0$. Since $\lambda_{\mathrm A}(\frac12\log5)>0$
(Theorem~\ref{thm:log5A}), monotonicity gives $T_*>\frac12\log5$. By Lemma~\ref{lem:compact} the value $0$ is
an eigenvalue of $H_{T_*}$, and since $\ker H_{T_*}$ is $J$-invariant it contains a nonzero vector of
definite parity.
\end{proof}

We call a nonzero $F\in\ker H_T$ with $JF=\varepsilon F$, at a window where $H_T\ge0$, a
\emph{first-contact vector}, and $\ker H_T$ the \emph{radical}. Equivalently, $F\in\PT$ satisfies
$\mathfrak a_T(F,\mathrm{pr}_TG)=0$ for every $G$ in the form domain, i.e.\ the weak equation
\begin{equation}\label{eq:localweak}
 \mathfrak a_T(F,G)=\langle r,G\rangle\qquad(G\ \text{in the form domain of}\ L^2(\IT))
\end{equation}
holds for a unique $r$ in the Tate space. More generally, for $F$ in the domain of $H_T$ we define
$r\in(\PT)^\perp$ by $\mathfrak a_T(F,G)=\langle H_TF+r,G\rangle$.

\subsection{Operators on the line}

Let $XF(t)=tF(t)$ on the Schwartz space $\mathcal S(\R)$ (the \emph{Schwartz
core}). Let $\mathcal M=m(D)$ be the Fourier multiplier with the symbol $m$ of \eqref{eq:symbol},
$(\mathcal MF)\,\widehat{}\,=m\Fh$; we call it the completed Gamma multiplier. Let
$u_\infty(\tau)=\pi^{-i\tau}\Gamma(\frac14+\frac{i\tau}2)/\Gamma(\frac14-\frac{i\tau}2)$, a
unimodular function with $u_\infty(-\tau)=\overline{u_\infty(\tau)}$ and, by Tate's local
functional equation, $\frac{d}{d\tau}\arg u_\infty=m$. Put
\[
 \mathscr F=u_\infty(D)\,J .
\]
This is the additive Fourier transform in logarithmic coordinates: if $F(t)=e^{t/2}f(e^t)$ with $f$
an even Schwartz function on $\R$, then $\mathscr FF(t)=e^{t/2}\widehat f(e^t)$. On the Schwartz
core
\[
 \mathscr F^*=\mathscr F^{-1}=\mathscr F,\qquad
 JV_aJ=V_{-a},\qquad X+\mathscr FX\mathscr F=\mathcal M .
\]
The first identity follows from $Ju_\infty(D)J=u_\infty(D)^*$. For the third, if $U=e^{i\theta}(D)$
with $\theta$ real and of polynomial growth with its first two derivatives, then
$U^*XU=X+\theta'(D)$ (differentiate $\widehat{XG}=-i\partial_\tau\widehat G$, or see the time-delay
identity $\langle UG,XUG\rangle=\langle G,XG\rangle+\frac1{2\pi}\int\theta'|\widehat G|^2$); with
$JX=-XJ$ this gives $\mathscr FX\mathscr F=-UXU^*=-X+m(D)$.

Shifts are written with the right shift $V_a$ of \S\ref{subsec:conv}, acting on the whole line; so
$V_{-\log n}F(t)=F(t+\log n)$.

\subsection{The finite sum--product identity}

For integers $M,N\ge1$ define on the Schwartz core
\begin{align*}
 S_N&=\sum_{n\le N}n^{-1/2}V_{-\log n},&
 S'_N&=\sum_{n\le N}(\log n)n^{-1/2}V_{-\log n},\\
 L_M&=\sum_{d\le M}\Lambda(d)d^{-1/2}V_{-\log d},&
 A_M&=\mathcal M-(L_M+L_M^*),\\
 E_N&=S_N\mathscr F-JS_N,&
 H_{M,N}&=L_MS_N-S'_N,
\end{align*}
and $\{X,E\}=XE+EX$. For the full integer comb, Poisson summation formally reads
$S_\infty\mathscr F=JS_\infty$; thus $E_N$ is the finite Poisson defect, and $H_{M,N}$ is the
incomplete divisor hyperbola.

\begin{theorem}[finite Poisson--Chebyshev identity]\label{thm:PCfinite}
On the Schwartz core,
\begin{equation}\label{eq:PCfinite}
 S_NA_M\mathscr F
 =\{X,E_N\}-L_M^*E_N-JH_{M,N}-H_{M,N}\mathscr F.
\end{equation}
If $M\ge N$, then
\[
 H_{M,N}=\sum_{k>N}h_{M,N}(k)k^{-1/2}V_{-\log k},\qquad
 h_{M,N}(k)=\sum_{\substack{dm=k\\d\le M,\ m\le N}}\Lambda(d)\ge0.
\]
In particular $H_{M,N}$ has no coefficient at $k\le N$.
\end{theorem}
\begin{proof}
The commutator $[X,S_N]=-S'_N$ and the three Fourier identities above give
\[
 S_N\mathcal M\mathscr F=JS'_N+S'_N\mathscr F+\{X,E_N\}.
\]
Insert $S'_N=L_MS_N-H_{M,N}$, use $JL_M=L_M^*J$ and
$JS_N=S_N\mathscr F-E_N$, and subtract $S_N(L_M+L_M^*)\mathscr F$.
This proves \eqref{eq:PCfinite}. The coefficient of
$k^{-1/2}V_{-\log k}$ in $L_MS_N$ is the displayed divisor sum. For
$k\le N\le M$ every divisor occurs, so $\sum_{d\mid k}\Lambda(d)=\log k$ cancels the
coefficient of $S'_N$.
\end{proof}

The identity is specific to $\Lambda$: replacing $\Lambda$ by other coefficients $a(d)$ leaves, at
$k\le N$, the coefficient $\sum_{d\mid k}a(d)-\log k$. On the other hand the square of either side
of \eqref{eq:PCfinite} is invariant under $A_M\mapsto-A_M$, so squaring it cannot determine the sign
of the Weil form.

\subsection{A faithful observation by the Gamma poles}

Fix $\delta>0$ and put, for $m\ge1$,
\[
 a_m=2m+\tfrac12,\qquad
 b_m=\sqrt{a_m}\,e^{-a_m(T+\delta)},\qquad
 \mathcal L_F(a)=\int_{-T}^T e^{-at}F(t)\,dt,\qquad
 (\mathcal OF)_m=b_m\mathcal L_F(a_m).
\]
Since $\|e^{-a_mt}\|^2_{L^2(\IT)}\le2Te^{2a_mT}$, the map $\mathcal O:L^2(\IT)\to\ell^2$ is Hilbert--Schmidt,
and
\begin{equation}\label{eq:Gammaenergy}
 \|\mathcal OF\|^2
 =\sum_{m\ge1}a_me^{-2a_m(T+\delta)}|\mathcal L_F(a_m)|^2
 =\iint F(t)\overline{F(u)}\,\bigl[-(k_\Gamma^+)'(2T+2\delta+t+u)\bigr]\,dt\,du,
\end{equation}
where
\[
 k_\Gamma^+(v)=\sum_{m\ge1}e^{-a_mv}.
\]
This is a positive Hankel form, and it is faithful: $\mathcal OF=0$ implies $F=0$. Indeed, after the
substitution $x=e^{-2t}$, vanishing of every $\mathcal L_F(a_m)$ says that $x^{1/4}F(-\frac12\log x)\,x^{-1}$
is orthogonal to every power $x^m$, $m\ge1$, on the compact interval $[e^{-2T},e^{2T}]$. Since this
interval does not contain $0$, the functions $x\,p(x)$, $p$ a polynomial, are dense in its continuous
functions; hence $F=0$.

\begin{remark}[interpretation]\label{rem:residue}
The numbers $\mathcal L_F(a_m)$ are the residues at the positive odd poles of the Gamma factor. Write
$s=\frac12+z$, $\chi(s)=\pi^{s-1/2}\Gamma(\frac{1-s}2)/\Gamma(\frac s2)$, so that
$\zeta(s)=\chi(s)\zeta(1-s)$, and $d_\Gamma=-\chi'/\chi$. Near $s_m=2m+1$ one has
$\psi(\frac{1-s}2)=2/(s-s_m)+O(1)$, hence $\operatorname{Res}_{s=s_m}d_\Gamma=1$. For a
first-contact vector $F$ of parity $\varepsilon$ let $W^0_\varepsilon(s)$ be the corrected product
defect: the function $d_\Gamma(s)\mathcal L_F(s-\frac12)$ minus the Laplace transforms of the two compactly
supported Euler histories of $F$ and minus two polar corrections supported outside $\IT$. We use only
two properties of it: its inverse transform restricts on $\IT$ to the Tate vector $r$ of
\eqref{eq:localweak}, and every term other than $d_\Gamma \mathcal L_F$ is holomorphic at each $s_m$. Hence
\begin{equation}\label{eq:residue}
 \operatorname*{Res}_{s=s_m}\bigl[\zeta(s)W^0_\varepsilon(s)\bigr]
 =\zeta(s_m)\mathcal L_F(a_m).
\end{equation}
Consequently holomorphy of $\zeta W^0_\varepsilon$ at every $s_m$, and vanishing of
\eqref{eq:Gammaenergy}, are each sufficient for $F=0$. Formula \eqref{eq:residue} places the
obstruction outside the window: corrections supported in a compact interval have entire Laplace
transforms and cannot cancel these residues. The defect $W^0_\varepsilon$ is described here only through the two properties stated, and nothing below depends on this remark.
\end{remark}

\subsection{Compressed collars and the balance on the radical}\label{subsec:collars}

For $N\ge2$ put
\[
 B_N=Q_T\sum_{2\le n\le N}\Lambda(n)n^{-1/2}V_{\log n}Q_T,\qquad
 d_{m,N}=\sum_{2\le n\le N}\frac{\Lambda(n)}{n^{a_m+1/2}},\qquad
 D_N=\operatorname{diag}(d_{m,N})_{m\ge1}.
\]
For $N\ge e^{2T}$, $B_N$ contains every shift that acts on $L^2(\IT)$, and
$\mathfrak a_T(F,G)=\mathfrak m_T(F,G)-\langle(B_N+B_N^*)F,G\rangle$, where
\[
 \mathfrak m_T(F,G)=\frac1{2\pi}\int_\R m(\tau)\,\Fh(\tau)\overline{\widehat G(\tau)}\,d\tau
\]
is the compressed Gamma form, i.e.\ the form of $\mathcal M$ on the form domain of $L^2(\IT)$.
Thus \eqref{eq:localweak} is the weak form of $(\mathcal M-B_N-B_N^*)F=r$ on $\IT$.
The exterior collar $\mathcal E_N:L^2(\IT)\to\ell^2$ is
\[
 (\mathcal E_NF)_m=b_m\sum_{2\le n\le N}\Lambda(n)n^{-(a_m+1/2)}
 \int_{\max(-T,T-\log n)}^T e^{-a_mu}F(u)\,du.
\]
The substitution $u=t-\log n$ gives the exact identity
\begin{equation}\label{eq:collar}
 \mathcal OB_N=D_N\mathcal O-\mathcal E_N.
\end{equation}
Primes with $n\ge e^{2T}$ do not contribute to $B_N$ but do contribute to $\mathcal E_N$: once
$N\ge e^{2T}$, one has $\mathcal E_{N'}-\mathcal E_N=(D_{N'}-D_N)\mathcal O$ for $N'\ge N$, so these terms
cannot be discarded.

Let $N\ge e^{2T}$, let $\mathcal O_L,D_{N,L},\mathcal E_{N,L}$ denote the truncations to $1\le m\le L$, and put
$G_L=\mathcal O_L^*\mathcal O_L$ and $\Sigma_L=G_L+JG_LJ$. The operator $\Sigma_L$ commutes with $J$ and its kernel is a finite
sum of exponentials on $\overline{\IT}$, so $\mathrm{pr}_T\Sigma_LF$ lies in the form domain and is an admissible
test. Let $F$ be a first-contact vector with Tate vector $r$. Taking $G=\Sigma_LF$ in
\eqref{eq:localweak}, using $JB_NJ=B_N^*$, equal parity of $F$ and $\Sigma_LF$, and \eqref{eq:collar}, one
obtains
\begin{equation}\label{eq:Qbalance}
 0=2c_{N,L}(F)+Q_{N,L}(F),\qquad
 c_{N,L}(F)=\langle \mathcal O_LF,D_{N,L}\mathcal O_LF\rangle,
\end{equation}
where, with no further boundary term,
\begin{align}\label{eq:Qfull}
 Q_{N,L}(F)={}&2\Rel\langle \mathcal O_LF,\mathcal O_LB_N^*F\rangle
 -2\Rel\langle \mathcal O_LF,\mathcal E_{N,L}F\rangle\notag\\
 &-\Rel\mathfrak m_T(F,\Sigma_LF)+\Rel\langle r,\Sigma_LF\rangle.
\end{align}
The backward term involves the truncated moments
\[
 (\mathcal O_LB_N^*F)_m=b_m\sum_{2\le n<e^{2T}}
 \Lambda(n)n^{a_m-1/2}\int_{-T+\log n}^T e^{-a_mu}F(u)\,du;
\]
replacing them by full moments would remove a collar.

Identity \eqref{eq:Qbalance} carries no sign information. Indeed, for $F$ in the domain of $H_T$
with $r$ defined as above, the same computation gives
\begin{equation}\label{eq:Qrecombine}
 Q_{N,L}=-2c_{N,L}-\Rel\mathfrak a_T(F,\mathrm{pr}_T\Sigma_LF),
\end{equation}
and for an eigenvector $H_TF=\lambda F$ of parity $\varepsilon$, since
$\langle F,\Sigma_LF\rangle=2\|\mathcal O_LF\|^2$,
\begin{equation}\label{eq:Qspectral}
 Q_{N,L}(F)=-2\sum_{m\le L}(d_{m,N}+\lambda)|(\mathcal O_LF)_m|^2.
\end{equation}
Thus $Q_{N,L}$ is negative on every eigenvector with $\lambda\ge0$ and $\mathcal O_LF\ne0$. In particular, the
inequality $Q_{N,L}\ge0$ on the radical is equivalent to $\mathcal O_LF=0$ there, and $Q_{N,L}$ admits no
representation as a sum of squares valid on all of $\PT$.

\subsection{The Gamma--Euler contraction and what it does not prove}

The Euler channel $D_N$ is a strict contraction, uniformly in $m$ and $N$.

\begin{proposition}[strict Gamma--Euler contraction]\label{prop:GEgap}
For all $m\ge1$ and $N\ge2$,
\[
 0<d_{m,N}\le q_0:=\frac{\log2}{4}+\frac1{16}<0.236<1.
\]
Consequently
\begin{equation}\label{eq:gapop}
 \Delta_{N,L}:=\mathcal O_L^*(I-D_{N,L})\mathcal O_L
 \ge(1-q_0)\mathcal O_L^*\mathcal O_L,
\end{equation}
and, keeping the collar,
\begin{equation}\label{eq:gapcollar}
 \langle F,\Delta_{N,L}F\rangle
 =\|\mathcal O_LF\|^2-\Rel\langle \mathcal O_LF,\mathcal O_LB_NF\rangle
  -\Rel\langle \mathcal O_LF,\mathcal E_{N,L}F\rangle.
\end{equation}
For a first-contact vector $F$,
\begin{equation}\label{eq:gapQ}
 2\langle F,\Delta_{N,L}F\rangle=2\|\mathcal O_LF\|^2+Q_{N,L}(F).
\end{equation}
\end{proposition}
\begin{proof}
Since $a_m+\frac12=2m+1\ge3$, $\Lambda(n)\le\log n$, and $x^{-3}\log x$ is decreasing on $[2,\infty)$,
\[
 d_{m,N}\le\sum_{n\ge2}\frac{\log n}{n^3}
 \le\frac{\log2}{8}+\int_2^\infty\frac{\log x}{x^3}\,dx
 =\frac{\log2}{4}+\frac1{16}.
\]
Positivity holds because the term $n=2$ is present. Equation \eqref{eq:gapop} follows coordinatewise,
and \eqref{eq:gapcollar} follows by pairing \eqref{eq:collar} with $\mathcal O_LF$. Finally, on the radical
\eqref{eq:Qbalance} gives $Q_{N,L}=-2\langle \mathcal O_LF,D_{N,L}\mathcal O_LF\rangle$.
\end{proof}

Letting $N\to\infty$, $d_m:=\lim_Nd_{m,N}=-\zeta'(s_m)/\zeta(s_m)$. Hence the residue vector
$x_m=\zeta(s_m)b_m\mathcal L_F(a_m)$ and the Chebyshev jet $j_m=-\zeta'(s_m)b_m\mathcal L_F(a_m)$ satisfy, with
$D=\operatorname{diag}(d_m)$,
\begin{equation}\label{eq:smallgain}
 j=D x,\qquad \|j\|\le q_0\|x\|.
\end{equation}
This is a quantitative identity, but it supplies no independent leverage toward vanishing.  For a
fixed $L$ injective on the radical and $F\in\ker H_T$, the following are equivalent:
\begin{equation}\label{eq:equivgate}
 \mathcal O_LF=0
 \quad\Longleftrightarrow\quad
 \langle F,\Delta_{N,L}F\rangle\le0
 \quad\Longleftrightarrow\quad
 Q_{N,L}(F)\ge0.
\end{equation}
The first equivalence follows from \eqref{eq:gapop}; the second follows from
$Q_{N,L}=-2\langle \mathcal O_LF,D_{N,L}\mathcal O_LF\rangle$ and strict positivity of $D_{N,L}$.
Likewise, because $j=Dx$ with every diagonal entry of $D$ in $(0,q_0]$,
\[
 \|x\|^2\le\Rel\langle x,j\rangle\quad\Longleftrightarrow\quad x=0.
\]
Thus neither the finite return inequality nor its infinite-channel version is an intermediate
estimate: each is the desired conclusion in another notation.  Cauchy--Schwarz and $q_0<1$ merely
verify this equivalence.

\begin{problem}[Gamma-moment annihilation at first contact]\label{prob:annihilate}
Starting from the finite Poisson--Chebyshev identity \eqref{eq:PCfinite}, prove directly that for every primitive first-contact
vector $F$ there is a finite $L$, injective on the radical, for which
\begin{equation}\label{eq:annihilate}
 \mathcal O_LF=0.
\end{equation}
The proof must use all three pieces absent from the formal radical equation: (i) minimality of the
contact window, not merely $H_TF=0$; (ii) an oriented causal or Poisson--Tate identity which changes
under reversal of the Weil sign; and (iii) $\Lambda*1=\log$ in the step which forces vanishing or
fixes the sign, not only in the preliminary contraction estimate.  It must retain both Fourier
orientations, the continuum polar term, the full exterior collars and the crossed Tate endpoint
before taking a limit.
\end{problem}

Condition (i) is a falsification test: by \eqref{eq:Qspectral}, reasoning which also applied to
positive eigenvalues would give the wrong conclusion.  Condition (ii) excludes norm-squaring of
\eqref{eq:PCfinite}, which is invariant under $A_M\mapsto-A_M$.  Condition (iii) retains the genuine
arithmetic advantage of this route: changing one von Mangoldt weight creates an interior divisor
defect, while Poisson and a functional equation alone also occur for Epstein/Davenport--Heilbronn
examples without an Euler product.

A finite $L$ with $\mathcal O_L$ injective on $\ker H_T$ exists at each fixed window. The kernel is
finite-dimensional because $H_T$ has compact resolvent; the subspaces $\ker H_T\cap\ker \mathcal O_L$ decrease
with $L$ and have zero intersection by faithfulness of $\mathcal O$; a decreasing chain of subspaces of a
finite-dimensional space stabilizes, so one of them is $\{0\}$. The argument gives no bound on $L$
that is uniform in $T$.

\section{Open problems}\label{sec:open}

\begin{enumerate}
\item \emph{Window positivity.} Prove \eqref{eq:target} for $T>\frac12\log5$. By Theorem~\ref{thm:coer-schur}
this is equivalent, at each window and for any $k\ge k_0(T)$, to $S_{T,k}\succeq0$; by
Lemma~\ref{lem:firstcontact} it is equivalent to the absence of a first-contact window. A first test case is
the window $T=\frac12\log5+\epsilon$ at which the contact $5$ enters.
\item \emph{Gamma-moment annihilation at first contact} (Problem~\ref{prob:annihilate}). The return
inequalities are equivalent to its conclusion by \eqref{eq:equivgate}; they are not intermediate estimates.
\item \emph{Class-dependent value bounds} (Problem~\ref{prob:A}): an explicit deficit for lattice
certificates, first for one contact and then for two, and the rate of convergence of the ceilings of
Theorem~\ref{thm:R}.
\item \emph{The polar capacity beyond $\frac12\log5$.} Whether $\mathfrak G_T\succ0$ at some $T>\frac12\log5$;
by Proposition~\ref{prop:onewindow} a positive answer at one window gives it at all smaller ones, and by
Theorem~\ref{thm:parityschur} it is the complete-form counterpart of item 1.
\item \emph{A certified lower bound for $S_{T,k}$} at a window where $\lamA$ is not already certified
(Remark~\ref{rem:coer-honest}(2)).
\item Whether the optimal multipliers of the primitive problem converge to $\beta=(2,0)$, i.e.\ whether
$\lamP=\lamW$ at the windows considered (Remark~\ref{rem:Kprim}).
\end{enumerate}

\appendix

\section{Computer verifications}\label{app:encl}

This appendix describes the computer verifications used in Theorems~\ref{thm:C}, \ref{thm:N},
\ref{thm:Np}, \ref{thm:index}, \ref{thm:log5A}, \ref{thm:log5W}, Proposition~\ref{prop:archneg} and
Corollary~\ref{cor:convfail}. All of them are carried out in outward-rounded interval arithmetic
(\texttt{mpmath}), except those of Theorems~\ref{thm:log5A} and \ref{thm:log5W}, which use ball arithmetic
(\S\ref{subsec:encllog5}). The data (atoms, coefficient vectors, test kernels and interval endpoints) and the
scripts accompany the paper; each verification runs in at most a few minutes on one core with Python~3,
\texttt{mpmath}~1.3.0 and \texttt{scipy}~1.11.4, and the ball-arithmetic certificates use
\texttt{python-flint}~0.9.0 (Arb) and take about seven minutes per parity.

\subsection{The time--band limiting operator}
\begin{lemma}[a Ky Fan ceiling for the time--band limiting operator]\label{lem:KyFan}
Let $S_c=Q_TB_sQ_T$ on $L^2(\IT)$, where $Q_T$ is restriction to $\IT$ and $B_s$ the band projection
to $|\tau|<s$, and put $c=Ts$. Let $j_n$ be the spherical Bessel functions and
$G_k(x)=\sum_{j<k}(2j+1)j_j(x)^2\in[0,1]$. Then, in the orthonormal Legendre basis
$e_j(t)=\sqrt{(2j+1)/2T}\,P_j(t/T)$,
\[
\langle S_ce_j,e_j\rangle=\frac{2(2j+1)}\pi\int_0^cj_j(x)^2dx,\qquad
\sum_{j\ge0}\langle S_ce_j,e_j\rangle=\frac{2c}\pi=\operatorname{tr}S_c ,
\]
and for every $k\ge0$
\[
\lambda_k(S_c)\ \le\ \Lambda_k(c):=\frac2\pi\int_0^c\big(1-G_k(x)\big)\,dx .
\]
In particular $\Lambda_0(c)=2c/\pi$ is the trace and $\Lambda_k(c)=O(c^{2k+1})$ as $c\to0$.
\end{lemma}
\begin{proof}
From $\int_{-1}^1P_j(u)e^{-ixu}du=2(-i)^jj_j(x)$ one gets
$\widehat{e_j}(\tau)=\sqrt{(2j+1)/2T}\,\cdot2T(-i)^jj_j(\tau T)$, so
$|\widehat{e_j}(\tau)|^2=2T(2j+1)j_j(\tau T)^2$ and Plancherel gives the displayed diagonal entry.
The addition theorem $\sum_{j\ge0}(2j+1)j_j(x)^2=1$ gives the trace identity, so the basis loses no
trace. By the Ky Fan maximum principle $\sum_{j<k}\lambda_j(S_c)\ge\sum_{j<k}\langle
S_ce_j,e_j\rangle$; since $S_c=(B_sQ_T)^*(B_sQ_T)\ge0$ its eigenvalues are nonnegative and
non-increasing, whence
$\lambda_k\le\sum_{n\ge k}\lambda_n=\operatorname{tr}S_c-\sum_{j<k}\lambda_j\le\sum_{j\ge k}\langle
S_ce_j,e_j\rangle=\Lambda_k(c)$.
\end{proof}

\begin{corollary}[quadrature-free form]\label{cor:P1}
$|j_n(x)|\le x^n/(2n+1)!!$ for $x\ge0$, hence
$\lambda_k(S_c)\le\Lambda^*_k(c):=\frac2\pi\sum_{j\ge k}c^{2j+1}/\big((2j+1)!!\big)^2$, a convergent
series of rationals in $c$: the bound is evaluated by arithmetic alone, with no quadrature.
\end{corollary}


\subsection{The interval enclosures of Theorems~\ref{thm:N} and \ref{thm:Np}}
\emph{The window $T=\frac12\log3$.} All three integrals are computed in interval arithmetic at $25$ digits
(\texttt{mpmath}), the only non-interval inputs being the rational $\tau_0$ and the decimal $q$, which
are data of the bound and not approximations. The archimedean multiplier is summed from
Theorem~\ref{thm:S}(i), $m(\tau)=C_0+\sum_{k\ge0}\tau^2/(a_k(\tau^2+4a_k^2))$, with
$C_0=\psi(\frac14)-\log\pi=-\gamma-\frac\pi2-3\log2-\log\pi$ by the Gauss digamma theorem, the first
$100$ terms in interval arithmetic and the tail bracketed by the convex midpoint rule: with $A=K-\frac14$
and $t(a)=\tau^2/(\tau^2a+4a^3)$, the function $t$ is positive, decreasing and convex on $(0,\infty)$
(because $2u'^2-u''u=2\tau^4+24\tau^2a^2+192a^4>0$ for $u=\tau^2a+4a^3$), the points $a_k=k+\frac14$
are the midpoints of unit cells tiling $[A,\infty)$, so
$\sum_{k\ge K}t(a_k)\in[\Lambda-E,\Lambda]$ with $\Lambda=\frac12\log(1+\tau^2/4A^2)=\int_A^\infty t$
and $E=\frac1{24}|t'(A)|=\frac1{24}\tau^2(\tau^2+12A^2)/(\tau^2A+4A^3)^2$, at most $2.4\cdot10^{-7}$
on the range used. The multipliers are evaluated from the closed form
$\int_0^{2T}e^{su}\cos\tau u\,du=[e^{2Ts}(s\cos2T\tau+\tau\sin2T\tau)-s]/(s^2+\tau^2)$ at
$s=\pm\frac12$ and its $s$-derivative, without quadrature. The integrals over $W$ are composite
Simpson with the classical remainder $Ph^4C_4/180$, the fourth derivatives being bounded by
$|m^{(j)}|\le2^{-j}j!\,\zeta(j+1,\frac14)$, $|(v_2\cos(\tau\log2))^{(j)}|\le v_2(\log2)^j$,
$|\gh0^{(j)}|\le2\cosh T(2T)^{j+1}/(j+1)$, $|\gh1^{(j)}|\le2\sinh T(2T)^{j+2}/(j+2)$ and the Bernstein
bound $|\varphi^{(j)}|\le(d(M-1))^j(\sum_k|q_k|)^2$, combined by the Leibniz rule; the resolutions are
chosen so that the remainder is $7\cdot10^{-5}$ for $a$, which needs only $10^{-4}$, and below
$10^{-8}$ for $b_0,b_1$, which are multiplied by $B$. The denominator $\int_W\varphi=P\|q\|^2$ is
exact. An independent evaluation by adaptive quadrature and a library digamma, outside the interval
arithmetic, lands inside all three enclosures. Steps $d$ with $2T/d\notin\N$ are not covered by
Theorem~\ref{thm:N}: each such $d$ needs its own enclosure with $M'=\lceil2T/d\rceil$.

\emph{The windows $T=0.68$ and $T=\log2$.} The archimedean multiplier, the constant $C_0$, the closed forms of
$\gh0,\gh1$ and the exact denominator $P\|q\|^2$ are as in Theorem~\ref{thm:N}; nonnegativity of
$\varphi$ is structural, since it is evaluated as a sum of two squares, so no property of the
particular $q$ is used. Two things change. First, Simpson's rule is replaced by composite
Gauss--Legendre of order $p$ on panels of width $H$, with the classical remainder
$|E_{\mathrm{panel}}|\le H^{2p+1}(p!)^4\big/\big[(2p+1)((2p)!)^3\big]\cdot\sup|f^{(2p)}|$ and the
derivative bounds of Theorem~\ref{thm:N} extended from order $4$ to order $2p$, the archimedean one
being $|m^{(j)}|\le2^{-j}j!\,\zeta(j+1,\frac14)\le2^{-j}j!\,(4^{j+1}+\zeta(j+1))$ for $j\ge1$ and
$|m|\le\max(|C_0|,m(\tau_{\max}))$ for $j=0$; the nodes are Newton-refined at $60$ digits, widened
before entering interval arithmetic, and guarded by the archived residuals
$\max_i|P_p(x_i)|$ and $|\sum_iw_i-2|$, both below $10^{-60}$ in every run; this guard ensures that the
computed rule is the one to which the remainder formula applies. Second, $q$ is optimised: rather than the bottom eigenvector of
$C_M(\tau_0,\beta^\star)$, which makes $a$ small and leaves $b$ where it falls, one minimises
$a+B(|b_0|+|b_1|)$ over the span of the three or four lowest eigenvectors, which drives
$|b_0|+|b_1|$ to the $10^{-16}$--$10^{-18}$ level while $a$ stays within a fraction of a percent of
the double-precision ceiling. Consequently the $\beta$-allowance $(|b_0|+|b_1|)B$ is below
$1.3\cdot10^{-11}$ in every row, four orders below the smallest $|a|$, and the binding error is the
quadrature on $a$, between $0\%$ and $6.7\%$ of $|a|$. The bound survives any $B\le4.7\cdot10^9$,
a factor $2.5\cdot10^5$ above the constants of Step~1, so only the qualitative content of
Theorem~\ref{thm:Bbeta} --- that some finite a priori box exists --- is used. At two values of $M$
the minimax over $\tau_0$ has two active period placements, and no single pair $(\tau_0,q)$ can then
drive $b$ to zero; a scan over the best period placements supplies one that closes, at the cost of
recovering only a few percent of the double-precision ceiling there. As in Theorem~\ref{thm:N},
steps $d$ with $2T/d\notin\N$ are not covered: each needs its own enclosure with
$M'=\lceil2T/d\rceil$.

\subsection{Verification on the bounded part}

\begin{proposition}[second-order cover]\label{prop:TAY}
Write $f=m+q+r$ with $q=-\sum_{n\in\NT}v_n\cos(\tau\log n)+\nuh$ ($\nu$ atomic) and
$r=\beta_0\gh0+\beta_1\gh1$, and put
$M_2=\sum_{n\in\NT}v_n\log^2n+2\sum_j|a_j|u_j^2+|\beta_0|H_0+|\beta_1|H_1$ with
$H_0=\int_{-2T}^{2T}u^2\cosh\frac u2\,du$, $H_1=\int_{-2T}^{2T}|u|^3\sinh\frac{|u|}2\,du$. Then
$\sup|q''+r''|\le M_2$, and on every cell $[x_0-\frac w2,x_0+\frac w2]\subset[0,\infty)$
\[
\min_{\text{cell}}f\ \ge\ m\!\big(x_0-\tfrac w2\big)+(q+r)(x_0)-|(q+r)'(x_0)|\,\tfrac w2-M_2\tfrac{w^2}8 .
\]
The same bound applies to each row $\nuh+W_k$ of Lemmas~\ref{lem:F1}, \ref{lem:hybrid} with
$m(x+kP)$ and $-|\beta_i|k_i(x+kP)$ evaluated at the left endpoint. A cell is certified when the
right side exceeds the target and is bisected otherwise; near a contact of $f$ the admissible width
is $\approx\sqrt{8c/M_2}$, so the number of evaluations grows like $c^{-1/2}$ rather than the
$c^{-1}$ of a Lipschitz cover.
\end{proposition}
\begin{proof}
Taylor's formula with integral remainder for $q+r\in C^2$, and monotonicity of $m$ and of
$-|\beta_i|k_i$.
\end{proof}


\subsection{The certificates of Theorem~\ref{thm:index}}\label{subsec:enclindex}
In this subsection $\mathrm{par}\in\{\mathrm{ev},\mathrm{odd}\}$, $\chi_{\mathrm{ev}}(t)=\cosh\frac t2$ and
$\chi_{\mathrm{odd}}(t)=\sinh\frac t2$ on $\IT$, and $L^2_{\mathrm{par}}(\IT)$ is the space of real functions
of that parity, so that a real $F\in L^2_{\mathrm{par}}(\IT)$ lies in $\PT$ if and only if
$\langle F,\chi_{\mathrm{par}}\rangle=0$ (\S\ref{subsec:arch}). The basis $e_j$, the operators $S_c$ and the
tails $\Lambda_k$ are those of Lemma~\ref{lem:KyFan}. Every other letter is local to this subsection.

\begin{lemma}[bordered inertia]\label{lem:border}
Let $q$ be a real quadratic form on a real vector space $V$, $f$ a linear functional on
$V$ and $v_0\in V$ with $f(v_0)=1$ and $q(\cdot,v_0)$ finite. On $V\oplus\R$ put
$b(v,\xi)=q[v]+2\xi f(v)$. Then $n_-(b)=n_-(q|_{\ker f})+1$, $n_-$ being the largest dimension of a
subspace on which the form is negative definite.
\end{lemma}
\begin{proof}
Write $v=w+tv_0$ with $w\in\ker f$. Then $b=q[w]+2tq(w,v_0)+t^2q[v_0]+2\xi t$, and the invertible
substitution $\eta=\xi+q(w,v_0)$ gives $b=q[w]+(t^2q[v_0]+2t\eta)$: a $b$-orthogonal sum of
$q|_{\ker f}$ and a binary form of determinant $-1$, whose inertia is $(1,1)$.
\end{proof}

\begin{lemma}[multi-band minorant]\label{lem:multiband}
Let $0=\tau^{(0)}<\dots<\tau^{(N)}=\tau_N$, let $m^-$ be a lower bound for $m$ at the edges,
and let $M_a\ge\sup m''$ on the band $[\tau^{(a)},\tau^{(a+1)}]$, of width $\Delta_a$. Define the even
function $\underline m$ by $\underline m=C_0$ on $[0,\tau^{(1)}]$; on the band $a\ge1$, the chord through
$(\tau^{(a)},m^-(\tau^{(a)}))$ and $(\tau^{(a+1)},m^-(\tau^{(a+1)}))$ lowered by $\max(M_a,0)\Delta_a^2/8$;
and $\underline m=\underline m_\infty:=m^-(\tau_N)$ for $|\tau|\ge\tau_N$. Then $\underline m\le m$ on
$\R$, and $\mT\ge\underline{\mathfrak m}$ as forms, where
$\underline{\mathfrak m}(F)=\frac1{2\pi}\int\underline m|\Fh|^2$. If $m^-$ increases along the edges, then also
$\underline m\le\underline m_\infty$.
\end{lemma}
\begin{proof}
On $[0,\tau^{(1)}]$ and beyond $\tau_N$ use that $m$ is increasing with $m(0)=C_0$
(Lemma~\ref{lem:mbounds}(i)). On a band $a\ge1$ let $L_a$ be the chord through the exact values and
$h=m-L_a$: $h$ vanishes at the ends and $h''\le M_a$, so $h\ge0$ if $M_a\le0$, and otherwise
$h+\frac12M_a(\tau-\tau^{(a)})(\tau^{(a+1)}-\tau)$ is concave and vanishes at the ends, whence
$h\ge-M_a\Delta_a^2/8$; lowering the end values to $m^-$ lowers the chord. Every end value of a
lowered chord is at most $m^-$ at that edge, hence at most $m^-(\tau_N)$ when $m^-$ increases along the
edges, and the chords are affine. Finally $\mT-\underline{\mathfrak m}=\frac1{2\pi}\int(m-\underline m)|\Fh|^2\ge0$.
\end{proof}
In the verification, $m^-$ is the lower end of an interval enclosure of $C_0+\sum_ks/(a_k(a_k^2+s))$, $s=\tau^2/4$
(Theorem~\ref{thm:S}(i)), summed to $3000$ terms with the tail of the decreasing summand bracketed by the
integral test; its width is at most $7\cdot10^{-9}$, while consecutive edge values of $m$ differ by at
least $0.03$, so $m^-$ increases along the edges. The bound $M_a$ sums, over $k<3000$, the termwise
majorant $a_k(a_k^2-2s)/(2(a_k^2+s)^3)$ of the summand $a_k(a_k^2-3s)/(2(a_k^2+s)^3)$ of $m''$, which is
decreasing-then-increasing in $s$ and is therefore bounded by its values at the ends of an $s$-interval
containing the band, plus the positive tail $\sum_{k\ge K_0}(2a_k^3)^{-1}\le(4(K_0-\frac34)^2)^{-1}$,
$K_0=3000$. There are $28$ bands: band $0$ is $[0,0.15]$, twenty fixed edges reach $8$, and eight
geometric edges run from $9$ to $\tau_N=c_*/T$, $c_*=16$ (to six decimals); the largest offset is $0.05791$.

\begin{proposition}[the upper certificate]\label{prop:haynsworth}
Fix $\mathrm{par}$, integers $k_1<k_2$ ($22$ and $36$ in the verification) and $\underline m$ as in
Lemma~\ref{lem:multiband}, with $\underline m\le\underline m_\infty$, and put $c_*=T\tau_N$,
$\omega=\underline m_\infty-\min\underline m$ and $\mathcal M=Q_T(\underline m_\infty-\underline m)(D)Q_T$,
so that $\underline{\mathfrak m}(F)=\underline m_\infty\|F\|^2-\langle\mathcal MF,F\rangle$ and
$0\le\mathcal M\le\omega S_{c_*}$. Let $E$ be the span of the $e_j$, $j<k_1$, of parity $\mathrm{par}$,
$Y$ its orthogonal complement in $L^2_{\mathrm{par}}(\IT)$, $P_Y$ the orthogonal projection onto $Y$,
and $\mathfrak b(F,\xi)=\underline{\mathfrak m}(F)+2\xi\langle F,\chi_{\mathrm{par}}\rangle$ on
$L^2_{\mathrm{par}}(\IT)\oplus\R$.
\emph{(a)} On $Y$, $\mathfrak b\ge q_Y:=\underline m_\infty-\omega\Lambda_{k_1}(c_*)$.
\emph{(b)} Suppose $q_Y>0$, and let $\mathsf R$ be the symmetric matrix on $E\oplus\R$ with entries
\[
\mathsf R_{ij}=\underline{\mathfrak m}(e_i,e_j)-\frac{\langle P_Y\mathcal Me_i,P_Y\mathcal Me_j\rangle}{q_Y},\quad
\mathsf R_{i\xi}=\langle e_i,\chi_{\mathrm{par}}\rangle+\frac{\langle P_Y\mathcal Me_i,P_Y\chi_{\mathrm{par}}\rangle}{q_Y},\quad
\mathsf R_{\xi\xi}=-\frac{\|P_Y\chi_{\mathrm{par}}\|^2}{q_Y}.
\]
Then the Schur complement $\mathfrak s$ of $\mathfrak b|_Y$ in $\mathfrak b$ satisfies $\mathfrak s\ge\mathsf R$
and $n_-(\mathfrak b)=n_-(\mathfrak s)$.
\emph{(c)} If $\mathsf R$ is positive definite on a subspace of codimension $r+1$, then
$n_{\mathrm{par}}(T)\le r$.
\end{proposition}
\begin{proof}
(a) $\mathfrak b|_Y=\underline m_\infty-P_Y\mathcal MP_Y$, and
$\|P_Y\mathcal MP_Y\|\le\omega\|P_YS_{c_*}P_Y\|\le\omega\operatorname{tr}(P^\perp S_{c_*}P^\perp)
=\omega\sum_{j\ge k_1}\langle S_{c_*}e_j,e_j\rangle=\omega\Lambda_{k_1}(c_*)$, where $P^\perp$ removes every
$e_j$, $j<k_1$, of either parity: $Y$ lies in its range, the norm of a positive operator is at most its
trace, and the last equality is the trace identity of Lemma~\ref{lem:KyFan}. (b) For $i<k_1$ and $y\in Y$,
$\mathfrak b(e_i,y)=\underline{\mathfrak m}(e_i,y)=-\langle P_Y\mathcal Me_i,y\rangle$, since $\langle e_i,y\rangle=0$,
and the border pairs with $y$ through $+\langle P_Y\chi_{\mathrm{par}},y\rangle$. The congruence
$(x,y)\mapsto(x,y+(\mathfrak b|_Y)^{-1}\mathfrak b_{YX}x)$ gives $\mathfrak b\cong\mathfrak s\oplus\mathfrak b|_Y$
with $\mathfrak b|_Y>0$, so $n_-(\mathfrak b)=n_-(\mathfrak s)$; and
$\mathfrak b_{XY}(\mathfrak b|_Y)^{-1}\mathfrak b_{YX}$ is at most $q_Y^{-1}$ times the Gram matrix of the
$Y$-columns $-P_Y\mathcal Me_i$ and $+P_Y\chi_{\mathrm{par}}$, which is $\mathfrak s\ge\mathsf R$.
(c) Since $\mT\ge\underline{\mathfrak m}$ and $\PT^{\mathrm{par}}=\ker\langle\cdot,\chi_{\mathrm{par}}\rangle$,
Lemma~\ref{lem:border} gives
$n_{\mathrm{par}}(T)\le n_-(\underline{\mathfrak m}|_{\PT^{\mathrm{par}}})=n_-(\mathfrak b)-1=n_-(\mathfrak s)-1\le n_-(\mathsf R)-1\le r$.
\end{proof}
\emph{The computation.} For $i-l$ even, $\widehat{e_i}\overline{\widehat{e_l}}=2T\sqrt{(2i+1)(2l+1)}\,
(-1)^{(l-i)/2}j_ij_l(\tau T)$ (proof of Lemma~\ref{lem:KyFan}), so with $\underline m=\alpha_a+\beta_a|\tau|$
on the band $a$ and $I_p^{(il)}(c)=\int_0^cx^pj_ij_l\,dx$,
\[
\underline{\mathfrak m}(e_i,e_l)=\underline m_\infty\delta_{il}+\frac2\pi\sqrt{(2i+1)(2l+1)}\,(-1)^{\frac{l-i}2}
\Big[\sum_a\big(\alpha_a\,\Delta I^{(il)}_0+\tfrac{\beta_a}T\,\Delta I^{(il)}_1\big)-\underline m_\infty I^{(il)}_0(c_*)\Big],
\]
$\Delta$ denoting the increment over the band, and the entries with $i-l$ odd vanish. The moments
$I_p$ come from the power series of $j_ij_l$ with exact rational coefficients, the alternating tail being
bounded by its first omitted term once the index exceeds $2c$; $\langle e_l,\chi_{\mathrm{par}}\rangle$ comes
from the positive series of the modified spherical Bessel function $i_l(T/2)$ (valid for $T<1$), and
$\|\chi_{\mathrm{par}}\|^2=\sinh T\pm T$. The inner products over $Y$ are exact sums over $k_1\le l<k_2$ of
the given parity ($\mathcal M$ preserves parity) plus boxes from
$\|P^\perp_{k_2}\mathcal Me_i\|^2\le\rho\,\mathcal M_{ii}$, $\rho=\omega\Lambda_{k_2}(c_*)$, and
Cauchy--Schwarz. In the verification $\Lambda_{22}(16)\le7.2151\cdot10^{-6}$ and $\rho\le4.02\cdot10^{-20}$, so
$q_Y\ge\underline m_\infty-5\cdot10^{-5}>0$. For (c), the eigenvectors of the midpoint of $\mathsf R$,
read as decimal data, give a matrix $U_r$ omitting the $r+1$ lowest; interval Cholesky of
$U_r^{\mathsf T}\mathsf RU_r-\mu I$ over the whole interval box of $\mathsf R$ succeeds for a certified
$\mu>0$, which forces $U_r$ to be injective and $\mathsf R$ to be positive definite on its range.
Everything is interval arithmetic at $40$ digits, with no quadrature; the entry widths of $\mathsf R$ are
at most $1.7\cdot10^{-13}$. The first eigenvalue of $\mathsf R$ is always the hyperbolic direction of
the border (Lemma~\ref{lem:border}); $r$ counts the others. The verified cases are:
\begin{center}\small
\begin{tabular}{@{}lcccc@{}}
\toprule
$T$ & sector & $\underline m_\infty$ & certified & margin $\mu$\\
\midrule
$0.58$ & even & $1.479384$ & $n_{\mathrm{ev}}\le0$ & $0.015998$\\
$0.58$ & odd & $1.479384$ & $n_{\mathrm{odd}}\le0$ & $0.047516$\\
$\log2$ & odd & $1.301146$ & $n_{\mathrm{odd}}\le0$ & $0.041093$\\
$0.90$ & even & $1.039940$ & $n_{\mathrm{ev}}\le1$ & $0.185541$\\
$0.90$ & odd & $1.039940$ & $n_{\mathrm{odd}}\le1$ & $0.144840$\\
\bottomrule
\end{tabular}
\end{center}
The cases $\log2$ (even, $\le1$, $\mu=0.420225$) and $0.80$ (even and odd, $\le1$,
$\mu=0.298000$, $0.161624$) were also verified; they follow from the table by
Lemma~\ref{lem:zeroext}.

\begin{proposition}[an explicit negative plane]\label{prop:negplane}
Let $T<\operatorname{arccosh}\sqrt2=0.88137$ and let $F,H\in L^2_{\mathrm{par}}(\IT)$ be
explicit real polynomials, given by their Legendre coefficients, such that $\langle\cdot,\chi_{\mathrm{par}}\rangle$
does not vanish on $\mathrm{span}\{F,H\}$. If the matrix
\[
B_3=\begin{pmatrix}\mT[F]&\mT(F,H)&\langle F,\chi_{\mathrm{par}}\rangle\\
\mT(F,H)&\mT[H]&\langle H,\chi_{\mathrm{par}}\rangle\\
\langle F,\chi_{\mathrm{par}}\rangle&\langle H,\chi_{\mathrm{par}}\rangle&0\end{pmatrix}
\]
is negative definite on a two-dimensional subspace, then $\mathrm{span}\{F,H\}\cap\PT^{\mathrm{par}}$
contains an $F_0\ne0$ with $\mT[F_0]<0$, and $n_{\mathrm{par}}(T)\ge1$.
\end{proposition}
\begin{proof}
$B_3$ is the bordered form of $\mT|_{\mathrm{span}\{F,H\}}$ with border $\langle\cdot,\chi_{\mathrm{par}}\rangle$,
so Lemma~\ref{lem:border} in dimension two gives
$n_-(\mT|_{\mathrm{span}\{F,H\}\cap\PT})=n_-(B_3)-1\ge1$.
\end{proof}
In the verification $F$ is of degree $<18$, $H$ is the Legendre truncation of
$\chi_{\mathrm{par}}$, and $\langle H,\chi_{\mathrm{par}}\rangle\in[0.072,1.69]$. The quadratic values
$\mT[G]=C_0c_G(0)+\int_0^{2T}(c_G(0)-c_G)K+c_G(0)\int_{2T}^\infty K$ (Theorem~\ref{thm:A} with
$\NT=\emptyset$), for $G=F,H,F\pm H$, are enclosed as follows. (a) $c_G$ is computed as an exact polynomial with interval
coefficients, by rational arithmetic on the Legendre recurrence, with no quadrature. (b) $g:=c_G(0)-c_G\ge0$ on
$[0,2T]$ by Cauchy--Schwarz, which allows integrating an enclosure of $K$ against $g$ term by term. (c) On
$[\delta,2T]$, $K$ is squeezed between two affine functions on each panel, since $1/(2\sinh(u/2))$ is convex
and $1/(2\cosh(u/2))$ is concave for $u/2<\operatorname{arccosh}\sqrt2$; this is the range condition
$T<0.88137$; $\int g$ and $\int ug$ are evaluated exactly from the interval antiderivatives. (d) On
$[0,\delta]$, $\delta=10^{-6}$, use $K(u)\le u^{-1}+\frac12$ and $g(u)\le G_1u$ with $G_1$ from the exact
Chebyshev expansion of $g$. (e) The remaining terms are closed forms in intervals:
$C_0=-\gamma-3\log2-\frac\pi2-\log\pi$ and $\int_{2T}^\infty K=\log\coth\frac T2+\pi-2\arctan e^{T}$. The
off-diagonal entry is obtained by polarisation, $\mT(F,H)=\frac14(\mT[F+H]-\mT[F-H])$, and the
pairings with $\chi_{\mathrm{par}}$ from exact moment series.

\begin{lemma}[rounding in the polarisation]\label{lem:polround}
If $F\pm H$ are formed in double precision, the value of $\mT(F,H)$ computed from
the rounded combinations differs from the true one by at most $10^{-10}$, for $0.6\le T\le0.9$ and the
vectors used.
\end{lemma}
\begin{proof}
The rounded combinations are $F\pm H+\delta_\pm$ with Legendre coefficients $|\delta_{\pm,k}|\le3\cdot10^{-16}$
(all coefficients of $F\pm H$ are at most $2.2$ in modulus), so on $\IT$ one has
$\|\delta\|_\infty\le3\cdot10^{-14}$ and $\|\delta'\|_\infty\le4\cdot10^{-12}$ for degree $<18$. The
Dirichlet form of Proposition~\ref{prop:nodes}(i) with $K(u)\le u^{-1}+\frac12$ gives
$\Lev(c_\delta)\le\|\delta'\|_\infty^2T^2(2T+2T^2)+\frac12\|\delta\|_\infty^2\int_{\IT}\kappa\le10^{-22}$,
and then
$|\mT[v+\delta]-\mT[v]|\le|C_0|(2\|v\|\|\delta\|+\|\delta\|^2)+2\sqrt{\Lev(c_v)\Lev(c_\delta)}+\Lev(c_\delta)\le2\cdot10^{-10}$
with $\Lev(c_v)\le50$; divide by four.
\end{proof}
Two further effects are covered by the same inflation: the stored interval ends are rounded to double
precision, and the value of Euler's constant entering $C_0$ is correct only to about $10^{-51}$. The
verification therefore rebuilds $B_3$ from the stored enclosures, inflates every entry by $\pm10^{-9}$, certifies negative definiteness on
a two-plane over the inflated box by interval Cholesky with a certified margin, and re-imposes
$T<0.88137$. It certifies $n_{\mathrm{ev}}\ge1$ at $T=0.60$ (margin $0.006627$) and $n_{\mathrm{odd}}\ge1$
at $T=0.81$ (margin $0.005124$), which is what Theorem~\ref{thm:index} uses, and also $n_{\mathrm{ev}}\ge1$
at $\log2$ and $0.80$ and $n_{\mathrm{odd}}\ge1$ at $0.82$ (margins $0.147170$, $0.290056$, $0.012368$),
which follow from the first two by Lemma~\ref{lem:zeroext}. A run at $T=0.90$ lies outside the range


\begin{lemma}[a certificate with $r=0$ is strict]\label{lem:rzero}
If Proposition~\ref{prop:haynsworth}(c) holds with $r=0$ at $T$, then
$\lambda_{\mathrm{par}}(T)>0$.
\end{lemma}
\begin{proof}
Then $\underline{\mathfrak m}\ge0$ on $\PT^{\mathrm{par}}$. At the attained minimiser $F^*$ of
Lemma~\ref{lem:strict}, $\mT[F^*]-\underline{\mathfrak m}[F^*]=\frac1{2\pi}\int(m-\underline m)|\Fh^*|^2>0$,
because $m-\underline m=m-C_0=\varphi(\tau^2)>0$ on $0<|\tau|<\tau^{(1)}$ (Theorem~\ref{thm:S}(i)) and
$\Fh^*$ is entire and not identically zero.
\end{proof}


\subsection{The certificates at \texorpdfstring{$T=\frac12\log5$}{T = (1/2) log 5}}\label{subsec:encllog5}
The notation $\chi_{\mathrm{par}}$, $e_j$, $S_c$, $\Lambda_k$ is that of \S\ref{subsec:enclindex}, and
$\ell_{\mathrm{par}}$, $\varepsilon_{\mathrm{par}}$ are as in \S\ref{sec:complete}. We work with real functions,
which suffices as in \S\ref{sec:complete}. Unlike Proposition~\ref{prop:haynsworth}, the primes are kept:
the lower form is a frequency truncation of the full symbol $\sT$, and the operator subtracted from a
multiple of the identity is indefinite.

\emph{The lower form.} Fix $\tau_N>0$ and a number $s_\infty\le m(\tau_N)-\Aconst$. Since $m$ is increasing
with $m(0)=C_0$ and the comb is at least $-\Aconst$ (Lemma~\ref{lem:mbounds}(i),(iii)), the even function
$\underline s$ equal to $\sT$ on $|\tau|\le\tau_N$ and to $s_\infty$ on $|\tau|>\tau_N$ satisfies
$\underline s\le\sT$. Write $g=s_\infty-\underline s=(s_\infty-\sT)\mathbf 1_{|\tau|<\tau_N}=g_+-g_-$ with
$g_\pm\ge0$, and $\mathcal G_{(\pm)}=Q_Tg_{(\pm)}(D)Q_T$. Then
\[
\mathfrak a_T(F)\ \ge\ s_\infty\|F\|^2-\langle\mathcal GF,F\rangle,\qquad
0\le\mathcal G_\pm\le\omega_\pm S_{c_*},\quad c_*=T\tau_N,
\]
with $\omega_+=s_\infty-C_0+\Aconst$ and $\omega_-=m(\tau_N)+\Aconst-s_\infty$. For $\mu\in\R$ and
$c\in\{0,\frac12,-\frac12\}$ put, on $L^2_{\mathrm{par}}(\IT)\oplus\R$,
\[
\mathfrak b_{\mu,c}(F,\xi)=(s_\infty-\mu)\|F\|^2-\langle\mathcal GF,F\rangle+2\xi\,\ell_{\mathrm{par}}(F)+c\,\xi^2 .
\]

\begin{proposition}[the lower certificate]\label{prop:lowercert}
Let $k_1<k_2$, let $E$ be the span of the $e_j$, $j<k_1$, of parity $\mathrm{par}$, $Y$ its orthogonal
complement in $L^2_{\mathrm{par}}(\IT)$, and $P_Y$ the orthogonal projection onto $Y$.
\emph{(a)} On $Y$, $\mathfrak b_{\mu,c}\ge q_Y:=s_\infty-\mu-\omega_+\Lambda_{k_1}(c_*)$.
\emph{(b)} Suppose $q_Y>0$, and let $\mathsf R$ be the symmetric matrix on $E\oplus\R$ with
\[
\mathsf R_{ij}=(s_\infty-\mu)\delta_{ij}-\langle\mathcal Ge_i,e_j\rangle
-\frac{\langle P_Y\mathcal Ge_i,P_Y\mathcal Ge_j\rangle}{q_Y},\quad
\mathsf R_{i\xi}=\langle e_i,\chi_{\mathrm{par}}\rangle+\frac{\langle P_Y\mathcal Ge_i,P_Y\chi_{\mathrm{par}}\rangle}{q_Y},\quad
\mathsf R_{\xi\xi}=c-\frac{\|P_Y\chi_{\mathrm{par}}\|^2}{q_Y}.
\]
Then $n_-(\mathfrak b_{\mu,c})\le n_-(\mathsf R)$, and the same holds for any symmetric $\mathsf R'\le\mathsf R$.
\emph{(c)} If $n_-(\mathsf R')\le1$ with $c=0$, then $\mathfrak a_T\ge\mu$ on the unit sphere of
$\PT^{\mathrm{par}}$. If $n_-(\mathsf R')=0$ with $c=\frac12$ and $\mathrm{par}=\mathrm{odd}$, or $n_-(\mathsf R')\le1$ with
$c=-\frac12$ and $\mathrm{par}=\mathrm{ev}$, then $W_T\ge\mu\|\cdot\|^2$ on $\mathcal D^{\mathrm{par}}_T$.
\end{proposition}
\begin{proof}
(a) As in Proposition~\ref{prop:haynsworth}(a): $-P_Y\mathcal GP_Y\ge-P_Y\mathcal G_+P_Y$ and
$\|P_Y\mathcal G_+P_Y\|\le\omega_+\operatorname{tr}(P^\perp S_{c_*}P^\perp)=\omega_+\Lambda_{k_1}(c_*)$.
(b) For $i<k_1$ and $y\in Y$, $\mathfrak b_{\mu,c}(e_i,y)=-\langle P_Y\mathcal Ge_i,y\rangle$ and the border pairs
with $y$ through $+\langle P_Y\chi_{\mathrm{par}},y\rangle$. The congruence of the proof of
Proposition~\ref{prop:haynsworth}(b) gives $\mathfrak b_{\mu,c}\cong\mathfrak s\oplus\mathfrak b_{\mu,c}|_Y$ with
$\mathfrak b_{\mu,c}|_Y\ge q_Y>0$, so $n_-(\mathfrak b_{\mu,c})=n_-(\mathfrak s)$, and the subtracted term of the Schur
complement $\mathfrak s$ is at most $q_Y^{-1}$ times the Gram matrix of the $Y$-columns
$-P_Y\mathcal Ge_i$ and $P_Y\chi_{\mathrm{par}}$; hence $\mathfrak s\ge\mathsf R\ge\mathsf R'$.
(c) Put $\mathfrak q_\mu(F)=\mathfrak b_{\mu,0}(F,0)$, so that $\mathfrak a_T-\mu\|\cdot\|^2\ge\mathfrak q_\mu$. For
$c=0$, Lemma~\ref{lem:border} gives $n_-(\mathfrak q_\mu|_{\ker\ell_{\mathrm{par}}})=n_-(\mathfrak b_{\mu,0})-1=0$. For
$c=\frac12$, $\mathfrak b_{\mu,1/2}=\mathfrak q_\mu-2\ell_{\mathrm{par}}^2+\frac12(\xi+2\ell_{\mathrm{par}})^2$, so
$\mathfrak q_\mu-2\ell_{\mathrm{odd}}^2\ge0$, and $W_T=\mathfrak a_T-2\ell_{\mathrm{odd}}^2$ on odd $F$
(Theorem~\ref{thm:parityschur}(o)). For $c=-\frac12$,
$\mathfrak b_{\mu,-1/2}=\mathfrak q_\mu+2\ell_{\mathrm{par}}^2-\frac12(\xi-2\ell_{\mathrm{par}})^2$; the substitution
$\eta=\xi-2\ell_{\mathrm{par}}(F)$ is a congruence, so
$n_-(\mathfrak q_\mu+2\ell_{\mathrm{ev}}^2)=n_-(\mathfrak b_{\mu,-1/2})-1=0$, and $W_T=\mathfrak a_T+2\ell_{\mathrm{ev}}^2$ on even
$F$. The statements pass from $L^2_{\mathrm{par}}(\IT)$ to the form domain because $\mathfrak q_\mu$ is bounded.
\end{proof}

\emph{The computation.} For $i-l$ even,
$\langle\mathcal Ge_i,e_l\rangle=\frac{2T}\pi\sqrt{(2i+1)(2l+1)}\,(-1)^{(l-i)/2}\int_0^{\tau_N}g(\tau)j_i(\tau T)j_l(\tau T)\,d\tau$
(proof of Lemma~\ref{lem:KyFan}), and the entries with $i-l$ odd vanish. The integral is computed by composite
Gauss--Legendre quadrature with $46$ nodes on panels of length $\frac12$, with the remainder bound for
functions analytic in a Bernstein ellipse \cite[Thm.~19.3]{Trefethen2013}, $\rho=3.2$. The ellipse lies
within distance $0.439<\frac12$ of the panel, while the poles of $\psi(\frac14\pm\frac{i\tau}2)$ lie at
$|\Iml\tau|\ge\frac12$. The modulus on the ellipse is bounded with $|j_n(z)|\le e^{|\Iml z|}$,
$|\cos(z\log n)|\le\cosh(\Iml z\log n)$ and, for $\psi$, the recurrence $\psi(w)=\psi(w+8)-\sum_{k<8}(w+k)^{-1}$
together with Binet's bound $|\psi(u)|\le\log|u|+\frac\pi2+\frac1{2\Rel u}+\frac1{12(\Rel u)^2}$ for
$\Rel u\ge8$. The resulting remainder, $1.5\cdot10^{-43}$ per entry (below $2\cdot10^{-42}$ under either
point-count convention for the rule), is added as a ball to every entry. The inner products over $Y$ are
exact sums over $k_1\le l<k_2$ of the given parity plus tail boxes: by Cauchy--Schwarz for the positive
operators $\mathcal G_\pm$, $\|P^\perp_{k_2}\mathcal G_\pm e\|^2\le\|P^\perp_{k_2}\mathcal G_\pm P^\perp_{k_2}\|
\langle\mathcal G_\pm e,e\rangle\le\omega_\pm^2\Lambda_{k_2}(c_*)$, so
$\|P^\perp_{k_2}\mathcal Ge\|\le(\omega_++\omega_-)\Lambda_{k_2}(c_*)^{1/2}$, and
$\|P^\perp_{k_2}\chi_{\mathrm{par}}\|^2=\sinh T\pm T-\sum_{l<k_2}\langle e_l,\chi_{\mathrm{par}}\rangle^2$ with
$\langle e_l,\chi_{\mathrm{par}}\rangle=\sqrt{(2l+1)/2T}\,2T\,i_l(T/2)$. The tails $\Lambda_k$ are bounded by
$\Lambda^*_k$ of Corollary~\ref{cor:P1}. Everything is midpoint--radius ball arithmetic at $640$ bits
(Arb \cite{Johansson2017}, through \texttt{python-flint}), including the digamma and Bessel values and the
Gauss--Legendre nodes and weights; every Bessel ball is asserted to have radius below $10^{-60}$, and
comparisons are ball predicates.

\emph{Inertia.} In $\mathsf R$ the term $\|P_Y\chi_{\mathrm{par}}\|^2/q_Y$ is of order $10^{-190}$; it is replaced
by an exact upper bound $\delta'$, which gives $\mathsf R'\le\mathsf R$. When $n_-(\mathsf R')=0$ is to be shown,
interval Cholesky of $\mathsf R'$ is applied. When $n_-(\mathsf R')\le1$ is to be shown, a double-precision
Householder matrix $H$ that sends the midpoint of the column $(\mathsf R'_{i\xi})_i$ to a multiple of $e_0$ is
used as an exact matrix; $\|H^{\mathsf T}H-I\|_\infty<1$ is verified in balls, so $\mathrm{diag}(H,1)$ is a
congruence. In the new coordinates the $2\times2$ block $B$ on $(e_0,\xi)$ has $\det B<0$, certified, so its
inertia is $(1,1)$, and by Haynsworth's inertia additivity \cite{Haynsworth1968}
$n_-(\mathsf R')=1+n_-(S)$ for the Schur complement $S$ of $B$, whose positivity is verified by interval
Cholesky. The only non-rigorous objects are $H$ and the midpoints, which enter only through this exact
congruence.

\emph{Data.} $T=\frac12\log5$, $\tau_N=240$, $c_*=193.13$, $k_1=280$, $k_2=330$ ($\dim E=140$ in each parity).
From $m(\tau_N)=3.64276\ldots$: $s_\infty=0.700787608\ldots$, $\omega_+=9.0150$, $\omega_-=5.8839$,
$\Lambda^*_{k_1}(c_*)\le4.4\cdot10^{-20}$, $\Lambda^*_{k_2}(c_*)\le3.9\cdot10^{-70}$, $q_Y\ge0.70079$, tail
box $(\omega_++\omega_-)\Lambda^*_{k_2}(c_*)^{1/2}\le2.9\cdot10^{-34}$; the entry radii of $\mathsf R$ are at
most $4.4\cdot10^{-41}$. The verified cases are:
\begin{center}\small
\begin{tabular}{@{}llccc@{}}
\toprule
form & sector & certified $\mu$ & minimal pivot at that $\mu$ & fails at\\
\midrule
$\mathfrak a_T|\PT$ ($c=0$) & even & $2.5\cdot10^{-12}$ & $8.47\cdot10^{-3}$ & $3\cdot10^{-12}$\\
$\mathfrak a_T|\PT$ ($c=0$) & odd & $3.0\cdot10^{-10}$ & $2.03\cdot10^{-2}$ & $4\cdot10^{-10}$\\
$W_T$ ($c=-\frac12$) & even & $10^{-18}$ & $4.47\cdot10^{-4}$ & $10^{-16}$\\
$W_T$ ($c=\frac12$) & odd & $10^{-16}$ & $8.14\cdot10^{-14}$ & $10^{-14}$\\
\bottomrule
\end{tabular}
\end{center}
The radius of the smallest pivot in the last row is below $10^{-24}$. A failed Cholesky is not evidence of
negativity. A double-precision eigenvalue computation of the midpoint of $S$ at $\mu=0$ gives
$2.53\cdot10^{-12}$ (even) and $3.62\cdot10^{-10}$ (odd) for the primitive form, so the certified bounds are
within a small factor of the bottom of the lower form. As a check, the same procedure at
$T=\frac12\log3$ ($\tau_N=40$, $k_1=50$) certifies $\lambda^{\mathrm{ev}}_{\mathrm A}(T)\ge6.3\cdot10^{-4}$,
consistent with Theorem~\ref{thm:C}. Each run takes about seven minutes per parity on one core.


\section{Numerical values}\label{app:num}

The values in this appendix are computed in double precision and are not certified; no proof in the paper
depends on them. Galerkin values are upper bounds for the bottoms; the three values of $\lamA$ below are
converged in the basis size.

\begin{center}\small
\begin{tabular}{@{}lll@{}}
\toprule
quantity & value & certified bound in the paper\\
\midrule
$\lambda_{\mathrm A}(\frac12\log3)$ & $9.448\cdot10^{-4}$ & $\ge8.8577\cdot10^{-5}$ (Theorem~\ref{thm:C}), $\ge2.5\cdot10^{-12}$ (Theorem~\ref{thm:log5A})\\
$\lambda_{\mathrm A}(\log2)$ & $9.83\cdot10^{-8}$ & $\ge2.5\cdot10^{-12}$ (Theorem~\ref{thm:log5A})\\
$\lambda_{\mathrm A}(\frac12\log5)$ & $3.10\cdot10^{-12}$ & $\ge2.5\cdot10^{-12}$ (Theorem~\ref{thm:log5A})\\
\bottomrule
\end{tabular}
\end{center}

\begin{itemize}
\item Test-kernel ceilings at $T=\frac12\log3$: $\operatorname{Ceil}(8)=-5.4\cdot10^{-3}$,
$\operatorname{Ceil}(12)=7.6\cdot10^{-5}$, $\operatorname{Ceil}(24)=5.9\cdot10^{-4}$; the certificate of
Theorem~\ref{thm:C} realizes about $15\%$ of its class ceiling.
\item The lower form used in \S\ref{subsec:encllog5} has bottom $2.53\cdot10^{-12}$ (even) and
$3.62\cdot10^{-10}$ (odd) at $T=\frac12\log5$ (double-precision midpoint values, numerical); the certified bounds of Theorem~\ref{thm:log5A} lie below these, and the smaller one, $2.53\cdot10^{-12}$, is consistent with the value $3.10\cdot10^{-12}$ above.
\item Minimized over $k$, the right side of \eqref{eq:Bmargin} exceeds $1$ for $T<1.430$; the bound of
Theorem~\ref{thm:margin} is asymptotic and vacuous at the windows where the sign is known.
\end{itemize}


\begin{thebibliography}{99}
\bibitem{BoasKac1945} R.~P.~Boas and M.~Kac, \emph{Inequalities for Fourier transforms of positive
functions}, Duke Math.~J. \textbf{11} (1945), 189--206.
\bibitem{Bombieri2000} E.~Bombieri, \emph{Remarks on Weil's quadratic functional in the theory of
prime numbers, I}, Atti Accad.~Naz.~Lincei Cl.~Sci.~Fis.~Mat.~Natur.~Rend.~Lincei (9) Mat.~Appl.
\textbf{11} (2000), no.~3, 183--233.
\bibitem{CohnElkies2003} H.~Cohn and N.~Elkies, \emph{New upper bounds on sphere packings I},
Ann.~of Math. \textbf{157} (2003), 689--714.
\bibitem{ConnesConsani2021} A.~Connes and C.~Consani, \emph{Weil positivity and trace formula, the archimedean place}, Selecta Math.\ (N.S.) \textbf{27} (2021), no.~4, Paper No.~77; arXiv:2006.13771.
\bibitem{DLMF} NIST Digital Library of Mathematical Functions, \url{https://dlmf.nist.gov}, \S5.9.
\bibitem{Groskin2026} A.~Groskin, \emph{A finite Guinand--Weil dictionary and archimedean tail
order for the truncated Weil quadratic form}, arXiv:2607.02828 (2026).
\bibitem{Haynsworth1968} E.~V.~Haynsworth, \emph{Determination of the inertia of a partitioned Hermitian
matrix}, Linear Algebra Appl.\ \textbf{1} (1968), 73--81.
\bibitem{Katznelson} Y.~Katznelson, \emph{An Introduction to Harmonic Analysis}, 3rd ed.,
Cambridge Univ.~Press, 2004.
\bibitem{IK} H.~Iwaniec and E.~Kowalski, \emph{Analytic Number Theory}, Amer.~Math.~Soc.
Colloq.~Publ.~53, 2004.
\bibitem{Johansson2017} F.~Johansson, \emph{Arb: efficient arbitrary-precision midpoint-radius interval
arithmetic}, IEEE Trans.\ Comput.\ \textbf{66} (2017), 1281--1292.
\bibitem{Krein1940} M.~G.~Kre\u{\i}n, \emph{Sur le probl\`eme du prolongement des fonctions
hermitiennes positives et continues}, Dokl.~Akad.~Nauk SSSR \textbf{26} (1940), 17--22.
\bibitem{PlattTrudgian2021} D.~J.~Platt and T.~S.~Trudgian, \emph{The Riemann hypothesis is true up to
$3\cdot10^{12}$}, Bull.\ London Math.\ Soc.\ \textbf{53} (2021), 792--797.
\bibitem{Rockafellar} R.~T.~Rockafellar, \emph{Convex Analysis}, Princeton Univ.~Press, 1970.
\bibitem{Rosser1941} J.~B.~Rosser, \emph{Explicit bounds for some functions of prime numbers}, Amer.\ J.\
Math.\ \textbf{63} (1941), 211--232.
\bibitem{RosserYoheSchoenfeld1968} J.~B.~Rosser, J.~M.~Yohe and L.~Schoenfeld, \emph{Rigorous computation and
the zeros of the Riemann zeta-function}, Proc.\ IFIP Congress 1968, 70--76.
\bibitem{Sasvari1994} Z.~Sasv\'ari, \emph{Positive Definite and Definitizable Functions},
Akademie Verlag, Berlin, 1994.
\bibitem{SSV2012} R.~L.~Schilling, R.~Song and Z.~Vondra\v{c}ek, \emph{Bernstein Functions:
Theory and Applications}, 2nd ed., de Gruyter, 2012.
\bibitem{SongVondracek2003} R.~Song and Z.~Vondra\v{c}ek, \emph{Potential theory of subordinate
killed Brownian motion in a domain}, Probab.~Theory Related Fields \textbf{125} (2003), 578--592.
\bibitem{Suzuki2023jlms} M.~Suzuki, \emph{Aspects of the screw function corresponding to the
Riemann zeta function}, J.~Lond.~Math.~Soc.~(2) \textbf{107} (2023), 1846--1888.
\bibitem{Suzuki2023screwline} M.~Suzuki, \emph{The screw line of the Riemann zeta-function and
its applications}, arXiv:2209.04658.
\bibitem{Suzuki2025hilbert} M.~Suzuki, \emph{On the Hilbert space derived from the Weil
distribution}, arXiv:2301.00421.
\bibitem{Trefethen2013} L.~N.~Trefethen, \emph{Approximation Theory and Approximation Practice},
SIAM, Philadelphia, 2013.
\bibitem{Paper45} D.~A.~Trejo Pizzo, \emph{An explicit formula and a comparison theorem for the canonical trace on a correspondence
algebra of the idele class group}, preprint, 2026.
\bibitem{Weil52} A.~Weil, \emph{Sur les ``formules explicites'' de la th\'eorie des nombres
premiers}, Comm.~S\'em.~Math.~Univ.~Lund (1952), 252--265.
\bibitem{Yoshida1992} H.~Yoshida, \emph{On Hermitian forms attached to zeta functions}, in:
\emph{Zeta Functions in Geometry}, Adv.~Stud.~Pure Math.~\textbf{21}, Kinokuniya, Tokyo, 1992, 281--325.
\bibitem{Zhu2026} X.~Zhu, \emph{Weil positivity in compact windows: a finite reduction, certified
two-sided bounds, and a Landau--Widom decay law}, arXiv:2608.24827 (2026).
\end{thebibliography}

\end{document}

